% =============================================================================
% LG_Master_Unified.tex
% Flattened from the canonical LG_Master modular source on 2026-07-25.
% All files imported by main.tex are inlined below in canonical build order.
% The explicitly noncanonical development_archive is not part of this build.
% =============================================================================

\documentclass[11pt,reqno]{scrartcl}

\usepackage[margin=1in]{geometry}
\usepackage[T1]{fontenc}
\usepackage[utf8]{inputenc}
\usepackage{lmodern}
\usepackage{microtype}
\usepackage{setspace}
\usepackage{amsmath,amssymb,amsthm,mathtools,mathrsfs}
\usepackage{booktabs,tabularx,array,longtable}
\usepackage{enumitem}
\usepackage{xcolor}
\usepackage[colorlinks=true,linkcolor=blue!60!black,citecolor=blue!60!black,urlcolor=blue!60!black]{hyperref}
\hypersetup{
  pdftitle={The Luminous Ground: Emptiness, Awareness, and the Algebra of Appearance},
  pdfauthor={Alexander Yiannopoulos},
  pdfsubject={Cognitive realization of the universal transition architecture: cyclic grounds, Born correspondence, and the algebra of appearance},
  pdfkeywords={Awareness Operator, Volition, Krein space, luminosity, emptiness, reflexive awareness, Yogacara, Dharmakirti, graded operator algebra}
}
\usepackage{aliascnt}
\usepackage[nameinlink,capitalise,noabbrev]{cleveref}
\usepackage{scrlayer-scrpage}

% =============================================================================
% BEGIN INLINED FILE: macros.tex
% =============================================================================
% Core notation for The Luminous Ground
\newcommand{\CC}{\mathbb C}
\newcommand{\RR}{\mathbb R}
\newcommand{\NN}{\mathbb N}
\newcommand{\ZZ}{\mathbb Z}

\newcommand{\scrK}{\mathscr K}
\newcommand{\scrHP}{\mathscr H_{\Psi}}
\newcommand{\scrHF}{\mathscr H_{\Phi}}
\newcommand{\algP}{\mathfrak M_{\Psi}}
\newcommand{\algF}{\mathfrak M_{\Phi}}
\newcommand{\piP}{\pi_{\Psi}}
\newcommand{\piF}{\pi_{\Phi}}
\newcommand{\omegaP}{\omega_{0}^{\Psi}}
\newcommand{\omegaF}{\omega_{0}^{\Phi}}
\newcommand{\fundJ}{\mathfrak J}
\newcommand{\PP}{P_{\Psi}}
\newcommand{\PF}{P_{\Phi}}
\newcommand{\PzeroP}{P_{0}^{\Psi}}
\newcommand{\QzeroP}{Q_{0}^{\Psi}}
\newcommand{\PzeroF}{P_{0}^{\Phi}}
\newcommand{\QzeroF}{Q_{0}^{\Phi}}
\newcommand{\GzeroP}{\mathscr G_{\Psi}}
\newcommand{\GzeroF}{\mathscr G_{\Phi}}
\newcommand{\Kzero}{\mathscr K_{0}}
\newcommand{\Pizero}{\Pi_{0}}
\newcommand{\idop}{\mathbf 1}

\newcommand{\Aodd}{\mathsf A}
\newcommand{\Vodd}{\mathsf V}
\newcommand{\Atyped}{A}
\newcommand{\Vtyped}{V}
\newcommand{\polA}{u_A}
\newcommand{\absA}{|A|}
\newcommand{\Azero}{A_{0}}
\newcommand{\uzero}{u_{0}}
\newcommand{\Rzero}{\mathcal R_{0}}
\newcommand{\Izero}{\mathcal I_{0}}
\newcommand{\Qzero}{\mathcal Q_{0}}
\newcommand{\RA}{\mathcal R_{A}}
\newcommand{\IA}{\mathcal I_{A}}
\newcommand{\PiA}{\Pi_{A}}
\newcommand{\NzeroPlus}{\mathscr N_{0}^{+}}
\newcommand{\NzeroMinus}{\mathscr N_{0}^{-}}
\newcommand{\WzeroPlus}{W_{0}^{+}}
\newcommand{\WzeroMinus}{W_{0}^{-}}
\newcommand{\iP}{\iota_{\Psi}}
\newcommand{\iF}{\iota_{\Phi}}

\newcommand{\ket}[1]{\lvert #1\rangle}
\newcommand{\bra}[1]{\langle #1\rvert}
\newcommand{\braket}[2]{\langle #1\mid #2\rangle}
\newcommand{\inner}[2]{\langle #1,#2\rangle}
\newcommand{\kInner}[2]{[#1,#2]_{\fundJ}}
\newcommand{\proj}[1]{\ket{#1}\!\bra{#1}}

\DeclareMathOperator{\Dom}{Dom}
\DeclareMathOperator{\Ran}{Ran}
\DeclareMathOperator{\Ker}{Ker}
\DeclareMathOperator{\spec}{spec}
\DeclareMathOperator{\diag}{diag}
\DeclareMathOperator{\Graph}{Graph}
\DeclareMathOperator{\rank}{rank}
\DeclareMathOperator{\Span}{span}
\DeclareMathOperator{\Cyc}{Cyc}
\DeclareMathOperator{\ind}{ind}
\DeclareMathOperator{\coker}{coker}
\DeclareMathOperator{\supp}{supp}

\newcommand{\dsum}{\mathbin{[\dotplus]}}
\newcommand{\restr}[2]{#1\!\upharpoonright_{#2}}
\newcommand{\projective}[1]{\mathbb P\!\left(#1\right)}

% Status labels
\newcommand{\statuslabel}[1]{\par\smallskip\noindent\textsc{Status: #1.}\ }
\newcommand{\warninglabel}{\par\smallskip\noindent\textsc{Caution.}\ }
\newcommand{\passlabel}[1]{\par\smallskip\noindent\textsc{Pass #1 boundary.}\ }

% Pass 4: active-support BRLC and paired observables
\newcommand{\actHP}{\mathscr H_{\Psi}^{A}}
\newcommand{\actHF}{\mathscr H_{\Phi}^{A}}
\newcommand{\actK}{\mathscr K_{A}}
\newcommand{\suppProjP}{P_{A}^{\Psi}}
\newcommand{\suppProjF}{P_{A}^{\Phi}}
\newcommand{\defectProjF}{P_{A}^{\Phi,\perp}}
\newcommand{\piAplus}{\pi_{A,+}} % TAUM crosswalk only
\newcommand{\piAminus}{\pi_{A,-}} % TAUM crosswalk only
\newcommand{\NAplus}{\mathscr N_{A}^{+}}
\newcommand{\NAminus}{\mathscr N_{A}^{-}}
\newcommand{\NAtheta}{\mathscr N_{A}(\theta)}
\newcommand{\WA}{W_{A}}
\newcommand{\WAtheta}{W_{A}(\theta)}
\newcommand{\LBR}{\mathscr L_{\mathrm{BR}}(A)}
\newcommand{\alphaA}{\alpha_{A}}
\newcommand{\iotaA}{\iota_{A}}
\newcommand{\iotaMaxA}{\widetilde{\iota}_{A}}
\newcommand{\obsP}{\mathfrak O_{\Psi}^{A}}
\newcommand{\obsF}{\mathfrak O_{\Phi}^{A}}
\newcommand{\obsPair}{\mathfrak O_{\mathrm{pair}}^{A}}
\newcommand{\UA}{\mathcal U_{A}}
\newcommand{\PAPlus}{\mathcal P_{A}^{+}}
\newcommand{\PAMinus}{\mathcal P_{A}^{-}}
\newcommand{\PAtheta}{\mathcal P_{A}(\theta)}
\newcommand{\BRprob}{\mathsf p}
\newcommand{\sigmax}{\sigma_x}
\newcommand{\sigmay}{\sigma_y}
\newcommand{\sigmaz}{\sigma_z}
\newcommand{\eplus}{e_{+}}
\newcommand{\eminus}{e_{-}}
\DeclareMathOperator{\Tr}{Tr}

% Pass 5: algebra of appearance
\newcommand{\detAlg}{\mathfrak D_{\Phi}}
\newcommand{\nullIdeal}{\mathcal N_{\omega_{0}^{\Phi}}}
\newcommand{\centDet}{\mathfrak D_{\Phi}^{\circ}}
\newcommand{\appSpace}{\mathscr E_{\Phi}}
\newcommand{\vasSpace}{\mathcal P_{\mathrm{vas}}}
\newcommand{\vas}{\nu}
\newcommand{\vasEmpty}{\nu_{\varnothing}}
\newcommand{\genSet}{\mathcal G_{\Phi}}
\newcommand{\wordSet}{\mathfrak W_{\Phi}}
\newcommand{\Hvas}{\mathscr H_{\Phi}[\nu]}
\newcommand{\HvasEmpty}{\mathscr H_{\Phi}[\nu_{\varnothing}]}
\newcommand{\Evas}{\mathscr E_{\Phi}[\nu]}
\newcommand{\Gak}{\mathcal G_{\mathrm{ak}}}
\newcommand{\Gvik}{\mathcal G_{\mathrm{vik}}}
\newcommand{\Gsmr}{\mathcal G_{\mathrm{smr}}}
\newcommand{\Gaff}{\mathcal G_{\mathrm{aff}}}
\newcommand{\Gdv}{\mathcal G_{\mathrm{dv}}}
\newcommand{\contentAlg}{\mathfrak D_{\Phi}^{\mathrm{cont}}}
\newcommand{\affectAlg}{\mathfrak D_{\Phi}^{\mathrm{aff}}}
\newcommand{\Pdual}{P_{\mathrm{dv}}^{\phi}}
\newcommand{\Pgrasper}{P_{\mathrm{gr}}^{\phi}}
\newcommand{\Pgrasped}{P_{\mathrm{gd}}^{\phi}}
\newcommand{\Gammadual}{\Gamma_{\mathrm{dv}}^{\phi}}
\newcommand{\semsSet}{\mathscr S_{\mathrm{sems}}}
\newcommand{\rigRay}{\mathscr R_{\mathrm{rig}}}
\DeclareMathOperator{\Alg}{Alg}
\DeclareMathOperator{\Cov}{Cov}

% Pass 6: reflexivity, non-alteration, and deconditioning
\newcommand{\detAlgEmpty}{\mathfrak D_{\Phi,\varnothing}}
\newcommand{\groundChar}{\chi_{0}^{\Phi}}
\newcommand{\Decond}{\mathcal D}
\newcommand{\DynChan}{\mathcal T}
\newcommand{\Lzero}{\mathcal L_{0}}
\DeclareMathOperator{\Fix}{Fix}

% Pass 7/7.1: Volition, presentability, phenomenal defect, and index
\newcommand{\defectHF}{(\mathscr H_{\Phi}^{A})^{\perp}}
\newcommand{\actAppF}{\mathscr E_{\Phi}^{A}}
\newcommand{\presentWeight}{w_{A}}
\newcommand{\defectWeight}{w_{A}^{\perp}}
\newcommand{\returnOp}{\mathcal V_{A}}
\newcommand{\returnOpCond}{\widehat{\mathcal V}_{A}}
\newcommand{\returnGroundOp}{\mathcal V_{0}}
\newcommand{\cohZero}{\mathsf H^{0}(A)}
\newcommand{\cohOne}{\mathsf H^{1}(A)}
% =============================================================================
% END INLINED FILE: macros.tex
% =============================================================================

\setstretch{1.13}
\setlength{\emergencystretch}{3em}
\setkomafont{section}{\Large\bfseries}
\setkomafont{subsection}{\large\bfseries}
\setkomafont{disposition}{\normalfont}
\DeclareTOCStyleEntry[numwidth=3.8em]{tocline}{subsection}

\clearpairofpagestyles
\ihead{\textit{The Luminous Ground}}
\ohead{}
\cfoot{\pagemark}

\theoremstyle{plain}
\newtheorem{theorem}{Theorem}[section]
\newaliascnt{proposition}{theorem}
\newtheorem{proposition}[proposition]{Proposition}
\aliascntresetthe{proposition}
\newaliascnt{lemma}{theorem}
\newtheorem{lemma}[lemma]{Lemma}
\aliascntresetthe{lemma}
\newaliascnt{corollary}{theorem}
\newtheorem{corollary}[corollary]{Corollary}
\aliascntresetthe{corollary}

\theoremstyle{definition}
\newaliascnt{definition}{theorem}
\newtheorem{definition}[definition]{Definition}
\aliascntresetthe{definition}
\newaliascnt{axiom}{theorem}
\newtheorem{axiom}[axiom]{Realization axiom}
\aliascntresetthe{axiom}
\newaliascnt{identification}{theorem}
\newtheorem{identification}[identification]{Doctrinal identification}
\aliascntresetthe{identification}
\newaliascnt{modelchoice}{theorem}
\newtheorem{modelchoice}[modelchoice]{Model choice}
\aliascntresetthe{modelchoice}
\newaliascnt{openquestion}{theorem}
\newtheorem{openquestion}[openquestion]{Open question}
\aliascntresetthe{openquestion}

\theoremstyle{remark}
\newaliascnt{remark}{theorem}
\newtheorem{remark}[remark]{Remark}
\aliascntresetthe{remark}

\crefname{theorem}{theorem}{theorems}
\crefname{proposition}{proposition}{propositions}
\crefname{lemma}{lemma}{lemmas}
\crefname{corollary}{corollary}{corollaries}
\crefname{definition}{definition}{definitions}
\crefname{axiom}{realization axiom}{realization axioms}
\crefname{identification}{doctrinal identification}{doctrinal identifications}
\crefname{modelchoice}{model choice}{model choices}
\crefname{openquestion}{open question}{open questions}
\crefname{remark}{remark}{remarks}
\crefname{section}{section}{sections}
\crefname{appendix}{appendix}{appendices}

\title{\textbf{The Luminous Ground}\\[3pt]
\Large Emptiness, Awareness, and the Algebra of Appearance\\[7pt]
\large A Cognitive Realization of the Universal Transition Architecture}
\author{Alexander Yiannopoulos\\[2pt]
\small A$\lvert\Omega\rangle$ Institute, New Orleans, Louisiana, USA\\
\small\texttt{ayianno@ao.institute}}
\date{July 2026; revised November 10, 2026\\[5pt]\small Canonical reference manuscript}

\begin{document}
\maketitle
\vspace{-1.5\baselineskip}

\begin{abstract}
\noindent
This reference develops a specifically cognitive and Buddhist realization of
the interpretation-neutral transition architecture of \emph{Toward a
Universal Model}.  A densely defined closed channel
$A:\Dom A\subseteq\scrHP\to\scrHF$ and its adjoint are distinguished from
their grading-odd extensions, their supportwise polar orientations, and the
restricted rank-one ground channel.  Additional realization axioms select
unique cyclic physical and phenomenal grounds, impose the bridge
$A\ket0=a_0\ket{\Omega_0}$, and require asymmetric presentability through
$\Ker A=\{0\}$ and $\Ker A^*\neq\{0\}$.  The bridge determines a
Krein-neutral relational ground, while the polar orientation determines an
active-support Born Rule Locus Correspondence and a paired observable
representation preserving projective and effect-valued Born statistics.
Within the phenomenal sector, a cyclic determination algebra generates
differentiated appearance from the luminous ground; v\=asan\=a profiles index
image-forming, conceptual, mnemonic, affective, and duality-producing
operations.  Reflexive terminality, exact ground realization,
ground-non-alteration, and deconditioning are kept distinct, and a canonical
completely positive semigroup provides an explicit ground-attractive model.
Volition is realized as adjoint return and normalized reverse orientation; its
completely positive trace-nonincreasing operation has trace equal to the final
support weight.  The phenomenal domain decomposes into the luminous ground,
differentiated appearance within the final support, and the orthogonal defect
$\Ker A^*$, interpreted as phenomenal excess.  Under Fredholm hypotheses,
physical presentability yields $\ind A=-\dim\Ker A^*$.
A finite pair model, using the bare transfer component or a strictly
non-minimal symmetric positive completion, exhibits strong ground
correspondence with a dynamically accessible final defect.  A continuum
orbital parent supplies a ground-preserving source and a selected
three-state realization; preparation and readout restrictions remain
explicit.  Throughout, imported
theorems, realization axioms, model choices, and doctrinal identifications are
kept logically distinct.
\end{abstract}

\medskip
\noindent\textbf{Keywords:} Awareness Operator; Volition; Krein space; graded
operator algebra; polar decomposition; emptiness; luminosity; reflexive
awareness; Yog\=ac\=ara; Dharmak\=irti; nonduality.

\clearpage
\setcounter{tocdepth}{2}
\tableofcontents

\clearpage
% =============================================================================
% BEGIN INLINED FILE: sections/00_charter.tex
% =============================================================================
\section{Scope, status, and axiomatic spine}
\label{sec:charter}

\subsection{Purpose and logical hierarchy}

\emph{The Luminous Ground} develops a cognitive and Buddhist realization of
the interpretation-neutral operator architecture supplied by \emph{Toward a
Universal Model} (TAUM).  The inherited chassis consists of a bipartite
Hilbert/Krein carrier, a densely defined closed channel, its nilpotent graded
orientations, positive ordered returns, and the associated polar geometry.
The present document adds the structures required to interpret that chassis
in terms of matter, awareness, appearance, volition, and a luminous ground.

The governing logical hierarchy is
\begin{equation}
  \boxed{
  \text{TAUM chassis}
  \quad\Longrightarrow\quad
  \text{cognitive realization}
  \quad\Longrightarrow\quad
  \text{doctrinal interpretation}.}
  \label{eq:hierarchy}
\end{equation}
The arrows in \eqref{eq:hierarchy} do not denote logical equivalence.  A
theorem of the chassis may be imported into the realization; a realization
axiom supplies additional mathematical data; and a doctrinal identification
assigns a Buddhist philosophical or contemplative role to an already typed
mathematical object.  No implication is valid merely because the adjacent
layers use the same notation.

The scope is a single cognitive realization.  Families of realizations,
ground-line bundles, and moving-projection geometry require an independently
specified parameter space and are not assumed here.  Likewise, the
structural relative-phase operators receive no laboratory measurement
interpretation unless a concrete observable model is supplied.

\subsection{Status discipline}

\begin{definition}[Status taxonomy]
\label{def:status-taxonomy}
Each substantive statement in the canonical architecture has one of the
following statuses.
\begin{enumerate}[label=(\alph*)]
  \item An \emph{imported theorem} is proved in TAUM or is a standard result
  of operator theory under explicitly stated hypotheses.
  \item A \emph{derived proposition} follows from imported results together
  with the realization data explicitly in force.
  \item A \emph{realization axiom} adds mathematical structure not entailed
  by the interpretation-neutral chassis.
  \item A \emph{model choice} selects one admissible realization among
  mathematically available alternatives.
  \item A \emph{doctrinal identification} assigns a philosophical or
  contemplative meaning to a mathematical object or relation.
  \item An \emph{open question} records structure that has not been selected
  or derived.
\end{enumerate}
\end{definition}

This taxonomy is substantive rather than editorial.  For example,
\begin{equation}
  \Aodd^2=0
\end{equation}
is an imported algebraic statement, whereas reading nilpotency as reflexive
manifestness without an iterated Awareness regress is a doctrinal
identification.  Neither statement proves the other.  Historical and
superseded constructions are retained separately as development records and
are not premises of the canonical architecture.

\subsection{The five realization axioms}
\label{sec:axiomatic-spine}

The cognitive realization rests on five assumptions that are additional to
the closed-channel chassis.  Their complete formulations occur at their
points of use; they are collected here to make the dependency structure
explicit.

\begin{enumerate}[label=\textup{(R\arabic*)}]
  \item \textbf{Unique cyclic physical vacuum.}
  The physical representation has a normalized cyclic vector \(\ket0\) whose
  ray is the unique ground ray of its supplied physical ground generator;
  see \cref{ax:physical-vacuum}.

  \item \textbf{Unique cyclic luminous ground.}
  The phenomenal representation has a normalized cyclic vector
  \(\ket{\Omega_0}\) whose ray is the unique ground ray of its supplied
  phenomenal ground generator; see \cref{ax:phenomenal-ground}.

  \item \textbf{Ground bridge.}
  The typed channel carries the physical vacuum into the intrinsic luminous
  ground ray.  After fixing the phase of \(\ket{\Omega_0}\),
  \begin{equation}
    \ket\Omega:=A\ket0=a_0\ket{\Omega_0},
    \qquad
    a_0:=\|A\ket0\|>0.
    \label{eq:bridge-preview}
  \end{equation}
  This is \cref{ax:ground-bridge}, not a consequence of the bare channel.

  \item \textbf{Physical presentability.}
  No nonzero physical direction is annihilated by the channel:
  \begin{equation}
    \Ker A=\{0\}.
  \end{equation}
  Equivalently, the polar orientation has full initial support; see
  \cref{ax:physical-presentability}.

  \item \textbf{Strict phenomenal excess.}
  The final phenomenal support has a nontrivial orthogonal defect:
  \begin{equation}
    \Ker A^*\neq\{0\}.
  \end{equation}
  This is \cref{ax:strict-phenomenal-excess}.  Operator theory supplies the
  defect space; its interpretation as phenomenal excess is additional
  realization content.
\end{enumerate}

The stronger condition
\begin{equation}
  A^*A\ket0=a_0^2\ket0
\end{equation}
is not part of this five-axiom spine.  It is the optional strong polar ground
compatibility of \cref{choice:strong-polar-ground}.  The weak bridge already
defines the restricted ground orientation, while the stronger condition is
needed to identify that orientation with the compression of the full polar
geometry and to conclude \(u_A\ket0=\ket{\Omega_0}\).

\subsection{Core notation and type discipline}

The common carrier and its fundamental symmetry are
\begin{equation}
  \scrK=\scrHP\oplus\scrHF,
  \qquad
  \fundJ=\begin{pmatrix}\idop_{\Psi}&0\\0&-\idop_{\Phi}\end{pmatrix},
  \label{eq:carrier-charter}
\end{equation}
where \(\scrHP\) is the physical or matter sector and \(\scrHF\) is the
phenomenal sector.  The positive Hilbert product normalizes states.  The
indefinite product
\begin{equation}
  \kInner{\xi}{\eta}:=\inner{\xi}{\fundJ\eta}
\end{equation}
records the sector sign and is not a probability norm.

The bare symbol
\begin{equation}
  A:\Dom A\subseteq\scrHP\longrightarrow\scrHF
  \label{eq:typed-A-charter}
\end{equation}
denotes the densely defined closed Awareness channel.  Its Hilbert adjoint is
\begin{equation}
  V:=A^*:\Dom A^*\subseteq\scrHF\longrightarrow\scrHP.
  \label{eq:typed-V-charter}
\end{equation}
The sans-serif symbols \(\Aodd\) and \(\Vodd\) denote the corresponding
nilpotent odd operators on the graded carrier; their sign convention and
domains are fixed in \cref{def:odd-extensions}.

The distinguished vectors are
\begin{equation}
  \ket0\in\scrHP,
  \qquad
  \ket\Omega:=A\ket0\in\scrHF,
  \qquad
  \ket{\Omega_0}\in\scrHF.
  \label{eq:ground-charter}
\end{equation}
The first is the normalized cyclic physical vacuum, the second is its typed
Awareness image, and the third is the normalized cyclic luminous ground.
Their projective coincidence follows from the bridge
\eqref{eq:bridge-preview}; it is not a notational identity.  The bridge also
defines the bounded rank-one orientation
\begin{equation}
  u_0:=\frac{A P_0^{\Psi}}{\|A\ket0\|},
  \qquad
  \mathscr N_0^+:=\Graph(u_0).
\end{equation}

For the full polar decomposition \(A=u_A|A|\), write
\begin{equation}
  \suppProjP:=u_A^*u_A,
  \qquad
  \suppProjF:=u_Au_A^*,
  \qquad
  \defectProjF:=I_{\Phi}-\suppProjF.
  \label{eq:support-notation-charter}
\end{equation}
These are the initial-support, final-support, and final-defect projections.
Their ranges are \((\Ker A)^\perp\), \(\overline{\Ran A}\), and
\(\Ker A^*\), respectively.

\subsection{Typed, odd, and polar levels}

Three mathematically different roles must remain distinct:
\begin{enumerate}[label=(\roman*)]
  \item a typed channel between two Hilbert spaces;
  \item a nilpotent odd operator on their graded direct sum;
  \item a norm-preserving polar orientation carrying Born geometry on active
  support.
\end{enumerate}
They are represented by
\begin{equation}
  \boxed{
  A:\Dom A\subseteq\scrHP\to\scrHF,
  \qquad
  \Aodd:\Dom\Aodd\subseteq\scrK\to\scrK,
  \qquad
  u_A:(\Ker A)^\perp\to\overline{\Ran A}.}
  \label{eq:three-A-levels}
\end{equation}
The typed channel carries direction and transition magnitude, the odd
extension carries grading and nilpotency, and the polar partial isometry
carries the canonical norm-preserving correspondence.  Statements about one
of these objects cannot be transferred to another without an explicit
hypothesis.
% =============================================================================
% END INLINED FILE: sections/00_charter.tex
% =============================================================================
% =============================================================================
% BEGIN INLINED FILE: sections/02_typed_repair.tex
% =============================================================================
\section{Typed channel, odd extension, and polar orientation}
\label{sec:typed-repair}

\subsection{Typed channel, adjoint return, and odd extensions}

\begin{definition}[Typed Awareness and Volition]
\label{def:typed-AV}
Let \(\scrHP\) and \(\scrHF\) be complex Hilbert spaces.  A typed Awareness
channel is a densely defined closed operator
\begin{equation}
  A:\Dom A\subseteq\scrHP\longrightarrow\scrHF.
  \label{eq:typed-A}
\end{equation}
Its Hilbert adjoint
\begin{equation}
  V:=A^*:\Dom A^*\subseteq\scrHF\longrightarrow\scrHP
  \label{eq:typed-V}
\end{equation}
is the typed return channel.
\end{definition}

The notation \(V=A^*\) records the Hilbert-space adjoint relation at the typed
level.  The Krein sign appears only after the maps are embedded in the graded
carrier.

Let
\begin{equation}
  \iP\psi:=\psi\oplus0,
  \qquad
  \iF\phi:=0\oplus\phi
  \label{eq:embeddings}
\end{equation}
be the canonical embeddings, and let \(\PP\) and \(\PF\) be the corresponding
orthogonal projections.

\begin{definition}[TAUM-compatible odd extensions]
\label{def:odd-extensions}
On
\begin{equation}
  \Dom\Aodd=\Dom A\oplus\scrHF,
  \qquad
  \Dom\Vodd=\scrHP\oplus\Dom A^*,
\end{equation}
define
\begin{equation}
  \Aodd
  :=\begin{pmatrix}0&0\\-A&0\end{pmatrix},
  \qquad
  \Vodd
  :=\begin{pmatrix}0&A^*\\0&0\end{pmatrix}.
  \label{eq:odd-extensions}
\end{equation}
Thus
\begin{equation}
  \Aodd\iP\psi=-\iF(A\psi),
  \qquad
  \Vodd\iF\phi=\iP(A^*\phi).
  \label{eq:odd-action}
\end{equation}
\end{definition}

The minus sign in \(\Aodd\) is the convention used by the Krein branch of
TAUM.  At the level of projective rays it does not distinguish
\(A\psi\) from \(-A\psi\).  The typed equation
\begin{equation}
  \ket\Omega:=A\ket0
  \label{eq:typed-Omega}
\end{equation}
is therefore compatible with the embedded relation
\begin{equation}
  \Aodd\iP\ket0=-\iF\ket\Omega.
  \label{eq:embedded-Omega}
\end{equation}

\begin{proposition}[Imported oriented algebra]
\label{prop:imported-algebra}
The odd extensions in \cref{def:odd-extensions} satisfy
\begin{align}
  \Aodd^2&=0,
  &\Vodd^2&=0,
  \label{eq:odd-nilpotency}\\
  \Aodd^*&=-\Vodd,
  &\Aodd^{\sharp}&=\Vodd,
  \label{eq:odd-adjoints}
\end{align}
and their positive ordered returns are
\begin{equation}
  -\Vodd\Aodd
  =\begin{pmatrix}A^*A&0\\0&0\end{pmatrix},
  \qquad
  -\Aodd\Vodd
  =\begin{pmatrix}0&0\\0&AA^*\end{pmatrix}
  \label{eq:positive-returns}
\end{equation}
on their natural composition domains.
\end{proposition}

\statuslabel{Imported theorem}This is the TAUM orientation algebra, with the
abstract channel renamed \(A\) \cite{TAUM2026}.  The signs are grading data;
the positive operators are \(A^*A\) and \(AA^*\).

\subsection{Why the untyped nilpotent isometry fails}

\begin{theorem}[Nilpotent isometry obstruction]
\label{thm:nilpotent-isometry-obstruction}
Let \(T\) be a linear operator on a normed space with a \(T\)-invariant domain
\(\mathcal D\).  If
\begin{equation}
  T^2x=0
  \qquad\text{and}\qquad
  \|Tx\|=\|x\|
  \label{eq:nilpotent-isometry-hyp}
\end{equation}
for every \(x\in\mathcal D\), then \(\mathcal D=\{0\}\).
\end{theorem}

\begin{proof}
For \(x\in\mathcal D\), invariance gives \(Tx\in\mathcal D\).  Therefore
\begin{equation}
  \|x\|=\|Tx\|=\|T^2x\|=0.
\end{equation}
Hence \(x=0\).
\end{proof}

This theorem identifies the exact algebraic failure of the untyped reading:
\begin{equation}
  \Aodd^2=0,
  \qquad
  \|\Aodd\xi\|=\|\xi\|
  \quad\text{for all }\xi\in\scrK
  \label{eq:forbidden-global}
\end{equation}
cannot hold on a nontrivial common carrier.

The theorem does \emph{not} prohibit a norm-preserving map between different
spaces.  Nor does it prohibit an odd extension from preserving the norm on
one selected sector before becoming zero on the other.

\subsection{Typed isometry and nilpotent extension}

\begin{proposition}[Typed isometry with nilpotent odd extension]
\label{prop:typed-isometry}
Let \(A\in\mathcal B(\scrHP,\scrHF)\) be an isometry:
\begin{equation}
  A^*A=\idop_{\Psi}.
  \label{eq:A-isometry}
\end{equation}
Then its odd extension \(\Aodd\) is nilpotent on \(\scrK\), while for every
\(\psi\in\scrHP\),
\begin{equation}
  \|\Aodd\iP\psi\|_{\scrK}
  =\|A\psi\|_{\scrHF}
  =\|\psi\|_{\scrHP}.
  \label{eq:sector-isometry}
\end{equation}
\end{proposition}

\begin{proof}
Nilpotency follows from the off-diagonal block form.  The norm identity follows
from \(A^*A=\idop_{\Psi}\):
\begin{equation}
  \|A\psi\|^2
  =\inner{\psi}{A^*A\psi}
  =\|\psi\|^2.
\end{equation}
\end{proof}

Thus the conjunction
\begin{equation}
  \phi=A\psi,
  \qquad
  \|\phi\|=\|\psi\|
  \label{eq:original-conjunction}
\end{equation}
is mathematically sound when \(A\) is a typed isometric embedding.  Its odd
extension is nilpotent because it acts trivially on the phenomenal sector.

\begin{corollary}[Typed/odd separation]
\label{cor:typed-odd-separation}
The typed map \(A:\scrHP\to\scrHF\) and the odd extension
\(\Aodd:\scrK\to\scrK\) cannot be treated as one untyped operator carrying
the same norm-preservation requirement on the whole carrier.
\end{corollary}

\statuslabel{Derived proposition}The typed and carrier-level statements have
different domains and different algebraic roles.

\subsection{Strict phenomenal excess without two-sided inversion}

\begin{proposition}[Isometric asymmetry]
\label{prop:isometric-asymmetry}
Let \(A\in\mathcal B(\scrHP,\scrHF)\) satisfy \(A^*A=\idop_{\Psi}\).  Then
\begin{equation}
  AA^*=P_{\Ran A},
  \label{eq:range-projection}
\end{equation}
where \(\Ran A\) is closed.  Moreover,
\begin{equation}
  \scrHF=\Ran A\oplus\Ker A^*.
  \label{eq:phenomenal-decomposition}
\end{equation}
Consequently, \(A\) is a two-sided unitary equivalence if and only if
\(\Ker A^*=\{0\}\).
\end{proposition}

\begin{proof}
An isometry has closed range.  The operator \(AA^*\) is self-adjoint and
idempotent, acts as the identity on \(\Ran A\), and vanishes on
\((\Ran A)^\perp=\Ker A^*\).  Hence it is the orthogonal range projection.
\end{proof}

This theorem gives the mature form of the intuition that the phenomenal space
may be strictly larger:
\begin{equation}
  A^*A=\idop_{\Psi},
  \qquad
  AA^*<\idop_{\Phi}.
  \label{eq:strict-asymmetry}
\end{equation}
Every physical state has a norm-preserving phenomenal image, while the
phenomenal complement \(\Ker A^*\) has no physical preimage under the
isometric channel.

\warninglabel The subspace \(\Ker A^*\) is operator-theoretically present.
Its interpretation as genuine phenomenal possibility is additional
realization content, formalized only in \cref{ax:strict-phenomenal-excess}.

\subsection{The general closed channel and polar orientation}

The isometric model is too restrictive as a universal assumption.  For a
general densely defined closed channel, the standard polar decomposition
\cite{Kato,ReedSimon} gives
\begin{equation}
  A=u_A|A|,
  \qquad
  |A|=(A^*A)^{1/2},
  \label{eq:polar-decomposition}
\end{equation}
where \(u_A\) is a partial isometry from the active physical support
\begin{equation}
  \scrHP^{\mathrm{act}}:=(\Ker A)^\perp
\end{equation}
to the active phenomenal support
\begin{equation}
  \scrHF^{\mathrm{act}}:=\overline{\Ran A}.
\end{equation}

\begin{proposition}[Canonical supportwise norm correspondence]
\label{prop:polar-survival}
For all \(\psi,\chi\in\scrHP^{\mathrm{act}}\),
\begin{equation}
  \inner{u_A\psi}{u_A\chi}_{\Phi}
  =\inner{\psi}{\chi}_{\Psi}.
  \label{eq:polar-inner-product}
\end{equation}
In particular,
\begin{equation}
  \|u_A\psi\|=\|\psi\|.
  \label{eq:polar-norm}
\end{equation}
\end{proposition}

\begin{proof}
By the defining properties of the polar partial isometry,
\(u_A^*u_A\) is the support projection onto \((\Ker A)^\perp\).  Hence it acts
as the identity on \(\scrHP^{\mathrm{act}}\), and
\begin{equation}
  \inner{u_A\psi}{u_A\chi}
  =\inner{\psi}{u_A^*u_A\chi}
  =\inner{\psi}{\chi}.
\end{equation}
\end{proof}

Norm-preserving correspondence therefore belongs canonically to the
orientation \(u_A\), while \(|A|\) measures transition magnitude:
\begin{equation}
  \boxed{
  \text{Awareness channel}
  =
  \text{orientation}\times\text{magnitude}
  =u_A|A|.}
  \label{eq:orientation-magnitude}
\end{equation}

At the projective level,
\begin{equation}
  [\psi]\longmapsto[u_A\psi]
  \label{eq:projective-polar}
\end{equation}
preserves transition amplitudes on the active support.  The full Born Rule Locus Correspondence and its Krein-neutral graph are
established in \cref{sec:active-support-brlc-geometry,sec:born-rule-locus-correspondence}.

\subsection{Nilpotency and reflexive terminality}

\begin{proposition}[No nonzero fixed vector]
\label{prop:no-fixed-vector}
If \(\Aodd^2=0\), then
\begin{equation}
  \Aodd\xi=\xi
  \quad\Longrightarrow\quad
  \xi=0.
  \label{eq:no-fixed-vector}
\end{equation}
The same holds for \(\Vodd\).
\end{proposition}

\begin{proof}
Apply \(\Aodd\) to \(\Aodd\xi=\xi\):
\begin{equation}
  0=\Aodd^2\xi=\Aodd\xi=\xi.
\end{equation}
\end{proof}

For a typed physical reference \(\ket0\in\Dom A\), define
\begin{equation}
  \ket\Omega:=A\ket0.
  \label{eq:relative-luminosity}
\end{equation}
Then
\begin{equation}
  \Aodd\iF\ket\Omega=0
  \label{eq:odd-terminal}
\end{equation}
because \(\Aodd\) has no forward block from \(\scrHF\).  Equivalently,
\begin{equation}
  \Aodd^2\iP\ket0=0.
  \label{eq:iterated-awareness}
\end{equation}

\statuslabel{Imported algebra plus doctrinal identification}The algebraic
statement is nilpotency.  Its cognitive reading is that a phenomenal state
does not require a second Awareness operation to make it manifest.  This is
the no-regress component of reflexivity.  It does not yet distinguish
relative luminosity \(\ket\Omega\) from the pure luminous ground
\(\ket{\Omega_0}\).

\subsection{Three objects, not one}

The architecture fixes the following distinction:
\begin{equation}
  \boxed{
  \begin{aligned}
  A&:\Dom A\subseteq\scrHP\to\scrHF
  &&\text{typed cognitive channel},\\
  \Aodd&:\Dom\Aodd\subseteq\scrK\to\scrK
  &&\text{nilpotent odd orientation},\\
  u_A&:\scrHP^{\mathrm{act}}\to\scrHF^{\mathrm{act}}
  &&\text{norm-preserving polar orientation}.
  \end{aligned}}
  \label{eq:three-object-summary}
\end{equation}
These objects must not be conflated: their domains, algebraic laws, and
interpretive roles are different.
% =============================================================================
% END INLINED FILE: sections/02_typed_repair.tex
% =============================================================================
% =============================================================================
% BEGIN INLINED FILE: sections/03_doctrinal_map.tex
% =============================================================================
\section{Doctrinal map and interpretive constraints}
\label{sec:doctrinal-map}

\subsection{R\=upadh\=atu and dharmadh\=atu}

\begin{identification}[Fundamental sectors]
\label{id:fundamental-sectors}
The basic graded sectors are interpreted as
\begin{equation}
  \boxed{
  \scrHP\longleftrightarrow\textit{r\=upadh\=atu}\;\text{(matter)},
  \qquad
  \scrHF\longleftrightarrow\textit{dharmadh\=atu}\;\text{(phenomenal domain)}.}
  \label{eq:rupa-dharma}
\end{equation}
\end{identification}

Here \emph{r\=upa} is rendered as matter rather than the vaguer English
``form.''  The grading makes the distinction mathematically real:
\begin{equation}
  \fundJ\restr{}{\scrHP}=+\idop_{\Psi},
  \qquad
  \fundJ\restr{}{\scrHF}=-\idop_{\Phi}.
\end{equation}
The common carrier prevents the distinction from hardening into two unrelated
substances:
\begin{equation}
  \scrK=\scrHP\oplus\scrHF.
\end{equation}
The odd channel then makes the relation internal to the carrier.

The intended nonduality is therefore neither identity nor independence:
\begin{equation}
  \boxed{
  \text{difference by grading; unity by carrier and transition}.}
  \label{eq:nonduality-slogan}
\end{equation}
Matter and mind are distinguishable orientations or state types inside one
structured whole.  This is nonduality without homogenization.

\subsection{Two different dualities}

The fundamental matter--mind grading must not be confused with the
phenomenological bifurcation into grasper and grasped:
\begin{equation}
  \boxed{
  \text{matter/mind grading}
  \neq
  \text{grasper/grasped appearance}.}
  \label{eq:two-dualities}
\end{equation}
The former belongs to the carrier architecture.  The latter is a
v\=asan\=a-conditioned distortion internal to ordinary phenomenal appearance.

The dissertation reconstructs Dharmak\=irti's final position as follows:
differentiated appearances arise through latent karmic dispositions, appear
necessarily through subject--object structure, and ultimately dissolve into
the undifferentiated luminosity of reflexive awareness
\cite[pp.~27--28]{Yiannopoulos2020}.  It also emphasizes that cognition is in
reality ontologically singular and that the apparent subject--object relation
is erroneous rather than a relation between two genuinely discrete entities
\cite[pp.~40--41]{Yiannopoulos2020}.

\begin{identification}[Location of dualistic distortion]
\label{id:duality-location}
The gr\=ahya--gr\=ahaka bifurcation is not represented by the fundamental
projection pair \((\PP,\PF)\).  It is generated within \(\scrHF\) by the
algebra of phenomenal determination developed in
\cref{sec:phenomenal-determination-gns,sec:derived-duality-frames}.
\end{identification}

This constraint is load-bearing.  If subject and object were represented as
the two fundamental sectors, then the formalism would reify precisely the
duality the doctrinal analysis treats as erroneous.

\subsection{Structural sector projections}

\begin{modelchoice}[Structural-sector-projection reading]
\label{choice:ssp}
The projections
\begin{equation}
  \PP=\begin{pmatrix}\idop_{\Psi}&0\\0&0\end{pmatrix},
  \qquad
  \PF=\begin{pmatrix}0&0\\0&\idop_{\Phi}\end{pmatrix}
  \label{eq:sector-projections}
\end{equation}
are structural typing projections.  They are not assumed to be observables
whose measurement yields the mutually exclusive outcomes ``matter'' or
``mind.''
\end{modelchoice}

Thus, for a paired state \(\Xi=\psi\oplus\phi\),
\begin{equation}
  \PP\Xi=\psi\oplus0,
  \qquad
  \PF\Xi=0\oplus\phi
\end{equation}
extract the two conventional descriptions or aspects.  They do not by
themselves describe a collapse in which one aspect becomes real and the other
vanishes.

The stronger paired-observable algebra and the active neutral graph
implementing this reading are developed in
\cref{sec:active-support-brlc-geometry,sec:paired-observable-algebra,sec:ssp-factorization-phase}.

\subsection{Emptiness, vacuum, and the reference ket}

\begin{identification}[Physical vacuum and emptiness]
\label{id:vacuum-emptiness}
The distinguished physical ket \(\ket0\) is reserved for the physical vacuum
as the matter-sector representative of emptiness.
\end{identification}

The bare zero vector and the vacuum ket are categorically different:
\begin{equation}
  0\in\scrHP,
  \qquad
  \|0\|=0,
  \qquad
  \ket0\in\scrHP,
  \qquad
  \|\ket0\|=1.
  \label{eq:zero-vacuum}
\end{equation}
The zero vector is not a state.  The vacuum is a normalized state.

The physical-ground axiom imposes the cyclic condition
\begin{equation}
  \overline{\piP(\algP)\ket0}=\scrHP,
  \label{eq:physical-cyclicity-preview}
\end{equation}
which gives mathematical content to ``matter is vacuum; vacuum is matter.''
Every matter state in the selected representation is generated from the
vacuum; the vacuum itself is a matter state.  The full axiom and its
consequences appear in \cref{ax:physical-vacuum,id:matter-vacuum}.

The term emptiness does not mean nothingness.  The dissertation's opening
Mah\=amudr\=a citation states that mind is empty of its own nature and yet,
being empty and unobstructed, can appear as anything
\cite[front matter]{Yiannopoulos2020}.  The intended mathematical analogue is
a nonzero generative ground, not the zero vector.

\subsection{Relative luminosity and the luminous ground}

The typed image
\begin{equation}
  \boxed{\ket\Omega:=A\ket0}
  \label{eq:Omega-doctrinal}
\end{equation}
is interpreted as relative luminosity, or the awareness of emptiness.  Its
selection principle is relational: it is whatever the Awareness channel does
to the physical vacuum.

The symbol
\begin{equation}
  \ket{\Omega_0}
  \label{eq:Omega0-doctrinal}
\end{equation}
uses a different selection principle: the unique luminous phenomenal ground.
The phenomenal-ground axiom characterizes it intrinsically by a phenomenal
generator and cyclic representation.  The bridge
\begin{equation}
  A\ket0=c\ket{\Omega_0},
  \qquad c\neq0,
  \label{eq:ground-bridge-preview2}
\end{equation}
formalizes the claim
\begin{equation}
  \boxed{
  \text{the awareness of emptiness is the natural state}.}
  \label{eq:awareness-emptiness-natural}
\end{equation}

The notation must remain distinct because \(\ket\Omega\) and
\(\ket{\Omega_0}\) are identified by different routes.  Their projective equality is substantive content of the ground bridge, not a notational convention.

\subsection{Rang rig and algebraic reflexivity}

The dissertation describes reflexive awareness as the inherently
self-presenting nature of cognition, irrespective of the causal status of its
content \cite[p.~31]{Yiannopoulos2020}.  It also records the critical
traditional distinction between intrinsic nondual reflexivity and subsequent
definitive recognition: awareness may already apprehend cognition's nondual
nature while the seeds of error prevent that nature from being ascertained as
such \cite[p.~425]{Yiannopoulos2020}.

\begin{identification}[Nilpotency as the no-regress law]
\label{id:nilpotency-rangrig}
The equation
\begin{equation}
  \Aodd^2=0
  \label{eq:nilpotency-rangrig}
\end{equation}
is interpreted as the algebraic no-regress component of \emph{rang rig}.  A
phenomenal state does not require a second Awareness operation in order to be
manifest.
\end{identification}

This identification is deliberately limited.  Since \(\Aodd\) vanishes on
the entire phenomenal sector, nilpotency characterizes the terminality of all
phenomenal states.  It does not select the pure luminous ground.  A ground
condition is therefore supplied independently by
\cref{ax:phenomenal-ground}, which distinguishes \(\ket{\Omega_0}\) from
contentful experience.

\subsection{Gzhi, gnas lugs, rig pa, and sems}

The MA thesis argues that all cognitions are ultimately luminous even when
nondual luminosity is not ordinarily ascertained
\cite[pp.~162--163]{Yiannopoulos2012}.  In the Four Yogas, differentiated
appearances disappear at the Third Yoga, leaving nondual and undifferentiated
luminosity, while the reflexive nature of awareness remains the same through
the progression \cite[pp.~182--183]{Yiannopoulos2012}.  The Fourth Yoga is
nonappearance, not a second state below luminosity
\cite[pp.~184--185]{Yiannopoulos2012}.

The realization adopts the following assignments.

\begin{identification}[Ground dictionary]
\label{id:ground-dictionary}
\begin{align}
  \textit{gzhi},\ \textit{gnas lugs}
  &\longleftrightarrow
  \text{the luminous ground ray or its full ground structure},
  \label{eq:gzhi-map}\\
  \ket{\Omega_0}
  &\longleftrightarrow
  \text{a normalized vector representative of that ground},
  \label{eq:Omega0-map}\\
  \textit{rig pa}
  &\longleftrightarrow
  \text{the realization of the ground},
  \label{eq:rigpa-map}\\
  \textit{sems}
  &\longleftrightarrow
  \text{differentiated phenomenal articulation or excitation}.
  \label{eq:sems-map}
\end{align}
\end{identification}

The rank-one ground projection and its orthogonal complement supply the exact
conditions
\begin{equation}
  \PzeroF\phi=\phi
  \quad\Longleftrightarrow\quad
  \textit{rig pa},
  \qquad
  (\idop-\PzeroF)\phi\neq0
  \quad\Longleftrightarrow\quad
  \textit{sems}.
  \label{eq:rigpa-sems-preview}
\end{equation}

\warninglabel ``Contentless'' must remain qualified.  Pure luminosity lacks
differentiated \=ak\=ara-content and subject--object marks; it is not the blank
zero vector or a phenomenal absence.  The Third and Fourth Yogas therefore do
not receive distinct kets such as \(\ket{\Omega_{-1}}\).

\subsection{The algebra of appearance}

The generators of the phenomenal determination algebra encode
v\=asan\=a-conditioned operations producing
\begin{enumerate}[label=(\roman*)]
  \item differentiated cognitive images or \=ak\=aras;
  \item conceptual determination, memory, and concatenation;
  \item affective and perspectival variation;
  \item the erroneous grasper--grasped structure.
\end{enumerate}
This is doctrinally motivated by the dissertation's account of ordinary
appearance as internally generated by defiled dispositions and necessarily
structured through dualistic distortion \cite[pp.~27--28]{Yiannopoulos2020}.

The generator system and its operator stages are specified in
\cref{sec:vasana-generator-system,sec:operator-stages-appearance}.
Their structural location is fixed by
\begin{equation}
  \boxed{
  \text{the algebra of dualistic appearance acts inside }\scrHF,
  \text{ not between }\scrHP\text{ and }\scrHF.}
  \label{eq:appearance-location}
\end{equation}

\subsection{Concordance table}

\begin{center}
\small
\begin{tabularx}{\textwidth}{>{\raggedright\arraybackslash}p{0.19\textwidth}
>{\raggedright\arraybackslash}p{0.29\textwidth}
>{\raggedright\arraybackslash}X}
\toprule
\textbf{Term} & \textbf{Mathematical role} & \textbf{Constraint or status} \\
\midrule
\emph{r\=upadh\=atu} & \(\scrHP\), physical or matter sector & Not the grasper pole of cognition. \\
\emph{dharmadh\=atu} & \(\scrHF\), phenomenal sector & Not identical with differentiated content alone. \\
\emph{\'s\=unyat\=a} & \(\ket0\) and ultimately its cyclic vacuum structure & Not the bare zero vector. \\
\emph{'od gsal}/luminosity & Ground and shared phenomenal manifestness & Not represented by adding \(\ket{\Omega_0}\) to every state. \\
\emph{rang rig} & Reflexive terminality, algebraically shadowed by \(\Aodd^2=0\) & Nilpotency alone does not select the ground. \\
\emph{gzhi}/\emph{gnas lugs} & Luminous ground ray or full pointed ground structure & Represented by the unique cyclic ground structure. \\
\emph{rig pa} & Exact realization of the ground & Not partial ground overlap. \\
\emph{sems} & Differentiated phenomenal articulation or excitation & A projective condition, not a linear sector. \\
\emph{gr\=ahya--gr\=ahaka} & Derived internal bifurcation of appearance & Must not be identified with the fundamental grading. \\
\bottomrule
\end{tabularx}
\end{center}
% =============================================================================
% END INLINED FILE: sections/03_doctrinal_map.tex
% =============================================================================
% =============================================================================
% BEGIN INLINED FILE: sections/05_pointed_ground_representations.tex
% =============================================================================
\section{Pointed cyclic representations and ground-state data}
\label{sec:pointed-ground-representations}

The typed transition channel, its odd extension, and its polar orientation are
mathematically distinct.  The additional structures below make it possible to
speak precisely about a physical vacuum, a luminous phenomenal ground, and the
relation between their vector, ray, projection, and state-functional
representations.

\subsection{Represented algebras and pointed representations}

\begin{definition}[Pointed representation]
\label{def:pointed-representation}
Let $\mathfrak M$ be a unital $C^*$-algebra, let
\begin{equation}
  \pi:\mathfrak M\longrightarrow\mathcal B(\mathscr H)
\end{equation}
be a nondegenerate $*$-representation, and let $\Omega\in\mathscr H$ be a
unit vector.  The quadruple
\begin{equation}
  (\mathfrak M,\pi,\mathscr H,\Omega)
\end{equation}
is a \emph{pointed representation}.  It is \emph{cyclic} when
\begin{equation}
  \overline{\pi(\mathfrak M)\Omega}=\mathscr H.
  \label{eq:generic-cyclicity}
\end{equation}
\end{definition}

Because $\mathfrak M$ is a linear algebra containing the identity,
$\pi(\mathfrak M)\Omega$ is already a linear subspace.  The closure in
\eqref{eq:generic-cyclicity} is therefore the cyclic subspace generated by
$\Omega$.  The phrase ``orbit of the vacuum'' will sometimes be used
informally, but ``cyclic subspace generated by the vacuum'' is the more
standard expression.  A group orbit alone generally does not generate the
whole state space.

\begin{definition}[Vector state]
\label{def:vector-state}
The unit vector $\Omega$ determines the state functional
\begin{equation}
  \omega_{\Omega}(X)
  :=\inner{\Omega}{\pi(X)\Omega},
  \qquad X\in\mathfrak M.
  \label{eq:generic-vector-state}
\end{equation}
\end{definition}

\begin{proposition}[Positivity and normalization]
\label{prop:vector-state-positive}
The functional $\omega_{\Omega}$ is linear, positive, and normalized:
\begin{equation}
  \omega_{\Omega}(X^*X)\geq0,
  \qquad
  \omega_{\Omega}(\idop)=1.
\end{equation}
\end{proposition}

\begin{proof}
Linearity is immediate.  Moreover,
\begin{equation}
  \omega_{\Omega}(X^*X)
  =\|\pi(X)\Omega\|^2\geq0,
\end{equation}
and nondegeneracy together with unitality gives
$\pi(\idop)=\idop_{\mathscr H}$, so
$\omega_{\Omega}(\idop)=\|\Omega\|^2=1$.
\end{proof}

The state functional is invariant under a phase change of the representing
vector:
\begin{equation}
  \omega_{e^{i\theta}\Omega}=\omega_{\Omega}.
\end{equation}
Thus a vector representative is never the intrinsic object by itself.  The
same vector state and Hilbert-space direction are represented by the ray
\begin{equation}
  [\Omega]=\CC\Omega
\end{equation}
or, after normalization, by the rank-one projection
\begin{equation}
  P_{\Omega}=\ket\Omega\!\bra\Omega.
\end{equation}

\subsection{Ground generators and unique ground rays}

\begin{definition}[Ground datum]
\label{def:ground-datum}
Let $H$ be self-adjoint and bounded below on $\mathscr H$, and let
\begin{equation}
  E_0:=\inf\spec(H).
\end{equation}
A unit vector $\Omega_0\in\Dom H$ is a \emph{ground vector} when
\begin{equation}
  H\Omega_0=E_0\Omega_0.
\end{equation}
The ground ray is \emph{unique} when
\begin{equation}
  \Ker(H-E_0\idop)=\CC\Omega_0.
  \label{eq:generic-unique-ground}
\end{equation}
Equivalently, the spectral projection
\begin{equation}
  P_0:=E_H(\{E_0\})
\end{equation}
has rank one.
\end{definition}

The qualifier ``unique'' belongs to the ray, not to a particular phase choice.
If $\Omega_0$ is normalized, then every other normalized ground vector is
$e^{i\theta}\Omega_0$.  The different phases do not encode differences among
sentient beings, observers, or cognitive continua.  A family of distinct
continua would require a family of representations or ground fibers; no such
family-level structure is assumed here.

\begin{remark}[Uniqueness, cyclicity, and a gap are independent]
\label{rem:ground-properties-independent}
A unique ground ray need not be cyclic for a given algebra, and a cyclic
vector need not be a unique ground vector.  Likewise, uniqueness does not
imply a positive spectral gap.  One may have a one-dimensional zero-energy
eigenspace while continuous spectrum accumulates immediately above it.
Accordingly, the following are treated as separate hypotheses:
\begin{equation}
  \rank P_0=1,
  \qquad
  \overline{\pi(\mathfrak M)\Omega_0}=\mathscr H,
  \qquad
  H\geq E_0\idop+\Delta(\idop-P_0).
\end{equation}
Only the first two are adopted as basic ground axioms; the gap estimate is
optional.
The underlying spectral distinctions are standard; see
\cite{Kato,ReedSimon}.
\end{remark}

\subsection{GNS canonicality}

Every state $\omega$ on a unital $C^*$-algebra has a Gelfand--Naimark--Segal
triple
\begin{equation}
  (\pi_{\omega},\mathscr H_{\omega},\Omega_{\omega})
\end{equation}
with
\begin{equation}
  \omega(X)
  =\inner{\Omega_{\omega}}{\pi_{\omega}(X)\Omega_{\omega}},
  \qquad
  \overline{\pi_{\omega}(\mathfrak M)\Omega_{\omega}}
  =\mathscr H_{\omega}.
\end{equation}
The triple is unique up to a unitary intertwiner carrying one cyclic vector
to the other \cite{BratteliRobinson}.

\begin{proposition}[Cyclic realization as a GNS realization]
\label{prop:cyclic-is-gns}
Let $(\mathfrak M,\pi,\mathscr H,\Omega)$ be a cyclic pointed
representation, and let $\omega_{\Omega}$ be its vector state.  Then the
pointed representation is unitarily equivalent to the GNS representation of
$\omega_{\Omega}$.
\end{proposition}

This is the standard sense in which a state functional canonically determines
a cyclic vector representation.  The vector is canonical only relative to the
state and up to unitary equivalence and phase.  No claim of numerical identity
across distinct representations follows.

\subsection{Technical purity and faithfulness}

\begin{definition}[Pure and faithful algebraic states]
\label{def:pure-faithful-state}
A state $\omega$ on $\mathfrak M$ is \emph{pure} when it is an extreme point
of the convex state space.  Equivalently, its GNS representation is
irreducible.  It is \emph{faithful} when
\begin{equation}
  \omega(X^*X)=0
  \quad\Longrightarrow\quad
  X=0.
\end{equation}
For the concrete represented algebra $\pi(\mathfrak M)$, the vector state is
faithful exactly when $\Omega$ is separating for $\pi(\mathfrak M)$.  Its
pullback to the abstract algebra $\mathfrak M$ is faithful when, in addition,
the representation $\pi$ is faithful.
\end{definition}

\warninglabel The doctrinal expression ``pure luminosity'' does not by itself
assert purity in the technical convex-geometric sense of
\cref{def:pure-faithful-state}.  Technical purity means non-mixture; it does
not mean absence of differentiated content.  Likewise, cyclicity, purity,
faithfulness, ground uniqueness, and dynamical invariance are distinct
properties.  The present realization assumes cyclicity and unique ground
rays, but does not assume that either ground state functional is technically
pure or faithful.

\subsection{Zero is not a state}

The bare zero vector and the zero functional fail normalization:
\begin{equation}
  \|0\|=0,
  \qquad
  0(\idop)=0.
\end{equation}
Neither is a state.  This provides a precise asymmetry that will matter on the
phenomenal side: pure luminosity may be a state, whereas complete absence of
phenomenal instantiation is not a competing phenomenal state.
% =============================================================================
% END INLINED FILE: sections/05_pointed_ground_representations.tex
% =============================================================================
% =============================================================================
% BEGIN INLINED FILE: sections/06_physical_cyclic_vacuum.tex
% =============================================================================
\section{The physical cyclic vacuum}
\label{sec:physical-cyclic-vacuum}

The physical vacuum is not the zero vector and is not assumed to be a zero
mode of every operator in the transition architecture.  It is selected by a
physical ground generator and by the cyclic representation of a matter
algebra.

\subsection{Physical ground datum}

\begin{axiom}[Unique cyclic physical vacuum]
\label{ax:physical-vacuum}
Let $\algP$ be a unital $C^*$-algebra with nondegenerate representation
\begin{equation}
  \piP:\algP\longrightarrow\mathcal B(\scrHP).
\end{equation}
Let $H_{\Psi}$ be self-adjoint and bounded below, with
\begin{equation}
  E_0^{\Psi}:=\inf\spec(H_{\Psi}).
\end{equation}
There exists a normalized vector $\ket0\in\Dom H_{\Psi}$ such that
\begin{align}
  \Ker(H_{\Psi}-E_0^{\Psi}\idop_{\Psi})
  &=\CC\ket0,
  \label{eq:physical-unique-ground}\\
  \overline{\piP(\algP)\ket0}
  &=\scrHP.
  \label{eq:physical-cyclic-vacuum}
\end{align}
\end{axiom}

Define
\begin{equation}
  \PzeroP:=\ket0\!\bra0,
  \qquad
  \QzeroP:=\idop_{\Psi}-\PzeroP.
  \label{eq:physical-ground-projections}
\end{equation}
The invariant physical ground object is the ray
\begin{equation}
  \mathscr G_{\Psi}:=\Ran\PzeroP=\CC\ket0.
\end{equation}
The ket $\ket0$ is a normalized representative of this ray.

\statuslabel{Realization axiom} Neither the existence of $H_{\Psi}$ nor the
rank-one and cyclicity conditions in \cref{ax:physical-vacuum} follow from the
bare TAUM channel algebra.  They are additional data of the cognitive
realization.

\subsection{The physical vacuum state functional}

The vacuum vector determines
\begin{equation}
  \boxed{
  \omegaP(X)
  :=\bra0\piP(X)\ket0,
  \qquad X\in\algP.}
  \label{eq:physical-vacuum-functional}
\end{equation}
By \cref{prop:vector-state-positive}, $\omegaP$ is a state.  Because the
vacuum is cyclic, the pointed representation
\begin{equation}
  (\algP,\piP,\scrHP,\ket0)
  \label{eq:physical-pointed-representation}
\end{equation}
is unitarily equivalent to the GNS representation of $\omegaP$.

The notation $\omegaP$ is standard operator-algebraic notation.  It should
not be confused with the vector $\ket0$, the ray $\mathscr G_{\Psi}$, or the
projection $\PzeroP$.  These encode the same selected ground state at
different mathematical levels:
\begin{equation}
  \boxed{
  \omegaP
  \quad\longleftrightarrow\quad
  (\piP,\scrHP,\ket0)
  \quad\longleftrightarrow\quad
  \mathscr G_{\Psi}
  \quad\longleftrightarrow\quad
  \PzeroP.}
  \label{eq:physical-ground-levels}
\end{equation}
The arrows denote representation correspondence, not literal equality of a
functional, a vector, a subspace, and an operator.

\subsection{Cyclic generation of matter states}

\begin{definition}[Physical cyclic subspace]
\label{def:physical-cyclic-subspace}
For a vector $\psi\in\scrHP$, define
\begin{equation}
  \Cyc_{\piP(\algP)}(\psi)
  :=\overline{\piP(\algP)\psi}.
\end{equation}
\end{definition}

The vacuum axiom is precisely
\begin{equation}
  \Cyc_{\piP(\algP)}(\ket0)=\scrHP.
\end{equation}
Thus every physical state in the selected vacuum sector can be approximated
arbitrarily well by vectors of the form
\begin{equation}
  \piP(X)\ket0,
  \qquad X\in\algP.
\end{equation}
No claim is made that every vector is obtained from one preferred operator or
one symmetry-group action.  The full represented algebra supplies the cyclic
generation.

\begin{identification}[Matter is vacuum; vacuum is matter]
\label{id:matter-vacuum}
The cyclicity equation
\begin{equation}
  \overline{\piP(\algP)\ket0}=\scrHP
\end{equation}
is the physical content assigned to the traditional formula ``matter is
emptiness; emptiness is matter.''  Matter states are determinations or
excitations of the cyclic vacuum, and the vacuum itself belongs to the matter
sector.
\end{identification}

The identification does not collapse all matter vectors into one ket.  It
states a generative relation:
\begin{equation}
  \boxed{
  \text{the physical representation is generated from its empty ground}.}
  \label{eq:physical-generative-relation}
\end{equation}

\subsection{Centered matter operations}

For $X\in\algP$, define its vacuum-centered part by
\begin{equation}
  X^{\circ}_{\Psi}
  :=X-\omegaP(X)\idop.
  \label{eq:physical-centered-operator}
\end{equation}
Then
\begin{equation}
  \bra0\piP(X^{\circ}_{\Psi})\ket0=0,
\end{equation}
so
\begin{equation}
  \piP(X)\ket0
  =\omegaP(X)\ket0
  +\piP(X^{\circ}_{\Psi})\ket0,
  \label{eq:physical-centered-decomposition}
\end{equation}
with
\begin{equation}
  \piP(X^{\circ}_{\Psi})\ket0\perp\ket0.
\end{equation}

\begin{proposition}[Vacuum and excitation decomposition]
\label{prop:physical-centered-cyclic}
Under \cref{ax:physical-vacuum},
\begin{equation}
  \QzeroP\scrHP
  =\overline{
  \left\{
    \piP(X^{\circ}_{\Psi})\ket0:
    X\in\algP
  \right\}}.
  \label{eq:physical-centered-span}
\end{equation}
\end{proposition}

\begin{proof}
Applying $\QzeroP$ to \eqref{eq:physical-centered-decomposition} gives
\begin{equation}
  \QzeroP\piP(X)\ket0
  =\piP(X^{\circ}_{\Psi})\ket0.
\end{equation}
The uncentered cyclic orbit is dense in $\scrHP$ by
\eqref{eq:physical-cyclic-vacuum}; applying the bounded projection
$\QzeroP$ and taking closures yields \eqref{eq:physical-centered-span}.
\end{proof}

This proposition gives a clean physical decomposition:
\begin{equation}
  \scrHP
  =\mathscr G_{\Psi}\oplus\QzeroP\scrHP,
\end{equation}
where the second summand is generated by centered matter operations acting on
the vacuum.

\subsection{Uniqueness and symmetry}

Suppose a unitary representation $U(g)$ of a symmetry group commutes with
$H_{\Psi}$.  Then $U(g)$ preserves the ground eigenspace.  Because that
eigenspace is one-dimensional,
\begin{equation}
  U(g)\ket0=\chi(g)\ket0,
  \qquad |\chi(g)|=1.
  \label{eq:vacuum-symmetry-phase}
\end{equation}
Thus a unique vacuum ray is automatically invariant under every symmetry of
the ground generator, although a normalized representative may acquire a
phase.  This phase freedom is gauge freedom of the ray; it is not a model of
individual differences among sentient beings.

\begin{openquestion}[Choice of matter algebra]
\label{q:matter-algebra}
The algebra $\algP$ is supplied realization data.  Taking
$\algP=\mathcal B(\scrHP)$ would make every nonzero vector cyclic and would
therefore trivialize much of the content of \cref{ax:physical-vacuum}.  A
concrete physical realization must specify the smaller field, observable, or
matter-generating algebra whose action genuinely builds the vacuum sector.
\end{openquestion}
% =============================================================================
% END INLINED FILE: sections/06_physical_cyclic_vacuum.tex
% =============================================================================
% =============================================================================
% BEGIN INLINED FILE: sections/07_luminous_cyclic_ground.tex
% =============================================================================
\section{The luminous phenomenal ground}
\label{sec:luminous-cyclic-ground}

The phenomenal side has the same formal shape as the physical vacuum
representation, but a different doctrinal role.  The luminous ground is the
intrinsic phenomenal reference from which determinate appearances are
algebraically generated.

\subsection{Phenomenal ground datum}

\begin{axiom}[Unique cyclic luminous ground]
\label{ax:phenomenal-ground}
Let $\algF$ be a unital $C^*$-algebra with nondegenerate representation
\begin{equation}
  \piF:\algF\longrightarrow\mathcal B(\scrHF).
\end{equation}
Let $H_{\Phi}$ be self-adjoint and bounded below, with
\begin{equation}
  E_0^{\Phi}:=\inf\spec(H_{\Phi}).
\end{equation}
There exists a normalized vector $\ket{\Omega_0}\in\Dom H_{\Phi}$ such that
\begin{align}
  \Ker(H_{\Phi}-E_0^{\Phi}\idop_{\Phi})
  &=\CC\ket{\Omega_0},
  \label{eq:phenomenal-unique-ground}\\
  \overline{\piF(\algF)\ket{\Omega_0}}
  &=\scrHF.
  \label{eq:phenomenal-cyclic-ground}
\end{align}
\end{axiom}

Define
\begin{equation}
  \PzeroF:=\ket{\Omega_0}\!\bra{\Omega_0},
  \qquad
  \QzeroF:=\idop_{\Phi}-\PzeroF,
  \label{eq:phenomenal-ground-projections}
\end{equation}
and
\begin{equation}
  \mathscr G_{\Phi}:=\Ran\PzeroF=\CC\ket{\Omega_0}.
  \label{eq:phenomenal-ground-ray}
\end{equation}
The ray $\mathscr G_{\Phi}$ is the intrinsic ground object; the ket
$\ket{\Omega_0}$ is a normalized representative.

\statuslabel{Realization axiom} The unique cyclic luminous ground is not
implied by the TAUM chassis or by the existence of a nonzero Awareness image.
It is additional structure selected by $H_{\Phi}$ and $\algF$.

\subsection{The luminous ground state functional}

The vector $\ket{\Omega_0}$ determines
\begin{equation}
  \boxed{
  \omegaF(X)
  :=\bra{\Omega_0}\piF(X)\ket{\Omega_0},
  \qquad X\in\algF.}
  \label{eq:phenomenal-ground-functional}
\end{equation}
The cyclic pointed representation
\begin{equation}
  (\algF,\piF,\scrHF,\ket{\Omega_0})
  \label{eq:phenomenal-pointed-representation}
\end{equation}
is unitarily equivalent to the GNS representation of $\omegaF$.

The standard notation is therefore
\begin{equation}
  \omegaF
  \quad\text{for the phenomenal ground-state functional,}
\end{equation}
and
\begin{equation}
  \ket{\Omega_0}
  \quad\text{for its normalized cyclic ground representative.}
\end{equation}
The subscript $0$ carries substantive information: it marks the state selected
intrinsically as the phenomenal ground.  It should not ordinarily be
suppressed in favor of $\ket\Omega$, which is selected by the transition
channel.

\subsection{All phenomenal states are determinations of luminosity}

\begin{identification}[Cyclic luminosity]
\label{id:cyclic-luminosity}
The cyclicity equation
\begin{equation}
  \overline{\piF(\algF)\ket{\Omega_0}}=\scrHF
  \label{eq:cyclic-luminosity}
\end{equation}
is interpreted as the claim that every phenomenal state is an articulation,
determination, or excitation of luminosity.
\end{identification}

The statement is generative rather than additive.  A phenomenal state need
not have the form
\begin{equation}
  \ket{\Omega_0}+\ket{\text{content}}.
\end{equation}
Nor need it have nonzero overlap with $\ket{\Omega_0}$.  Its luminosity lies
in its membership in the cyclic representation generated from the luminous
ground.

This formal distinction matches the doctrinal distinction between the
luminosity intrinsic to all cognition and the explicit realization or
ascertainment of that luminosity.  The MA thesis states that all cognitions
are always ultimately luminous even though nondual luminosity is not
ordinarily ascertained as such \cite[pp.~162--163]{Yiannopoulos2012}.  The
dissertation likewise distinguishes the nondual reflexive nature already
present in cognition from the subsequent determination that cognition is
nondual \cite[p.~425]{Yiannopoulos2020}.

\subsection{No autonomous dark phenomenal sector}

\begin{definition}[Dark complement relative to the luminous cyclic orbit]
\label{def:dark-complement}
Define
\begin{equation}
  \scrHF^{\mathrm{dark}}
  :=
  \left(
  \overline{\piF(\algF)\ket{\Omega_0}}
  \right)^{\perp}.
  \label{eq:dark-complement}
\end{equation}
\end{definition}

\begin{proposition}[Absence of a separate non-luminous sector]
\label{prop:no-dark-sector}
Under \cref{ax:phenomenal-ground},
\begin{equation}
  \scrHF^{\mathrm{dark}}=\{0\}.
\end{equation}
\end{proposition}

\begin{proof}
The cyclicity equation \eqref{eq:phenomenal-cyclic-ground} identifies the
closed subspace in parentheses in \eqref{eq:dark-complement} with all of
$\scrHF$.  Its orthogonal complement is therefore zero.
\end{proof}

This proposition does not assert that every phenomenal vector has positive
ground overlap.  It excludes only an independent phenomenal subspace that
cannot be generated from luminosity at all.  A vector orthogonal to the ground
can still be a centered determination of the ground.

Complete absence of luminosity is represented neither by a second normalized
ket nor by a second state functional.  The bare zero vector and zero
functional are not states.  Thus the formalism contains the asymmetry
\begin{equation}
  \boxed{
  \begin{gathered}
  \text{pure luminosity is a phenomenal state,}\\
  \text{while pure absence of luminosity is no phenomenal state.}
  \end{gathered}}
  \label{eq:no-nonluminosity-state}
\end{equation}

\subsection{Centered phenomenal determinations}

For $X\in\algF$, define
\begin{equation}
  X^{\circ}_{\Phi}
  :=X-\omegaF(X)\idop.
  \label{eq:phenomenal-centered-operator}
\end{equation}
Then
\begin{equation}
  \bra{\Omega_0}\piF(X^{\circ}_{\Phi})\ket{\Omega_0}=0,
\end{equation}
and
\begin{equation}
  \piF(X)\ket{\Omega_0}
  =\omegaF(X)\ket{\Omega_0}
  +\piF(X^{\circ}_{\Phi})\ket{\Omega_0},
  \label{eq:phenomenal-centered-decomposition}
\end{equation}
with
\begin{equation}
  \piF(X^{\circ}_{\Phi})\ket{\Omega_0}
  \perp
  \ket{\Omega_0}.
\end{equation}

\begin{proposition}[Ground-orthogonal phenomenal sector]
\label{prop:phenomenal-centered-cyclic}
Under \cref{ax:phenomenal-ground},
\begin{equation}
  \QzeroF\scrHF
  =\overline{
  \left\{
    \piF(X^{\circ}_{\Phi})\ket{\Omega_0}:
    X\in\algF
  \right\}}.
  \label{eq:phenomenal-centered-span}
\end{equation}
\end{proposition}

\begin{proof}
Applying $\QzeroF$ to \eqref{eq:phenomenal-centered-decomposition} gives
\begin{equation}
  \QzeroF\piF(X)\ket{\Omega_0}
  =\piF(X^{\circ}_{\Phi})\ket{\Omega_0}.
\end{equation}
The cyclic orbit is dense in $\scrHF$; applying the bounded projection
$\QzeroF$ and taking closures gives \eqref{eq:phenomenal-centered-span}.
\end{proof}

This decomposition supplies the algebraic input for the sems construction in
\cref{sec:sems-exhaustion-ground-recovery}:
\begin{equation}
  \scrHF
  =\mathscr G_{\Phi}\oplus\QzeroF\scrHF.
  \label{eq:phenomenal-ground-excitation-split}
\end{equation}
The relevant centered generators are the v\=asan\=a-conditioned operations of
\cref{sec:vasana-generator-system}, including the production of \=ak\=aras,
memory, conceptual concatenation, affective differentiation, and
grasper--grasped distortion.  The generic centered decomposition alone does
not assign a doctrinal interpretation to every element of $\algF$.

\subsection{Technical purity remains open}

\begin{openquestion}[Operator-algebraic purity of the luminous state]
\label{q:luminous-purity}
Is $\omegaF$ pure in the technical sense of
\cref{def:pure-faithful-state}?  Equivalently, is its GNS representation
irreducible?  The doctrinal expression ``pure luminosity'' does not decide
this question.  Both technical purity and faithfulness remain open.
\end{openquestion}

For example, on the full algebra $\mathcal B(\scrHF)$ with
$\dim\scrHF>1$, the vector state of $\ket{\Omega_0}$ is pure but not
faithful.  A concrete realization may distinguish a global phenomenal algebra
from smaller determination algebras on which the ground vector is cyclic and
separating.

\begin{openquestion}[Choice of phenomenal algebra]
\label{q:phenomenal-algebra}
The algebra $\algF$ is supplied realization data.  It must be large enough
for $\ket{\Omega_0}$ to be cyclic but not so large that cyclicity becomes
vacuous.  The selected determination layer and its generator system are given
in \cref{sec:phenomenal-determination-gns,sec:vasana-generator-system} without
forcing a unique global choice of $\algF$.
\end{openquestion}
% =============================================================================
% END INLINED FILE: sections/07_luminous_cyclic_ground.tex
% =============================================================================
% =============================================================================
% BEGIN INLINED FILE: sections/08_ground_fidelity_excitation.tex
% =============================================================================
\section{Ground fidelity, orthogonality, and excitation}
\label{sec:ground-fidelity}

The two unique ground rays determine canonical decompositions of the physical
and phenomenal Hilbert spaces.  This section isolates what ground overlap does
and does not mean.

\subsection{Ground fidelity}

Let $S\in\{\Psi,\Phi\}$, with Hilbert space $\mathscr H_S$, normalized ground
vector $\Omega_{0,S}$, and rank-one ground projection $P_0^S$.  Thus
\begin{equation}
  \Omega_{0,\Psi}=\ket0,
  \qquad
  \Omega_{0,\Phi}=\ket{\Omega_0}.
\end{equation}
For a normalized vector $\xi\in\mathscr H_S$, define the ground fidelity
\begin{equation}
  \boxed{
  F_0^S(\xi)
  :=\inner{\xi}{P_0^S\xi}
  =|\inner{\Omega_{0,S}}{\xi}|^2.}
  \label{eq:ground-fidelity-general}
\end{equation}
Then
\begin{equation}
  0\leq F_0^S(\xi)\leq1.
\end{equation}
Moreover,
\begin{equation}
  F_0^S(\xi)=1
  \quad\Longleftrightarrow\quad
  [\xi]=[\Omega_{0,S}],
  \label{eq:fidelity-one-ray}
\end{equation}
and
\begin{equation}
  F_0^S(\xi)=0
  \quad\Longleftrightarrow\quad
  \xi\perp\Omega_{0,S}.
  \label{eq:fidelity-zero-orthogonal}
\end{equation}

The canonical decomposition is
\begin{equation}
  \xi=P_0^S\xi+Q_0^S\xi,
  \qquad
  Q_0^S:=\idop-P_0^S.
  \label{eq:ground-orthogonal-decomposition}
\end{equation}
If $F_0^S(\xi)<1$, define
\begin{equation}
  \eta_{\xi}
  :=\frac{Q_0^S\xi}{\sqrt{1-F_0^S(\xi)}}.
  \label{eq:conditional-excitation-vector}
\end{equation}
Then $\eta_{\xi}\perp\Omega_{0,S}$ and $\|\eta_{\xi}\|=1$.  If
$F_0^S(\xi)>0$, a phase $\alpha$ may be chosen so that
\begin{equation}
  \xi
  =e^{i\alpha}\sqrt{F_0^S(\xi)}\,\Omega_{0,S}
  +\sqrt{1-F_0^S(\xi)}\,\eta_{\xi}.
  \label{eq:fidelity-state-decomposition}
\end{equation}
At $F_0^S(\xi)=0$, the first term vanishes and no phase is required.  At
$F_0^S(\xi)=1$, the conditional vector $\eta_{\xi}$ is not defined and is not
needed.

\subsection{Exact excitation-energy factorization}

Let $H_S$ be the corresponding ground generator, bounded below by
$E_0^S$, with
\begin{equation}
  \Ker(H_S-E_0^S\idop)=\Ran P_0^S.
\end{equation}
For a normalized vector in the quadratic-form domain, define
\begin{equation}
  E_{\mathrm{exc}}^S(\xi)
  :=\left\|
  (H_S-E_0^S\idop)^{1/2}\xi
  \right\|^2.
  \label{eq:excitation-energy-general}
\end{equation}

\begin{proposition}[Ground-fidelity factorization]
\label{prop:ground-fidelity-factorization}
If $F_0^S(\xi)<1$, then
\begin{equation}
  \boxed{
  E_{\mathrm{exc}}^S(\xi)
  =
  \bigl(1-F_0^S(\xi)\bigr)
  Z_S(\xi),}
  \label{eq:ground-fidelity-factorization}
\end{equation}
where
\begin{equation}
  Z_S(\xi)
  :=\left\|
  (H_S-E_0^S\idop)^{1/2}\eta_{\xi}
  \right\|^2.
  \label{eq:conditional-excitation-energy}
\end{equation}
If $F_0^S(\xi)=1$, then $E_{\mathrm{exc}}^S(\xi)=0$.
\end{proposition}

\begin{proof}
The spectral projection $P_0^S$ reduces $H_S$, and
$(H_S-E_0^S\idop)^{1/2}P_0^S=0$.  Hence
\begin{equation}
  E_{\mathrm{exc}}^S(\xi)
  =\left\|
  (H_S-E_0^S\idop)^{1/2}Q_0^S\xi
  \right\|^2.
\end{equation}
Substituting
$Q_0^S\xi=\sqrt{1-F_0^S(\xi)}\,\eta_{\xi}$ gives
\eqref{eq:ground-fidelity-factorization}.  If $F_0^S(\xi)=1$, then
$Q_0^S\xi=0$ and the excitation energy vanishes.
\end{proof}

\begin{corollary}[Gap estimate]
\label{cor:ground-gap-fidelity}
If, in the quadratic-form sense,
\begin{equation}
  H_S-E_0^S\idop
  \geq
  \Delta_S Q_0^S
  \qquad
  \text{for some }\Delta_S>0,
  \label{eq:ground-gap-form}
\end{equation}
then
\begin{equation}
  E_{\mathrm{exc}}^S(\xi)
  \geq
  \Delta_S\bigl(1-F_0^S(\xi)\bigr).
  \label{eq:ground-gap-bound}
\end{equation}
\end{corollary}

\statuslabel{Imported theorem} The factorization and gap estimate are the
two-ground specialization of the reference-fidelity result already established
in TAUM \cite{TAUM2026}.  The cognitive realization assigns the physical and
phenomenal ground roles.

\subsection{Phenomenal orthogonality does not mean non-luminosity}

For a normalized phenomenal state $\phi\in\scrHF$,
\begin{equation}
  F_0^{\Phi}(\phi)
  =|\braket{\Omega_0}{\phi}|^2
  \label{eq:phenomenal-ground-fidelity}
\end{equation}
measures overlap with the undifferentiated luminous ground.  It does not
measure the amount of luminosity present in the state.

A state with
\begin{equation}
  F_0^{\Phi}(\phi)=0
\end{equation}
is fully ground-orthogonal, but by phenomenal cyclicity it remains in the
closure of states generated from $\ket{\Omega_0}$.  Under
\cref{prop:phenomenal-centered-cyclic}, it can be approximated by centered
determinations
\begin{equation}
  \piF(X^{\circ}_{\Phi})\ket{\Omega_0}.
\end{equation}
Thus exact orthogonality is compatible with luminosity as a structural
property of the representation.

\begin{modelchoice}[No positive lower bound on ordinary ground fidelity]
\label{choice:no-fidelity-lower-bound}
No axiom of the form
\begin{equation}
  F_0^{\Phi}(\phi)\geq\varepsilon>0
\end{equation}
for all normalized phenomenal states.  Exact ground-orthogonal states are
mathematically permitted.  Whether a specific cognitive realization ever
attains exact orthogonality rather than arbitrarily small positive fidelity
remains open.
\end{modelchoice}

The doctrinal reading is therefore:
\begin{equation}
  \boxed{
  \begin{aligned}
  F_0^{\Phi}(\phi)=1
  &\quad\text{means exact ground realization},\\
  0<F_0^{\Phi}(\phi)<1
  &\quad\text{means ground and differentiated components coexist},\\
  F_0^{\Phi}(\phi)=0
  &\quad\text{means ground orthogonality, not absence of luminosity}.
  \end{aligned}}
  \label{eq:phenomenal-fidelity-reading}
\end{equation}

Exact ground realization is identified with \emph{rig pa} in
\cref{id:rigpa-density}.  Accordingly, $F_0^{\Phi}$ is called ground
fidelity, not a degree of rig pa.

\subsection{Physical ground fidelity}

Likewise,
\begin{equation}
  F_0^{\Psi}(\psi)
  =|\braket{0}{\psi}|^2
\end{equation}
measures overlap with the physical vacuum.  It does not measure the amount of
matter in a state.  Matter is represented by membership in the vacuum sector
and cyclic generation from $\ket0$, while $F_0^{\Psi}$ measures the vacuum
component relative to the chosen ground decomposition.

\subsection{The grounds are state types, not enduring substances}

The equations of this section are compatible with Buddhist momentariness.
The rays $\mathscr G_{\Psi}$ and $\mathscr G_{\Phi}$ specify canonical ground
state types within a representation.  They do not by themselves posit an
enduring substantial self or one numerically persistent mental object through
time.  Family- and continuum-level interpretations require additional
realization data and are not assumed in this section.
% =============================================================================
% END INLINED FILE: sections/08_ground_fidelity_excitation.tex
% =============================================================================
% =============================================================================
% BEGIN INLINED FILE: sections/09_ground_bridge.tex
% =============================================================================
\section{Awareness of emptiness and the luminous ground bridge}
\label{sec:ground-bridge}

The two cyclic-ground axioms do not yet relate the grounds.  The relation is
supplied by the typed Awareness channel.  This section adopts the weak
projective bridge and separates it from the stronger polar bridge.  The weak
bridge canonically determines the restricted ground orientation developed in
\cref{sec:ground-channel-orientation}; the stronger condition controls only its
embedding into the full polar geometry.

\subsection{Relative and intrinsic luminosity}

Assume throughout this section that
\begin{equation}
  \ket0\in\Dom A,
  \qquad
  A\ket0\neq0.
  \label{eq:vacuum-in-A-domain}
\end{equation}
Define the relative phenomenal image
\begin{equation}
  \boxed{\ket\Omega:=A\ket0.}
  \label{eq:relative-Omega}
\end{equation}
The vector $\ket\Omega$ is selected relationally by the channel $A$.  The
vector $\ket{\Omega_0}$ is selected intrinsically by the phenomenal ground
data in \cref{ax:phenomenal-ground}.  These are different selection
principles even when they select the same ray.

\begin{axiom}[Ground bridge]
\label{ax:ground-bridge}
The Awareness image of the physical vacuum lies on the intrinsic luminous
ground ray:
\begin{equation}
  \boxed{
  A\ket0=c\ket{\Omega_0},
  \qquad c\in\CC\setminus\{0\}.}
  \label{eq:ground-bridge-axiom}
\end{equation}
Equivalently,
\begin{equation}
  [\ket\Omega]=[\ket{\Omega_0}].
  \label{eq:ground-bridge-rays}
\end{equation}
\end{axiom}

\statuslabel{Realization axiom} The bridge is not implied by the existence of
the channel or by either cyclic-ground axiom separately.  It is the central
cognitive realization condition connecting emptiness and luminosity.

Because the phase of $\ket{\Omega_0}$ is free, it may be chosen so that
\begin{equation}
  c=a_0>0,
  \qquad
  a_0:=\|A\ket0\|.
  \label{eq:ground-bridge-phase-fix}
\end{equation}
After that phase convention,
\begin{equation}
  \ket\Omega=a_0\ket{\Omega_0}.
  \label{eq:Omega-scaled-Omega0}
\end{equation}
The scalar $a_0$ measures the raw transition magnitude on the physical
ground.  It does not create a second phenomenal state.

\subsection{Equivalent rank-one formulations}

Because $\Ran\PzeroP=\CC\ket0\subseteq\Dom A$, the operator
\begin{equation}
  A\PzeroP:\scrHP\longrightarrow\scrHF
\end{equation}
is bounded and rank one even when $A$ itself is unbounded.  Under the phase
choice \eqref{eq:ground-bridge-phase-fix},
\begin{equation}
  A\PzeroP
  =a_0\ket{\Omega_0}\!\bra0.
  \label{eq:ground-channel-rank-one}
\end{equation}

\begin{proposition}[Equivalent ground-bridge identities]
\label{prop:ground-bridge-equivalences}
Under \cref{ax:ground-bridge}, the following identities hold:
\begin{align}
  \PzeroF A\PzeroP
  &=A\PzeroP,
  \label{eq:bridge-compression}\\
  \QzeroF A\PzeroP
  &=0,
  \label{eq:bridge-no-excitation}\\
  (A\PzeroP)(A\PzeroP)^*
  &=a_0^2\PzeroF,
  \label{eq:bridge-rank-one-product}\\
  \PzeroF
  &=\frac{\ket\Omega\!\bra\Omega}{\braket{\Omega}{\Omega}}
   =\frac{(A\PzeroP)(A\PzeroP)^*}{\|A\ket0\|^2}.
  \label{eq:bridge-projector-formula}
\end{align}
Conversely, \eqref{eq:bridge-no-excitation} together with
$A\ket0\neq0$ implies that $A\ket0$ lies on the ground ray.
\end{proposition}

\begin{proof}
Equation \eqref{eq:ground-channel-rank-one} gives all four displayed
identities directly.  Conversely, \eqref{eq:bridge-no-excitation} gives
$A\ket0=\PzeroF A\ket0\in\Ran\PzeroF=\CC\ket{\Omega_0}$, and nonvanishing
gives a nonzero proportionality scalar.
\end{proof}

The most invariant form of the bridge is the projector identity
\eqref{eq:bridge-projector-formula}.  It removes phase and normalization
conventions.  The most transparent doctrinal form remains
\eqref{eq:ground-bridge-axiom}.

\begin{identification}[Awareness of emptiness]
\label{id:awareness-of-emptiness}
The bridge
\begin{equation}
  A\ket0\propto\ket{\Omega_0}
\end{equation}
is interpreted as
\begin{equation}
  \boxed{
  \text{the awareness of emptiness is the natural luminous state}.}
  \label{eq:awareness-of-emptiness-slogan}
\end{equation}
The vector $\ket\Omega$ is relative luminosity because it is identified by
its relation to $A$ and $\ket0$; the vector $\ket{\Omega_0}$ is ultimate
luminosity because it is identified by the intrinsic phenomenal ground
structure.  The two names track the two routes to one ray.
\end{identification}

This is why the subscript should ordinarily remain visible.  Suppressing it
would conceal the fact that projective coincidence is substantive content of
the realization rather than a definition.

\subsection{The strong polar ground bridge}

The weak bridge does not by itself imply that $\ket0$ is a singular vector of
$A$ or that the polar orientation sends the physical ground directly to the
phenomenal ground.  That requires an additional condition.

\begin{modelchoice}[Strong polar ground compatibility]
\label{choice:strong-polar-ground}
One may impose
\begin{equation}
  \ket0\in\Dom(A^*A),
  \qquad
  A^*A\ket0=a_0^2\ket0.
  \label{eq:strong-polar-ground}
\end{equation}
This condition is not adopted as a basic realization axiom.  It remains a
separate model choice governing compatibility with the full polar geometry.
\end{modelchoice}

\begin{proposition}[Consequences of strong polar compatibility]
\label{prop:strong-polar-consequences}
Assume \cref{ax:ground-bridge,choice:strong-polar-ground}, with the phase
fixed by \eqref{eq:ground-bridge-phase-fix}.  Then
\begin{align}
  |A|\ket0
  &=a_0\ket0,
  \label{eq:absolute-A-ground}\\
  u_A\ket0
  &=\ket{\Omega_0},
  \label{eq:polar-ground-map}\\
  A^*\ket{\Omega_0}
  &=a_0\ket0,
  \label{eq:adjoint-ground-return}\\
  AA^*\ket{\Omega_0}
  &=a_0^2\ket{\Omega_0}.
  \label{eq:phenomenal-return-ground}
\end{align}
\end{proposition}

\begin{proof}
Equation \eqref{eq:strong-polar-ground} and functional calculus give
$|A|\ket0=a_0\ket0$.  The polar decomposition
$A=u_A|A|$, together with $A\ket0=a_0\ket{\Omega_0}$, gives
$u_A\ket0=\ket{\Omega_0}$.  Since
$\ket0\in\Dom(A^*A)$, the vector $A\ket0=a_0\ket{\Omega_0}$ lies in
$\Dom A^*$; dividing
$A^*A\ket0=a_0^2\ket0$ by $a_0$ gives
\eqref{eq:adjoint-ground-return}.  Applying $A$ then gives
\eqref{eq:phenomenal-return-ground}.
\end{proof}

These are standard consequences of polar decomposition and positive
functional calculus; see \cite{Kato,ReedSimon}.

Under the strong condition, the ground bridge is no longer merely a
coincidence of rays.  It is part of the canonical polar orientation of the
full channel.  This full-channel statement is distinct from the canonical
polar orientation of the restricted rank-one ground channel, which already
follows from the weak bridge.

\subsection{The return-kernel obstruction and its correct reading}

\begin{proposition}[A nonzero Awareness image is not a zero return mode]
\label{prop:vacuum-not-return-zero-mode}
If $\ket0\in\Dom(A^*A)$ and $A\ket0\neq0$, then
\begin{equation}
  A^*A\ket0\neq0.
\end{equation}
In particular,
\begin{equation}
  A^*A\ket0=0
  \quad\Longrightarrow\quad
  A\ket0=0.
  \label{eq:return-kernel-obstruction}
\end{equation}
\end{proposition}

\begin{proof}
For $\ket0\in\Dom(A^*A)$,
\begin{equation}
  \bra0 A^*A\ket0=\|A\ket0\|^2.
\end{equation}
If $A^*A\ket0=0$, the right-hand side vanishes.
\end{proof}

This does not conflict with calling $\ket0$ the vacuum.  The label $0$ means
zero excitation relative to the physical ground generator $H_{\Psi}$ after a
choice of energy origin.  It does not mean zero eigenvalue for every positive
return operator in the architecture.

\subsection{Return operators are not automatically Hamiltonians}

The structural positive returns inherited from TAUM \cite{TAUM2026} are
\begin{equation}
  R_{\Psi}:=A^*A,
  \qquad
  R_{\Phi}:=AA^*.
  \label{eq:structural-returns}
\end{equation}
They should not be silently identified with $H_{\Psi}$ and $H_{\Phi}$.

If one imposed
\begin{equation}
  H_{\Psi}=A^*A,
  \qquad
  H_{\Psi}\ket0=0,
\end{equation}
then \cref{prop:vacuum-not-return-zero-mode} would force
$A\ket0=0$, contradicting the ground bridge.  A shifted active-support model
may instead use
\begin{equation}
  H_{\Psi}=A^*A-a_0^2\idop
  \label{eq:shifted-physical-return-model}
\end{equation}
provided $a_0^2$ is genuinely the lowest relevant return value.

The phenomenal side is more delicate.  If $\Ker A^*\neq\{0\}$, then
$AA^*$ has zero modes in the phenomenal excess, while
$\ket{\Omega_0}$ has return value $a_0^2>0$ under the strong polar bridge.
Therefore $AA^*$ cannot automatically be the global phenomenal Hamiltonian
for which $\ket{\Omega_0}$ is the unique ground.

\begin{modelchoice}[Independent ground generators]
\label{choice:independent-ground-generators}
The realization keeps
\begin{equation}
  H_{\Psi},\ H_{\Phi}
\end{equation}
separate from
\begin{equation}
  A^*A,\ AA^*.
\end{equation}
Relations between them may be added in a concrete realization, but no global
identity is assumed.
\end{modelchoice}

\subsection{Normalized relative luminosity}

Under the weak bridge,
\begin{equation}
  \widehat\Omega
  :=\frac{\ket\Omega}{\|\ket\Omega\|}
  =e^{i\theta}\ket{\Omega_0}
\end{equation}
for some phase.  Hence $\widehat\Omega$ and $\ket{\Omega_0}$ determine the
same projector and the same phenomenal vector state:
\begin{equation}
  \proj{\widehat\Omega}=\PzeroF,
  \qquad
  \bra{\widehat\Omega}\piF(X)\ket{\widehat\Omega}
  =\omegaF(X).
\end{equation}
Again, this equivalence holds because the bridge has been imposed.  It is not
a consequence of notation alone.
% =============================================================================
% END INLINED FILE: sections/09_ground_bridge.tex
% =============================================================================
% =============================================================================
% BEGIN INLINED FILE: sections/11_ground_channel_orientation.tex
% =============================================================================
\section{The restricted ground channel and its canonical orientation}
\label{sec:ground-channel-orientation}

Assume the weak bridge
\begin{equation}
  A\ket0=a_0\ket{\Omega_0},
  \qquad
  a_0=\|A\ket0\|>0,
  \label{eq:ground-bridge-fixed}
\end{equation}
after fixing the phase of \(\ket{\Omega_0}\).  The question is whether this
weak bridge already supplies a canonical norm-preserving ground orientation,
or whether the stronger singular-vector condition from
\cref{choice:strong-polar-ground} is required.

The answer is precise: the weak bridge is sufficient on the two ground rays.
The stronger condition is needed only if the ground orientation is to be the
restriction of the full polar orientation \(u_A\).

\subsection{The rank-one ground channel}

Define the physical and phenomenal ground spaces
\begin{equation}
  \GzeroP:=\Ran\PzeroP=\CC\ket0,
  \qquad
  \GzeroF:=\Ran\PzeroF=\CC\ket{\Omega_0}.
  \label{eq:ground-spaces}
\end{equation}
Because \(\GzeroP\subseteq\Dom A\), the restricted operator
\begin{equation}
  \boxed{\Azero:=A\PzeroP:\scrHP\longrightarrow\scrHF}
  \label{eq:ground-channel-definition}
\end{equation}
is bounded even when \(A\) is unbounded.

\begin{lemma}[Bounded rank-one ground channel]
\label{lem:ground-channel-bounded}
Under \cref{ax:ground-bridge},
\begin{equation}
  \Azero=a_0\ket{\Omega_0}\!\bra0,
  \qquad
  \|\Azero\|=a_0,
  \qquad
  \rank\Azero=1.
  \label{eq:restricted-ground-channel-rank-one}
\end{equation}
Moreover,
\begin{equation}
  \Azero^*\Azero=a_0^2\PzeroP,
  \qquad
  \Azero\Azero^*=a_0^2\PzeroF.
  \label{eq:ground-channel-products}
\end{equation}
\end{lemma}

\begin{proof}
For \(\psi\in\scrHP\),
\begin{equation}
  \Azero\psi
  =A\ket0\braket{0}{\psi}
  =a_0\ket{\Omega_0}\braket{0}{\psi}.
\end{equation}
The remaining identities follow from normalization of the two ground
vectors.
\end{proof}

\subsection{The canonical ground partial isometry}

\begin{definition}[Ground orientation]
\label{def:ground-orientation}
Define
\begin{equation}
  \boxed{
  \uzero:=a_0^{-1}\Azero
  =\ket{\Omega_0}\!\bra0.}
  \label{eq:u0-definition}
\end{equation}
Thus \(\uzero\) is regarded as a bounded operator from \(\scrHP\) to
\(\scrHF\), vanishing on \(\GzeroP^\perp\).
\end{definition}

\begin{theorem}[Polar decomposition of the ground channel]
\label{thm:ground-channel-polar}
The operator \(\uzero\) is the unique partial isometry satisfying
\begin{equation}
  \uzero^*\uzero=\PzeroP,
  \qquad
  \uzero\uzero^*=\PzeroF,
  \qquad
  \Azero=a_0\uzero.
  \label{eq:u0-support-identities}
\end{equation}
The polar decomposition of the bounded ground channel is
\begin{equation}
  \boxed{
  \Azero=\uzero|\Azero|,
  \qquad
  |\Azero|=a_0\PzeroP.}
  \label{eq:A0-polar-decomposition}
\end{equation}
Restricted to \(\GzeroP\),
\begin{equation}
  \uzero:\GzeroP\longrightarrow\GzeroF
  \label{eq:u0-ground-unitary}
\end{equation}
is unitary and satisfies
\begin{equation}
  \uzero\ket0=\ket{\Omega_0},
  \qquad
  \uzero^*\ket{\Omega_0}=\ket0.
  \label{eq:u0-ground-actions}
\end{equation}
\end{theorem}

\begin{proof}
The support identities follow immediately from
\eqref{eq:u0-definition}.  By \cref{lem:ground-channel-bounded},
\begin{equation}
  |\Azero|
  =(\Azero^*\Azero)^{1/2}
  =a_0\PzeroP,
\end{equation}
and therefore
\begin{equation}
  \uzero|\Azero|
  =a_0\uzero\PzeroP
  =a_0\uzero
  =\Azero.
\end{equation}
Uniqueness is the uniqueness of the partial-isometric factor in the polar
decomposition of a bounded operator.  The final claims follow because the
initial and final spaces are one-dimensional.
\end{proof}

\statuslabel{Derived theorem} No stronger assumption on \(A^*A\) is required
to obtain a canonical orientation between the two ground rays.  The bridge
itself makes \(A\PzeroP\) a rank-one operator, and that restricted operator
has its own canonical polar decomposition.
This is the rank-one instance of standard polar decomposition
\cite{Kato,ReedSimon}.

\subsection{Phase independence}

The vector formulas in \eqref{eq:u0-definition} appear to use phase choices,
but the operator \(\uzero\) does not depend on them once the bridge is fixed.
Indeed, the invariant formula
\begin{equation}
  \boxed{
  \uzero=\frac{A\PzeroP}{\|A\ket0\|}}
  \label{eq:u0-invariant-formula}
\end{equation}
uses only the channel and the physical ground projection.  If
\begin{equation}
  \ket0\longmapsto e^{i\alpha}\ket0,
  \qquad
  \ket{\Omega_0}\longmapsto e^{i\alpha}\ket{\Omega_0},
\end{equation}
then \(\PzeroP\), \(\PzeroF\), and \(\uzero\) are unchanged.  The common
phase changes vector representatives but not the bridge-selected ground
relation.

\begin{proposition}[Uniqueness of the bridge-selected ground orientation]
\label{prop:u0-uniqueness}
Let \(w:\scrHP\to\scrHF\) be a partial isometry with initial projection
\(\PzeroP\), final projection \(\PzeroF\), and
\begin{equation}
  A\PzeroP=a_0 w.
\end{equation}
Then \(w=\uzero\).
\end{proposition}

\begin{proof}
Division by \(a_0>0\) gives \(w=a_0^{-1}A\PzeroP=\uzero\).
\end{proof}

Thus the weak bridge selects not only a pair of rays but also their relative
phase as encoded by the restricted channel.

\subsection{Weak bridge versus full polar compatibility}

Write the polar decomposition of the full closed channel as
\begin{equation}
  A=\polA\absA.
  \label{eq:full-A-polar}
\end{equation}
The restricted ground orientation \(\uzero\) need not equal
\(\polA\PzeroP\).  The reason is that \(\absA\ket0\) need not lie on the
physical ground ray even though \(A\ket0\) lies on the phenomenal ground
ray.

\begin{theorem}[Equivalent forms of strong polar ground compatibility]
\label{thm:strong-ground-equivalences}
Assume \cref{ax:ground-bridge} and the phase convention
\eqref{eq:ground-bridge-fixed}.  The following statements are
equivalent:
\begin{enumerate}[label=(\roman*)]
  \item \(\ket0\in\Dom(A^*A)\) and
  \begin{equation}
    A^*A\ket0=a_0^2\ket0;
  \end{equation}
  \item
  \begin{equation}
    \absA\ket0=a_0\ket0;
  \end{equation}
  \item
  \begin{equation}
    \polA\ket0=\ket{\Omega_0};
  \end{equation}
  \item
  \begin{equation}
    \polA\PzeroP=\uzero;
  \end{equation}
  \item \(\ket{\Omega_0}\in\Dom A^*\) and
  \begin{equation}
    A^*\ket{\Omega_0}=a_0\ket0.
  \end{equation}
\end{enumerate}
\end{theorem}

\begin{proof}
The equivalence of (i) and (ii) follows from positive functional calculus.
If (ii) holds, then
\begin{equation}
  a_0\polA\ket0
  =\polA\absA\ket0
  =A\ket0
  =a_0\ket{\Omega_0},
\end{equation}
which gives (iii).  Statements (iii) and (iv) are equivalent because both
operators vanish on \(\GzeroP^\perp\).  Statement (i) implies (v) by applying
\(A^*\) to \(A\ket0=a_0\ket{\Omega_0}\).  Conversely, (v) gives
\begin{equation}
  A^*A\ket0
  =a_0A^*\ket{\Omega_0}
  =a_0^2\ket0.
\end{equation}
Finally, if (iii) holds, then \(\|\polA\ket0\|=1\), so \(\ket0\) lies in
the initial support of \(\polA\).  Both \(\absA\ket0\) and
\(a_0\ket0\) lie in that support and have the same image under the isometry
\(\polA\); hence they are equal, giving (ii).
\end{proof}

\begin{remark}[Scope of the strong-bridge condition]
\label{rem:strong-ground-compatibility}
The weak bridge is sufficient to construct the canonical ground orientation
\(\uzero\) and the neutral ground line.  Strong polar compatibility is needed
only for the additional claim
\begin{equation}
  \uzero=\polA\PzeroP,
\end{equation}
which makes the ground construction a literal slice of the full active-support
polar correspondence.  The strong condition therefore remains optional rather
than a basic axiom.
\end{remark}
% =============================================================================
% END INLINED FILE: sections/11_ground_channel_orientation.tex
% =============================================================================
% =============================================================================
% BEGIN INLINED FILE: sections/12_neutral_ground_line.tex
% =============================================================================
\section{The Krein-neutral nondual ground line}
\label{sec:neutral-ground-line}

The ground orientation \(\uzero\) relates two distinct one-dimensional
sectoral grounds.  This section packages that relation as a line in the total
Krein carrier.  The graph and neutrality language is standard in indefinite
inner-product geometry; see \cite{Bognar}.

\subsection{The ground plane}

Define the ground support projection and ground plane by
\begin{equation}
  \Pizero
  :=\begin{pmatrix}\PzeroP&0\\0&\PzeroF\end{pmatrix},
  \qquad
  \Kzero:=\Ran\Pizero=\GzeroP\oplus\GzeroF.
  \label{eq:ground-plane-definition}
\end{equation}
The restriction of the fundamental symmetry to \(\Kzero\) has signature
\((1,1)\):
\begin{equation}
  \restr{\fundJ}{\Kzero}
  =\begin{pmatrix}1&0\\0&-1\end{pmatrix}
  \label{eq:J-ground-plane}
\end{equation}
relative to the normalized basis \(\ket0,\ket{\Omega_0}\).

\begin{definition}[Aligned and anti-aligned ground graphs]
\label{def:ground-graph-lines}
Define
\begin{align}
  \NzeroPlus
  &:=\Graph(\uzero|_{\GzeroP})
    =\{\psi\oplus\uzero\psi:\psi\in\GzeroP\},
  \label{eq:N0-plus-definition}\\
  \NzeroMinus
  &:=\Graph(-\uzero|_{\GzeroP})
    =\{\psi\oplus(-\uzero\psi):\psi\in\GzeroP\}.
  \label{eq:N0-minus-definition}
\end{align}
The line \(\NzeroPlus\) is the bridge-aligned ground graph.  The line
\(\NzeroMinus\) is its canonical anti-aligned companion.
\end{definition}

\subsection{Total neutrality}

\begin{theorem}[Neutrality of the ground graphs]
\label{thm:ground-graphs-neutral}
For either sign \(\varepsilon\in\{+1,-1\}\),
\begin{equation}
  \kInner{\xi}{\eta}=0
  \qquad
  \text{for all }\xi,\eta\in\mathscr N_0^{\varepsilon}.
  \label{eq:ground-total-neutrality}
\end{equation}
Thus \(\NzeroPlus\) and \(\NzeroMinus\) are totally neutral linear
subspaces of \(\Kzero\).  Each is maximal neutral in the ground plane.
\end{theorem}

\begin{proof}
Let
\begin{equation}
  \xi=\psi\oplus\varepsilon\uzero\psi,
  \qquad
  \eta=\chi\oplus\varepsilon\uzero\chi.
\end{equation}
Because \(\uzero:\GzeroP\to\GzeroF\) is unitary,
\begin{align}
  \kInner{\xi}{\eta}
  &=\inner{\psi}{\chi}_{\Psi}
    -\inner{\varepsilon\uzero\psi}{\varepsilon\uzero\chi}_{\Phi}\\
  &=\inner{\psi}{\chi}_{\Psi}-\inner{\psi}{\chi}_{\Psi}=0.
\end{align}
A neutral subspace of the two-dimensional nondegenerate plane \(\Kzero\)
cannot have dimension two, so each one-dimensional neutral graph is maximal.
\end{proof}

\statuslabel{Derived theorem} The ground is neutral because the bridge-selected
orientation is isometric.  Neutrality means equality of sector norms in the
indefinite metric; it does not mean that the total vector is zero or has zero
Hilbert norm.

\subsection{Normalized relational ground vectors}

Define Hilbert-isometric embeddings
\begin{align}
  \WzeroPlus\psi
  &:=\frac{1}{\sqrt2}\bigl(\psi\oplus\uzero\psi\bigr),
  \label{eq:W0-plus-definition}\\
  \WzeroMinus\psi
  &:=\frac{1}{\sqrt2}\bigl(\psi\oplus(-\uzero\psi)\bigr),
  \label{eq:W0-minus-definition}
\end{align}
for \(\psi\in\GzeroP\).  Their ranges are \(\NzeroPlus\) and
\(\NzeroMinus\), respectively.

\begin{proposition}[Normalized ground pair]
\label{prop:normalized-ground-pair}
The vectors
\begin{align}
  \ket{\Xi_0^+}
  &:=\WzeroPlus\ket0
   =\frac1{\sqrt2}\bigl(\ket0\oplus\ket{\Omega_0}\bigr),
  \label{eq:Xi0-plus}\\
  \ket{\Xi_0^-}
  &:=\WzeroMinus\ket0
   =\frac1{\sqrt2}\bigl(\ket0\oplus(-\ket{\Omega_0})\bigr)
  \label{eq:Xi0-minus}
\end{align}
satisfy
\begin{align}
  \|\Xi_0^+\|_{\scrK}
  &=\|\Xi_0^-\|_{\scrK}=1,
  \label{eq:Xi-Hilbert-normalization}\\
  \kInner{\Xi_0^+}{\Xi_0^+}
  &=\kInner{\Xi_0^-}{\Xi_0^-}=0,
  \label{eq:Xi-Krein-null}\\
  \|\PP\Xi_0^{\pm}\|^2
  &=\|\PF\Xi_0^{\pm}\|^2=\frac12.
  \label{eq:Xi-equal-sector-weights}
\end{align}
\end{proposition}

\begin{proof}
All identities follow from normalization of \(\ket0\) and
\(\ket{\Omega_0}\), orthogonality of the two Hilbert summands, and the
opposite signs in the Krein form.
\end{proof}

The factor \(2^{-1/2}\) normalizes the total Hilbert vector.  It does not
reduce either sectoral state after structural compression and
renormalization:
\begin{equation}
  \frac{\PP\ket{\Xi_0^+}}{\|\PP\Xi_0^+\|}=\ket0,
  \qquad
  \frac{\PF\ket{\Xi_0^+}}{\|\PF\Xi_0^+\|}=\ket{\Omega_0}.
  \label{eq:Xi-conditional-sector-states}
\end{equation}

\subsection{Nonduality without identity}

\begin{proposition}[Distinct sectors, one graph relation]
\label{prop:nonduality-without-identity}
The aligned ground line satisfies
\begin{equation}
  \NzeroPlus\cap\scrHP=\{0\},
  \qquad
  \NzeroPlus\cap\scrHF=\{0\}.
  \label{eq:N0-trivial-sector-intersections}
\end{equation}
Nevertheless, the restricted projections
\begin{equation}
  \restr{\PP}{\NzeroPlus}:\NzeroPlus\to\GzeroP,
  \qquad
  \restr{\PF}{\NzeroPlus}:\NzeroPlus\to\GzeroF
  \label{eq:N0-sector-projection-isomorphisms}
\end{equation}
are linear isomorphisms.  Either sectoral component determines the entire
relational ground vector through \(\uzero\) or \(\uzero^*\).
\end{proposition}

\begin{proof}
If \(\psi\oplus\uzero\psi\) lies in \(\scrHP\), then
\(\uzero\psi=0\).  Since \(\uzero\) is unitary on \(\GzeroP\), this gives
\(\psi=0\).  The phenomenal case is identical.  The projection claims are
immediate from the graph definition.
\end{proof}

This is the exact ground-level form of nonduality without identity:
\begin{equation}
  \boxed{
  \begin{aligned}
  &\text{the physical and phenomenal grounds are neither identical}\\
  &\text{nor independent; they are projections of one graph relation.}
  \end{aligned}}
  \label{eq:nonduality-without-identity-slogan}
\end{equation}

\begin{identification}[The relational nondual ground]
\label{id:relational-nondual-ground}
Under the structural-sector-projection reading, the ray
\begin{equation}
  [\ket{\Xi_0^+}]=\projective{\NzeroPlus}
  \label{eq:relational-ground-ray}
\end{equation}
is interpreted as one relational ground state in the common carrier.  Its
structural projections yield the physical vacuum and luminous ground as the
two conventional sectoral descriptions.
\end{identification}

\statuslabel{Doctrinal identification and model choice} The total vector is
not interpreted as a premeasurement superposition that collapses into either
matter or mind.  The equal sector weights in
\eqref{eq:Xi-equal-sector-weights} express balanced representation in the
common Hilbert topology.  The sector projections remain structural
compressions, not assumed matter-or-mind observables.

\subsection{Ground cyclicity remains sectoral}

The total ground relation does not replace either cyclic representation.
The claims remain
\begin{equation}
  \overline{\piP(\algP)\ket0}=\scrHP,
  \qquad
  \overline{\piF(\algF)\ket{\Omega_0}}=\scrHF.
  \label{eq:sectoral-cyclicity-retained}
\end{equation}
No algebra acting cyclically on \(\ket{\Xi_0^+}\) in the full carrier is
assumed here.  The paired observable algebra associated with the active graph
is constructed in
\cref{sec:active-support-brlc-geometry,sec:born-rule-locus-correspondence,sec:paired-observable-algebra}.
% =============================================================================
% END INLINED FILE: sections/12_neutral_ground_line.tex
% =============================================================================
% =============================================================================
% BEGIN INLINED FILE: sections/13_ground_reflection_phase.tex
% =============================================================================
\section{Ground reflection, complex structure, and relative phase}
\label{sec:ground-reflection-phase}

The bridge-selected orientation produces the ground-level reflection and
quarter-turn already latent in the polar geometry of the universal chassis.
It also shows why Krein neutrality alone does not select a unique ground
ray.

\subsection{The ground polar operators}

\begin{definition}[Ground reflection and quarter-turn]
\label{def:ground-polar-operators}
Define bounded operators on \(\scrK\) by
\begin{equation}
  \boxed{
  \Rzero
  :=\begin{pmatrix}0&\uzero^*\\\uzero&0\end{pmatrix},
  \qquad
  \Izero:=\fundJ\Rzero,
  \qquad
  \Qzero:=-i\Izero.}
  \label{eq:ground-polar-operators}
\end{equation}
All three vanish on \(\Kzero^\perp\).
\end{definition}

\begin{theorem}[Ground polar algebra]
\label{thm:ground-polar-algebra}
The operators in \cref{def:ground-polar-operators} satisfy
\begin{align}
  \Rzero^*&=\Rzero,
  &\Rzero^2&=\Pizero,
  \label{eq:R0-properties}\\
  \fundJ\Rzero&=-\Rzero\fundJ,
  &\Izero^*&=-\Izero,
  \label{eq:J-R0-anticommute}\\
  \Izero^2&=-\Pizero,
  &\Qzero^*&=\Qzero,
  \label{eq:I0-Q0-properties}\\
  \Qzero^2&=\Pizero.
  \label{eq:Q0-square}
\end{align}
On the ground plane, the restrictions of
\(\fundJ,\Rzero,\Qzero\) form a Pauli triple:
\begin{equation}
  \fundJ\Rzero=i\Qzero,
  \qquad
  \Rzero\Qzero=i\fundJ,
  \qquad
  \Qzero\fundJ=i\Rzero,
  \label{eq:ground-Pauli-products}
\end{equation}
and the three operators anticommute pairwise.
\end{theorem}

\begin{proof}
The support identities \(\uzero^*\uzero=\PzeroP\) and
\(\uzero\uzero^*=\PzeroF\) give
\(\Rzero^2=\Pizero\).  The remaining block computations are immediate.
On \(\Kzero\), \(\Pizero\) is the identity, and the matrices relative to
\(\ket0,\ket{\Omega_0}\) are
\begin{equation}
  \fundJ=\begin{pmatrix}1&0\\0&-1\end{pmatrix},
  \qquad
  \Rzero=\begin{pmatrix}0&1\\1&0\end{pmatrix},
  \qquad
  \Qzero=\begin{pmatrix}0&-i\\i&0\end{pmatrix}.
\end{equation}
These are the standard Pauli matrices \(\sigma_z,\sigma_x,\sigma_y\).
\end{proof}

\subsection{The two reflection eigenlines}

\begin{proposition}[Ground graph eigenspaces]
\label{prop:ground-graph-eigenspaces}
On \(\Kzero\),
\begin{equation}
  \NzeroPlus=\Ker(\Rzero-\idop),
  \qquad
  \NzeroMinus=\Ker(\Rzero+\idop).
  \label{eq:N0-R0-eigenspaces}
\end{equation}
In particular,
\begin{equation}
  \Rzero\ket{\Xi_0^+}=\ket{\Xi_0^+},
  \qquad
  \Rzero\ket{\Xi_0^-}=-\ket{\Xi_0^-}.
  \label{eq:R0-Xi-eigenvalues}
\end{equation}
Moreover,
\begin{align}
  \fundJ\ket{\Xi_0^+}&=\ket{\Xi_0^-},
  &\fundJ\ket{\Xi_0^-}&=\ket{\Xi_0^+},
  \label{eq:J-exchanges-Xi}\\
  \Izero\ket{\Xi_0^+}&=\ket{\Xi_0^-},
  &\Izero\ket{\Xi_0^-}&=-\ket{\Xi_0^+}.
  \label{eq:I0-quarter-turn-Xi}
\end{align}
\end{proposition}

\begin{proof}
For \(\psi\in\GzeroP\),
\begin{align}
  \Rzero(\psi\oplus\uzero\psi)
  &=\uzero^*\uzero\psi\oplus\uzero\psi
   =\psi\oplus\uzero\psi,\\
  \Rzero(\psi\oplus-\uzero\psi)
  &=-\psi\oplus\uzero\psi
   =-(\psi\oplus-\uzero\psi).
\end{align}
The remaining identities follow from the definitions.
\end{proof}

The pair \(\ket{\Xi_0^+},\ket{\Xi_0^-}\) is Hilbert orthonormal but Krein
hyperbolic:
\begin{align}
  \inner{\Xi_0^+}{\Xi_0^-}&=0,
  \label{eq:Xi-Hilbert-orthogonal}\\
  \kInner{\Xi_0^+}{\Xi_0^-}&=1,
  \label{eq:Xi-Krein-cross-pairing}
\end{align}
with both diagonal Krein norms zero.  Thus the two neutral lines form a
normalized Witt pair for the ground plane:
\begin{equation}
  \Kzero=\NzeroPlus\oplus\NzeroMinus.
  \label{eq:ground-Witt-decomposition}
\end{equation}
The direct sum in \eqref{eq:ground-Witt-decomposition} is Hilbert orthogonal,
but it is not Krein orthogonal.

\subsection{The full neutral circle}

The null condition by itself does not select \(\NzeroPlus\).  Every normalized
neutral ray in the ground plane has the form
\begin{equation}
  \boxed{
  \ket{\Xi_0(\theta)}
  :=\frac1{\sqrt2}
  \bigl(\ket0\oplus e^{i\theta}\ket{\Omega_0}\bigr),
  \qquad
  \theta\in\RR/2\pi\ZZ,}
  \label{eq:neutral-ground-circle}
\end{equation}
up to an overall phase.

\begin{proposition}[Projectivized neutral locus]
\label{prop:projectivized-neutral-locus}
The projectivized neutral cone in \(\Kzero\) is a circle parameterized by
\(\theta\) in \eqref{eq:neutral-ground-circle}.  For every \(\theta\),
\begin{equation}
  \kInner{\Xi_0(\theta)}{\Xi_0(\theta)}=0,
  \qquad
  \|\PP\Xi_0(\theta)\|^2
  =\|\PF\Xi_0(\theta)\|^2=\frac12.
  \label{eq:neutral-circle-common-properties}
\end{equation}
The bridge-selected line is \(\theta=0\), while the companion reflection
line is \(\theta=\pi\).
\end{proposition}

\begin{proof}
A vector in \(\Kzero\) has the form
\(a\ket0\oplus b\ket{\Omega_0}\).  Hilbert normalization and Krein
neutrality give
\begin{equation}
  |a|^2+|b|^2=1,
  \qquad
  |a|^2-|b|^2=0,
\end{equation}
so \(|a|=|b|=2^{-1/2}\).  Removing a common phase yields
\eqref{eq:neutral-ground-circle}.
\end{proof}

This result is conceptually important:
\begin{equation}
  \boxed{
  \begin{aligned}
  &\text{equal sector norms define a neutral locus;}\\
  &\text{the Awareness bridge selects one line on that locus.}
  \end{aligned}}
  \label{eq:null-locus-versus-bridge-selection}
\end{equation}

\subsection{Relative-phase coordinates}

The structural operators \(\Rzero\) and \(\Qzero\) resolve the relative phase:
\begin{align}
  \bra{\Xi_0(\theta)}\Rzero\ket{\Xi_0(\theta)}
  &=\cos\theta,
  \label{eq:R0-phase-expectation}\\
  \bra{\Xi_0(\theta)}\Qzero\ket{\Xi_0(\theta)}
  &=\sin\theta,
  \label{eq:Q0-phase-expectation}\\
  \bra{\Xi_0(\theta)}\fundJ\ket{\Xi_0(\theta)}
  &=0.
  \label{eq:J-phase-expectation}
\end{align}
Thus \(\theta=0\) and \(\theta=\pi\) are the two reflection eigenlines,
while \(\theta=\pm\pi/2\) are the two \(\Qzero\)-eigenlines.

\warninglabel Equations \eqref{eq:R0-phase-expectation}--
\eqref{eq:J-phase-expectation} describe the internal geometry of the ground
plane.  They do not imply a measurable physiological phase offset.  Whether
cross-sector phase can be operationally accessed depends on the observable
algebra.

\subsection{Status of the companion branch}

\begin{identification}[Anti-aligned companion, not a lower state]
\label{id:companion-branch-status}
The vector \(\ket{\Xi_0^-}\) is interpreted as the canonical
anti-aligned relative-phase companion of the nondual ground.  It is not:
\begin{enumerate}[label=(\roman*)]
  \item a state below the ground;
  \item a state of non-luminosity;
  \item the Fourth Yoga;
  \item or a ket \(\ket{\Omega_{-1}}\).
\end{enumerate}
It has the same projective sectoral compressions as \(\ket{\Xi_0^+}\); only
the relative phase differs.
\end{identification}

Indeed,
\begin{equation}
  [\PP\Xi_0^-]=[\ket0],
  \qquad
  [\PF\Xi_0^-]=[\ket{\Omega_0}],
  \label{eq:companion-same-sector-rays}
\end{equation}
so purely sectoral data cannot distinguish the two branches.

The bridge itself selects the aligned line kinematically; no additional
alignment Hamiltonian is assumed.
% =============================================================================
% END INLINED FILE: sections/13_ground_reflection_phase.tex
% =============================================================================
% =============================================================================
% BEGIN INLINED FILE: sections/14_ground_in_full_polar_geometry.tex
% =============================================================================
\section{Embedding the ground line into the full polar geometry}
\label{sec:ground-full-polar}

The restricted ground channel has a canonical polar orientation under the
weak bridge.  This section identifies exactly what the stronger compatibility
condition adds.

\subsection{The full active-support polar reflection}

Let
\begin{equation}
  \suppProjP:=\polA^*\polA=\pi_{A,+},
  \qquad
  \suppProjF:=\polA\polA^*=\pi_{A,-},
  \qquad
  \PiA:=\begin{pmatrix}\suppProjP&0\\0&\suppProjF\end{pmatrix}.
  \label{eq:full-polar-supports}
\end{equation}
Define the full supportwise polar operators
\begin{equation}
  \RA:=\begin{pmatrix}0&\polA^*\\\polA&0\end{pmatrix},
  \qquad
  \IA:=\fundJ\RA.
  \label{eq:full-polar-RI}
\end{equation}
They satisfy
\begin{equation}
  \RA^2=\PiA,
  \qquad
  \IA^2=-\PiA,
  \label{eq:full-polar-squares}
\end{equation}
as imported from the polar layer of TAUM \cite{TAUM2026}.

The full aligned active graph is
\begin{equation}
  \mathscr N_A^+
  :=\{\psi\oplus\polA\psi:
       \psi\in\Ran\suppProjP\},
  \label{eq:full-active-aligned-graph-preview}
\end{equation}
and the anti-aligned graph is defined with \(-\polA\).  This graph description
uses the standard geometry of neutral subspaces in a Krein space
\cite{Bognar}.  The symbols $\pi_{A,+}$ and $\pi_{A,-}$ are retained here only as a
crosswalk to the TAUM support notation; the cognitive realization uses
$\suppProjP$ and $\suppProjF$ to avoid collision with the representation
symbols $\piP$ and $\piF$.  Their complete Born geometry is developed in
\cref{sec:active-support-brlc-geometry,sec:born-rule-locus-correspondence}.

\subsection{Strong compatibility as a compression theorem}

\begin{theorem}[Ground compression of the full polar geometry]
\label{thm:ground-compression-full-polar}
Assume any of the equivalent strong compatibility conditions in
\cref{thm:strong-ground-equivalences}.  Then
\begin{equation}
  \PzeroP\le\suppProjP,
  \qquad
  \PzeroF\le\suppProjF,
  \label{eq:ground-support-contained-active}
\end{equation}
and
\begin{align}
  \RA\Pizero
  &=\Pizero\RA=\Rzero,
  \label{eq:R0-compression-RA}\\
  \IA\Pizero
  &=\Pizero\IA=\Izero.
  \label{eq:I0-compression-IA}
\end{align}
Consequently,
\begin{equation}
  \NzeroPlus\subseteq\mathscr N_A^+,
  \qquad
  \NzeroMinus\subseteq\mathscr N_A^-.
  \label{eq:ground-graphs-subset-full-graphs}
\end{equation}
\end{theorem}

\begin{proof}
Strong compatibility gives
\(\polA\ket0=\ket{\Omega_0}\).  Since the image has norm one,
\(\ket0\) lies in the initial support of \(\polA\), and
\(\ket{\Omega_0}\) lies in its final support.  Hence
\eqref{eq:ground-support-contained-active}.  On the two ground rays,
\(\polA\PzeroP=\uzero\) and
\(\polA^*\PzeroF=\uzero^*\), which gives the compression formulas.  The graph
inclusions follow directly.
\end{proof}

\begin{corollary}[Strong compatibility is an embedding condition]
\label{cor:strong-is-embedding-condition}
The strong condition does not create the neutral ground line.  It states that
the already canonical line \(\NzeroPlus\) is the ground slice of the full
active-support neutral graph determined by \(\polA\).
\end{corollary}

This clarifies the logical hierarchy:
\begin{equation}
  \boxed{
  \begin{aligned}
  A\ket0\propto\ket{\Omega_0}
  &\quad\Longrightarrow\quad
  \text{canonical restricted ground graph},\\
  A^*A\ket0=a_0^2\ket0
  &\quad\Longrightarrow\quad
  \text{ground graph embeds into the full polar graph}.
  \end{aligned}}
  \label{eq:weak-strong-hierarchy}
\end{equation}

\subsection{The SSP hierarchy}

The structural-sector-projection reading now has a precise three-level form.

\begin{enumerate}[label=(\roman*)]
  \item The total ray \([\ket{\Xi_0^+}]\) is the relational nondual ground in
  the common carrier.
  \item The structural compressions \(\PP\) and \(\PF\) extract its physical
  and phenomenal aspects.
  \item Renormalizing those components gives the conventional sectoral ground
  states \([\ket0]\) and \([\ket{\Omega_0}]\).
\end{enumerate}

For \(\psi\in\GzeroP\),
\begin{equation}
  \WzeroPlus\psi
  =\frac1{\sqrt2}
   \bigl(\psi\oplus\uzero\psi\bigr).
  \label{eq:W0-SSP-summary}
\end{equation}
The total ground is not a third substance over and above matter and
luminosity.  It is their bridge relation represented as a state in the common
carrier.

\begin{remark}[What is and is not measurable]
\label{rem:ground-measurability}
At this structural level, \(\PP\) and \(\PF\) are not included in a physical
observable algebra.
If they were treated as mutually exclusive measurement outcomes, the null
ground would yield a literal \(1/2\)--\(1/2\) sector probability.  The SSP
reading rejects that interpretation.  The equal weights are structural.  The
operationally meaningful Born correspondence instead compares matched
physical and phenomenal effects through the active-support isometry in
\cref{sec:born-rule-locus-correspondence,sec:paired-observable-algebra}.
\end{remark}

\subsection{The ground line and doctrinal nonduality}

\begin{identification}[Emptiness--luminosity neutral relation]
\label{id:emptiness-luminosity-neutral-relation}
The bridge-selected null line
\begin{equation}
  \NzeroPlus
  =\{\psi\oplus\uzero\psi:\psi\in\GzeroP\}
  \label{eq:doctrinal-neutral-line}
\end{equation}
is interpreted as the nondual relation between the physical vacuum as
emptiness and the phenomenal ground as luminosity.
\end{identification}

This identification does not replace the separate claims that the vacuum is
cyclic for matter and the luminous ground is cyclic for appearance.  Rather,
it adds:
\begin{equation}
  \boxed{
  \begin{aligned}
  &\text{the two generative grounds are balanced sectoral aspects}\\
  &\text{of one bridge-selected relational ground.}
  \end{aligned}}
  \label{eq:two-generative-grounds-one-relation}
\end{equation}

The word ``null'' is therefore used in the Krein-geometric sense of zero net
grading imbalance.  It is not a synonym for nothingness, nonexistence, blank
awareness, or absence of luminosity.
% =============================================================================
% END INLINED FILE: sections/14_ground_in_full_polar_geometry.tex
% =============================================================================
% =============================================================================
% BEGIN INLINED FILE: sections/16_active_support_brlc_geometry.tex
% =============================================================================
\section{Active-support polar correspondence and the BRLC graph}
\label{sec:active-support-brlc-geometry}

The bridge-selected neutral line describes the relation on the two ground
rays.  The construction below extends to the full support on which the polar
orientation of the closed Awareness channel is norm preserving.  This is the
largest subspace on which a basis-free Born correspondence follows from the
channel alone.

\subsection{The active physical and phenomenal supports}

Let
\begin{equation}
  A=\polA\absA
  \label{eq:active-polar-decomposition}
\end{equation}
be the polar decomposition of the densely defined closed typed channel
\(A:\Dom A\subseteq\scrHP\to\scrHF\).  Its initial and final support
projections are
\begin{equation}
  \suppProjP:=\polA^*\polA,
  \qquad
  \suppProjF:=\polA\polA^*.
  \label{eq:active-support-projections}
\end{equation}
Define
\begin{equation}
  \boxed{
  \actHP:=\Ran\suppProjP=(\Ker A)^\perp,
  \qquad
  \actHF:=\Ran\suppProjF=\overline{\Ran A}.}
  \label{eq:active-support-spaces}
\end{equation}
The active carrier is
\begin{equation}
  \actK:=\actHP\oplus\actHF=\Ran\PiA,
  \qquad
  \PiA=
  \begin{pmatrix}
    \suppProjP&0\\
    0&\suppProjF
  \end{pmatrix}.
  \label{eq:active-carrier}
\end{equation}

\begin{theorem}[Unitary polar orientation on active support]
\label{thm:polar-unitary-active}
The restriction
\begin{equation}
  \restr{\polA}{\actHP}:\actHP\longrightarrow\actHF
  \label{eq:uA-active-unitary-map}
\end{equation}
is unitary.  In particular,
\begin{equation}
  \inner{\polA\psi}{\polA\chi}_{\Phi}
  =\inner{\psi}{\chi}_{\Psi}
  \qquad
  (\psi,\chi\in\actHP).
  \label{eq:uA-active-inner-product}
\end{equation}
Its inverse is \(\restr{\polA^*}{\actHF}\).
\end{theorem}

\begin{proof}
For \(\psi,\chi\in\actHP\),
\begin{equation}
  \inner{\polA\psi}{\polA\chi}_{\Phi}
  =\inner{\psi}{\polA^*\polA\chi}_{\Psi}
  =\inner{\psi}{\suppProjP\chi}_{\Psi}
  =\inner{\psi}{\chi}_{\Psi}.
\end{equation}
Thus the restriction is isometric.  Its range is
\(\Ran\suppProjF=\actHF\), and its adjoint restriction is its inverse.
\end{proof}

\statuslabel{Imported theorem} This is the supportwise unitary content of
polar decomposition \cite{Kato,ReedSimon,TAUM2026}.  No isometry of the raw
channel \(A\) is assumed.

\subsection{The aligned active graph}

\begin{definition}[Active graph embeddings]
\label{def:active-graph-embeddings}
For \(\theta\in\RR/2\pi\ZZ\), define
\begin{equation}
  \boxed{
  \WAtheta\psi
  :=\frac1{\sqrt2}
    \bigl(\psi\oplus e^{i\theta}\polA\psi\bigr),
  \qquad
  \psi\in\actHP.}
  \label{eq:WA-theta-definition}
\end{equation}
Write
\begin{equation}
  \WA:=W_A(0),
  \qquad
  W_A^-:=W_A(\pi).
  \label{eq:WA-plus-minus}
\end{equation}
The corresponding graph subspaces are
\begin{equation}
  \NAtheta:=\Ran\WAtheta
  =\left\{
  2^{-1/2}(\psi\oplus e^{i\theta}\polA\psi):
  \psi\in\actHP
  \right\}.
  \label{eq:NA-theta-definition}
\end{equation}
In particular,
\begin{equation}
  \NAplus=\mathscr N_A(0),
  \qquad
  \NAminus=\mathscr N_A(\pi).
  \label{eq:NA-plus-minus-def}
\end{equation}
\end{definition}

The factor \(2^{-1/2}\) normalizes the total Hilbert vector.  It does not
alter the underlying graph relation
\begin{equation}
  \phi=e^{i\theta}\polA\psi.
  \label{eq:active-graph-relation}
\end{equation}

\begin{theorem}[Isometric graph embedding]
\label{thm:WA-isometry}
For every \(\theta\), \(\WAtheta:\actHP\to\actK\) is an isometry with
closed range \(\NAtheta\).  More precisely,
\begin{equation}
  \inner{\WAtheta\psi}{\WAtheta\chi}_{\scrK}
  =\inner{\psi}{\chi}_{\Psi}
  \qquad
  (\psi,\chi\in\actHP).
  \label{eq:WA-inner-preservation}
\end{equation}
Its adjoint is
\begin{equation}
  \WAtheta^*(x\oplus y)
  =\frac1{\sqrt2}
   \bigl(x+e^{-i\theta}\polA^*y\bigr),
  \qquad x\in\actHP,\ y\in\actHF.
  \label{eq:WA-adjoint}
\end{equation}
\end{theorem}

\begin{proof}
Using \cref{thm:polar-unitary-active},
\begin{align}
  \inner{\WAtheta\psi}{\WAtheta\chi}_{\scrK}
  &=\frac12\left(
      \inner{\psi}{\chi}_{\Psi}
      +\inner{e^{i\theta}\polA\psi}
              {e^{i\theta}\polA\chi}_{\Phi}
    \right)\\
  &=\inner{\psi}{\chi}_{\Psi}.
\end{align}
An isometry has closed range.  The adjoint formula follows by direct
calculation.
\end{proof}

\subsection{Total and hypermaximal neutrality}

\begin{theorem}[Hypermaximal neutral active graphs]
\label{thm:active-graphs-hypermaximal-neutral}
For every \(\theta\),
\begin{equation}
  \kInner{\xi}{\eta}=0
  \qquad
  \text{for all }\xi,\eta\in\NAtheta.
  \label{eq:active-total-neutrality}
\end{equation}
Moreover, inside the nondegenerate active Krein space \(\actK\),
\begin{equation}
  \NAtheta^{[\perp]}=\NAtheta.
  \label{eq:active-hypermax-neutrality}
\end{equation}
Thus each \(\NAtheta\) is hypermaximal neutral.
\end{theorem}

\begin{proof}
Let
\begin{equation}
  \xi=\WAtheta\psi,
  \qquad
  \eta=\WAtheta\chi.
\end{equation}
Then
\begin{align}
  \kInner{\xi}{\eta}
  &=\frac12\left(
      \inner{\psi}{\chi}_{\Psi}
      -\inner{e^{i\theta}\polA\psi}
              {e^{i\theta}\polA\chi}_{\Phi}
    \right)=0.
\end{align}
Now let \(x\oplus y\in\actK\) be Krein-orthogonal to every vector in
\(\NAtheta\).  For all \(\psi\in\actHP\),
\begin{equation}
  0
  =\inner{x}{\psi}_{\Psi}
   -\inner{y}{e^{i\theta}\polA\psi}_{\Phi}
  =\inner{x-e^{-i\theta}\polA^*y}{\psi}_{\Psi}.
\end{equation}
Hence \(x=e^{-i\theta}\polA^*y\), or equivalently
\(y=e^{i\theta}\polA x\).  Therefore \(x\oplus y\in\NAtheta\).
\end{proof}

This is the standard neutral-graph mechanism for an isometry between the
positive and negative parts of a Krein space; see \cite{Bognar}.

The bare equal-norm cone
\begin{equation}
  \{\psi\oplus\phi:\|\psi\|=\|\phi\|\}
  \label{eq:bare-active-neutral-cone}
\end{equation}
is not linear.  Each \(\NAtheta\) is a closed linear slice of that cone,
selected by a specific isometric relation.

\subsection{Graph projections and the polar Pauli triple}

Define the Hilbert-orthogonal projection onto \(\NAtheta\) by
\begin{equation}
  \PAtheta:=\WAtheta\WAtheta^*.
  \label{eq:graph-projection-definition}
\end{equation}

\begin{proposition}[Projection onto a phase graph]
\label{prop:phase-graph-projection}
On \(\actK\),
\begin{equation}
  \boxed{
  \PAtheta
  =\frac12
   \begin{pmatrix}
     \idop_{\actHP}&e^{-i\theta}\polA^*\\
     e^{i\theta}\polA&\idop_{\actHF}
   \end{pmatrix}.}
  \label{eq:phase-graph-projection-block}
\end{equation}
As an operator on the full carrier, the diagonal identities in
\eqref{eq:phase-graph-projection-block} are replaced by
\(\suppProjP\) and \(\suppProjF\).  Equivalently,
\begin{equation}
  \boxed{
  \PAtheta
  =\frac12\left(
      \PiA+\cos\theta\,\RA+\sin\theta\,(-i\IA)
    \right).}
  \label{eq:phase-graph-projection-Pauli}
\end{equation}
In particular,
\begin{equation}
  \PAPlus=\frac12(\PiA+\RA),
  \qquad
  \PAMinus=\frac12(\PiA-\RA).
  \label{eq:aligned-anti-projections}
\end{equation}
\end{proposition}

\begin{proof}
Insert \eqref{eq:WA-adjoint} into
\(\PAtheta=\WAtheta\WAtheta^*\).  The second formula follows from
\begin{equation}
  \RA=
  \begin{pmatrix}0&\polA^*\\\polA&0\end{pmatrix},
  \qquad
  -i\IA=
  \begin{pmatrix}0&-i\polA^*\\i\polA&0\end{pmatrix}.
\end{equation}
\end{proof}

\begin{corollary}[Active Hilbert decomposition]
\label{cor:active-graph-decomposition}
The aligned and anti-aligned graph spaces are Hilbert orthogonal and
\begin{equation}
  \boxed{
  \actK=\NAplus\oplus\NAminus.}
  \label{eq:active-Hilbert-Witt-decomposition}
\end{equation}
They are the \(\pm1\) eigenspaces of \(\RA\) on \(\actK\).
\end{corollary}

\begin{proof}
The projection formulas in \eqref{eq:aligned-anti-projections} are
complementary on \(\actK\).  The eigenspace statement follows from
\(\RA\WA=\WA\) and \(\RA W_A^-=-W_A^-\).
\end{proof}

\subsection{The ground line inside the active graph}

The full active graph exists without strong ground compatibility.  Its
relation to the ground graph remains conditional.

\begin{corollary}[Ground-inclusive BRLC]
\label{cor:ground-inclusive-BRLC}
If the equivalent strong compatibility conditions in
\cref{thm:strong-ground-equivalences} hold, then
\begin{equation}
  \WA\ket0=\ket{\Xi_0^+},
  \qquad
  \NzeroPlus\subseteq\NAplus.
  \label{eq:ground-in-global-BRLC}
\end{equation}
Without strong compatibility, \(\NzeroPlus\) remains canonical by the weak
bridge, but it need not be the ground slice of \(\NAplus\).
\end{corollary}

\statuslabel{Conditional embedding} The active-support BRLC requires no
strong ground axiom.  Strong compatibility is needed only if the physical
vacuum and luminous ground are required to lie on that same global polar
correspondence.
% =============================================================================
% END INLINED FILE: sections/16_active_support_brlc_geometry.tex
% =============================================================================
% =============================================================================
% BEGIN INLINED FILE: sections/17_born_rule_locus_correspondence.tex
% =============================================================================
\section{The Born Rule Locus Correspondence}
\label{sec:born-rule-locus-correspondence}

The Born Rule Locus Correspondence expresses a basis-free requirement:
corresponding physical and phenomenal states carry the same probability
geometry.  The probability-preserving map is the polar orientation
\(\polA\), and the locus is its projectivized aligned neutral graph.

\subsection{Projective transition probabilities}

For nonzero vectors \(\xi,\eta\) in a Hilbert space, define the pure-state
transition probability
\begin{equation}
  \BRprob([\xi],[\eta])
  :=\frac{|\inner{\xi}{\eta}|^2}
          {\|\xi\|^2\|\eta\|^2}.
  \label{eq:projective-transition-probability}
\end{equation}
This depends only on the rays \([\xi]\) and \([\eta]\).

\begin{definition}[Born Rule Locus]
\label{def:born-rule-locus}
The active Born Rule Locus of the channel \(A\) is
\begin{equation}
  \boxed{
  \LBR:=\projective{\NAplus}
  =\left\{
      [\WA\psi]:0\neq\psi\in\actHP
    \right\}.}
  \label{eq:BR-locus-definition}
\end{equation}
The associated projective correspondence maps are
\begin{align}
  \beta_A:
  \projective{\actHP}&\longrightarrow\projective{\actHF},
  &\beta_A([\psi])&=[\polA\psi],
  \label{eq:betaA-projective-map}\\
  \gamma_A:
  \projective{\actHP}&\longrightarrow\LBR,
  &\gamma_A([\psi])&=[\WA\psi].
  \label{eq:gammaA-projective-map}
\end{align}
\end{definition}

\begin{theorem}[Born Rule Locus Correspondence]
\label{thm:BRLC}
The maps \(\beta_A\) and \(\gamma_A\) are bijections onto their displayed
ranges and preserve all pure-state transition probabilities.  Explicitly,
for nonzero \(\psi,\chi\in\actHP\),
\begin{equation}
  \boxed{
  \BRprob([\psi],[\chi])
  =\BRprob([\polA\psi],[\polA\chi])
  =\BRprob([\WA\psi],[\WA\chi]).}
  \label{eq:BRLC-transition-probability}
\end{equation}
Equivalently, for normalized \(\psi,\chi\),
\begin{equation}
  |\inner{\psi}{\chi}|^2
  =|\inner{\polA\psi}{\polA\chi}|^2
  =|\inner{\WA\psi}{\WA\chi}|^2.
  \label{eq:BRLC-normalized-form}
\end{equation}
\end{theorem}

\begin{proof}
The restriction of \(\polA\) to \(\actHP\) is unitary by
\cref{thm:polar-unitary-active}, and \(\WA\) is an isometry by
\cref{thm:WA-isometry}.  Both therefore preserve inner products and norms.
They are injective, and their ranges are respectively \(\actHF\) and
\(\NAplus\).  Projectivization gives the stated bijections and
\eqref{eq:BRLC-transition-probability}.
\end{proof}

\statuslabel{Derived theorem} The BRLC is not inferred from the existence of
two Hilbert spaces alone.  It follows from the polar partial isometry of the
closed channel after restriction to its active support.

\subsection{Fubini--Study geometry}

For normalized representatives, use the convention
\begin{equation}
  d_{\mathrm{FS}}([\xi],[\eta])
  :=\arccos|\inner{\xi}{\eta}|.
  \label{eq:FS-distance-convention}
\end{equation}

\begin{corollary}[Projective isometry]
\label{cor:BRLC-FS-isometry}
The maps \(\beta_A\) and \(\gamma_A\) preserve the Fubini--Study distance:
\begin{equation}
  d_{\mathrm{FS}}([\psi],[\chi])
  =d_{\mathrm{FS}}([\polA\psi],[\polA\chi])
  =d_{\mathrm{FS}}([\WA\psi],[\WA\chi]).
  \label{eq:BRLC-FS-preservation}
\end{equation}
Thus \(\LBR\) is a projective Hilbert-space copy of both active sectoral
state spaces.
\end{corollary}

The correspondence preserves not only norms but the entire pure-state Born
geometry.

\subsection{Rank-one propositions}

For a normalized \(\psi\in\actHP\), write
\begin{equation}
  P_\psi:=\ket\psi\!\bra\psi,
  \qquad
  P_{\polA\psi}:=\ket{\polA\psi}\!\bra{\polA\psi}.
  \label{eq:rank-one-corresponding-projectors}
\end{equation}

\begin{proposition}[Correspondence of pure-state projectors]
\label{prop:pure-projector-correspondence}
For normalized \(\psi\in\actHP\),
\begin{equation}
  \boxed{
  P_{\polA\psi}=\polA P_\psi\polA^*.}
  \label{eq:pure-projector-conjugation}
\end{equation}
For normalized \(\psi,\chi\in\actHP\),
\begin{equation}
  \Tr(P_\psi P_\chi)
  =\Tr(P_{\polA\psi}P_{\polA\chi})
  =\BRprob([\psi],[\chi]).
  \label{eq:projector-transition-trace}
\end{equation}
\end{proposition}

\begin{proof}
The first identity follows by evaluating both sides on \(\actHF\).  The trace
identity is the standard rank-one formula and is preserved by unitary
conjugation.
\end{proof}

\subsection{Effects and probability distributions}

Let \(\mathcal B(\actHP)\) and \(\mathcal B(\actHF)\) denote the bounded
operators on the active supports.  Define
\begin{equation}
  \boxed{
  \alphaA(X):=\polA X\polA^*,
  \qquad
  X\in\mathcal B(\actHP).}
  \label{eq:alphaA-definition}
\end{equation}
The identities in this section are written on the active spaces.  When viewed
on the full sector Hilbert spaces, all operators are extended by zero on the
corresponding inactive complements.

\begin{theorem}[Active observable conjugacy]
\label{thm:alphaA-star-isomorphism}
The map
\begin{equation}
  \alphaA:\mathcal B(\actHP)\longrightarrow\mathcal B(\actHF)
  \label{eq:alphaA-map}
\end{equation}
is a unital normal \(*\)-isomorphism with inverse
\begin{equation}
  \alphaA^{-1}(Y)=\polA^*Y\polA.
  \label{eq:alphaA-inverse}
\end{equation}
It preserves positivity, effects, projections, spectra, and bounded
functional calculus.
\end{theorem}

\begin{proof}
On the active supports \(\polA\) is unitary.  The result is therefore unitary
conjugation.
\end{proof}

If \(E\) is a physical effect,
\begin{equation}
  0\le E\le\idop_{\actHP},
\end{equation}
then \(\alphaA(E)\) is the corresponding phenomenal effect.

\begin{theorem}[Effect form of the BRLC]
\label{thm:BRLC-effect-form}
For every nonzero \(\psi\in\actHP\) and every physical effect
\(E\in\mathcal B(\actHP)\),
\begin{equation}
  \boxed{
  \frac{\inner{\psi}{E\psi}_{\Psi}}{\|\psi\|^2}
  =
  \frac{\inner{\polA\psi}{\alphaA(E)\polA\psi}_{\Phi}}
       {\|\polA\psi\|^2}.}
  \label{eq:BRLC-effect-probability}
\end{equation}
Thus every bounded active physical yes/no proposition has a canonical
phenomenal correspondent with exactly the same Born probability.
\end{theorem}

\begin{proof}
Using \(\polA^*\polA=\idop_{\actHP}\),
\begin{align}
  \inner{\polA\psi}{\alphaA(E)\polA\psi}_{\Phi}
  &=\inner{\polA\psi}{\polA E\psi}_{\Phi}\\
  &=\inner{\psi}{E\psi}_{\Psi}.
\end{align}
The norms agree by \cref{thm:polar-unitary-active}.
\end{proof}

\begin{corollary}[POVM correspondence]
\label{cor:BRLC-POVM}
Let \((\mathsf X,\Sigma)\) be a measurable outcome space and let
\begin{equation}
  \mathsf E_{\Psi}:\Sigma\to\mathcal B(\actHP)
\end{equation}
be a normalized positive-operator-valued measure.  Then
\begin{equation}
  \mathsf E_{\Phi}(\Delta)
  :=\alphaA\bigl(\mathsf E_{\Psi}(\Delta)\bigr)
  \label{eq:phenomenal-POVM-definition}
\end{equation}
is a normalized POVM on \(\actHF\), and
\begin{equation}
  \Pr_{\psi}^{\Psi}(\Delta)
  =\Pr_{\polA\psi}^{\Phi}(\Delta)
  \qquad
  (\Delta\in\Sigma).
  \label{eq:POVM-distribution-equality}
\end{equation}
Projection-valued measures correspond to projection-valued measures.
\end{corollary}

\begin{proof}
Normality and positivity of \(\alphaA\) preserve countable additivity and
normalization.  The probability identity is
\cref{thm:BRLC-effect-form} applied to each measurable effect.
\end{proof}

The BRLC therefore does not depend on a preferred basis.  It identifies full
probability distributions for matched measurement structures.

\subsection{The projective locus}

The term \emph{locus} denotes the closed projective subset
\begin{equation}
  \LBR=\projective{\NAplus}
  \label{eq:locus-not-spacetime}
\end{equation}
of the total projective state space on which the physical and phenomenal
components satisfy the canonical polar relation.  A point of this locus is a
ray
\begin{equation}
  \left[
  \frac1{\sqrt2}
  \left(\psi\oplus\polA\psi\right)
  \right].
  \label{eq:point-on-BR-locus}
\end{equation}
The physical ray, phenomenal ray, and relational ray are three exactly
isometric presentations of the same active Born geometry.

\begin{identification}[Born Rule Locus Correspondence]
\label{id:BRLC-doctrinal}
The projective isometry
\begin{equation}
  [\psi]
  \longleftrightarrow
  [\polA\psi]
  \longleftrightarrow
  [\WA\psi]
  \label{eq:BRLC-doctrinal-triangle}
\end{equation}
is interpreted as the probability-preserving correspondence among a
physically realized state, its phenomenal presentation, and their nondual
relational representative.
\end{identification}

The supportwise theorem applies to a general closed channel on
\(\actHP=(\Ker A)^\perp\).  Under the cognitive realization axiom
\cref{ax:physical-presentability}, however, \(\actHP=\scrHP\): no
nonzero direction in the chosen physical carrier is excluded by the
Awareness channel.  What remains open is the identification of that
carrier with physical degrees of freedom.  Taking it to represent a
cognitive subsystem does not identify it with all matter; taking it to
represent all matter makes the full-presentability assumption correspondingly
stronger.  A proper support inside the same chosen carrier would require
relaxing \cref{ax:physical-presentability}.
% =============================================================================
% END INLINED FILE: sections/17_born_rule_locus_correspondence.tex
% =============================================================================
% =============================================================================
% BEGIN INLINED FILE: sections/18_paired_observable_algebra.tex
% =============================================================================
\section{The paired observable algebra}
\label{sec:paired-observable-algebra}

The BRLC identifies matched effects by unitary conjugation.  The SSP reading
requires a corresponding algebra on the total active carrier: an observable
should ask the same question in the physical and phenomenal descriptions,
not ask which sector is ``really'' present.

\subsection{Selecting the physical correspondence algebra}

The algebraic formula \(X\mapsto\alphaA(X)\) is available for every bounded
operator on \(\actHP\).  Physical admissibility is a separate selection.

\begin{modelchoice}[Active physical observable algebra]
\label{choice:active-physical-observable-algebra}
Choose a unital \(C^*\)-subalgebra
\begin{equation}
  \boxed{
  \obsP\subseteq\mathcal B(\actHP)}
  \label{eq:active-physical-observable-algebra}
\end{equation}
whose elements are the physical observables or effects to which the
cross-domain correspondence is to be applied.
\end{modelchoice}

Define the corresponding phenomenal algebra by
\begin{equation}
  \boxed{
  \obsF:=\alphaA(\obsP)
  \subseteq\mathcal B(\actHF).}
  \label{eq:corresponding-phenomenal-algebra}
\end{equation}
The choice \(\obsP=\mathcal B(\actHP)\) is algebraically maximal.  A concrete
physical model may select a smaller local, gauge-invariant, or otherwise
admissible algebra.

The algebra \(\obsF\) is the \emph{corresponding} phenomenal algebra.  It is
not assumed to exhaust the full algebra of phenomenal determination selected
in \cref{sec:phenomenal-determination-gns,sec:vasana-generator-system}.

\subsection{The maximal and selected paired representations}

\begin{definition}[Maximal and selected paired embeddings]
\label{def:paired-observable-embedding}
For every $X\in\mathcal B(\actHP)$, define the maximal paired embedding
\begin{equation}
  \boxed{
  \iotaMaxA(X)
  :=X\oplus\alphaA(X)
  =\begin{pmatrix}
      X&0\\
      0&\polA X\polA^*
    \end{pmatrix}}
  \label{eq:iotaA-maximal-definition}
\end{equation}
as an operator on $\actK$.  The selected paired embedding is the restriction
\begin{equation}
  \boxed{
  \iotaA:=\restr{\iotaMaxA}{\obsP}.}
  \label{eq:iotaA-definition}
\end{equation}
The selected paired observable algebra is
\begin{equation}
  \boxed{
  \obsPair:=\iotaA(\obsP).}
  \label{eq:paired-observable-algebra-definition}
\end{equation}
Its unit is the active support identity $\PiA$.
\end{definition}

The maximal map is retained because bounded Borel functional calculus may
produce spectral projections in $\mathcal B(\actHP)$ that do not lie in the
selected $C^*$-algebra $\obsP$.  Physical admissibility remains encoded by the
restriction $\iotaA$; the extension supplies a mathematically unambiguous
operator domain when such projections are used.

\begin{theorem}[Paired algebra representations]
\label{thm:paired-algebra-representation}
The map
\begin{equation}
  \iotaMaxA:\mathcal B(\actHP)\longrightarrow\mathcal B(\actK)
  \label{eq:iotaA-maximal-map}
\end{equation}
is an injective unital normal $*$-homomorphism.  Its restriction
\begin{equation}
  \iotaA:\obsP\longrightarrow\obsPair
  \label{eq:iotaA-map}
\end{equation}
is an injective unital $*$-homomorphism.  Hence $\obsPair$ is canonically
$*$-isomorphic to both $\obsP$ and $\obsF$.
\end{theorem}

\begin{proof}
For bounded $X,Y$ on $\actHP$,
\begin{align}
  \iotaMaxA(XY)
  &=XY\oplus\alphaA(XY)\\
  &=(X\oplus\alphaA(X))(Y\oplus\alphaA(Y)),\\
  \iotaMaxA(X^*)&=\iotaMaxA(X)^*.
\end{align}
The unit is
\begin{equation}
  \idop_{\actHP}\oplus\idop_{\actHF}=\PiA
\end{equation}
when the operators are viewed on the full carrier.  Normality follows from
normality of unitary conjugation on the active supports, and injectivity
follows from physical compression.  The selected statement is the
restriction to $\obsP$.
\end{proof}

\subsection{The graph intertwining theorem}

\begin{theorem}[Paired operators preserve every phase graph]
\label{thm:paired-algebra-graph-intertwining}
For every $X\in\mathcal B(\actHP)$, every $\theta$, and every
$\psi\in\actHP$,
\begin{equation}
  \boxed{
  \iotaMaxA(X)\WAtheta\psi
  =\WAtheta X\psi.}
  \label{eq:paired-graph-intertwining}
\end{equation}
Equivalently,
\begin{equation}
  \WAtheta^*\iotaMaxA(X)\WAtheta=X.
  \label{eq:paired-compression-identity}
\end{equation}
For $X\in\obsP$, the same identities hold with $\iotaA(X)$.  Thus every
$\NAtheta$ is invariant under the maximal paired representation and, in
particular, under $\obsPair$.
\end{theorem}

\begin{proof}
Using $\polA^*\polA=\idop_{\actHP}$,
\begin{align}
  \iotaMaxA(X)\WAtheta\psi
  &=\frac1{\sqrt2}
    \left(
      X\psi\oplus
      e^{i\theta}\polA X\polA^*\polA\psi
    \right)\\
  &=\frac1{\sqrt2}
    \left(
      X\psi\oplus e^{i\theta}\polA X\psi
    \right)\\
  &=\WAtheta X\psi.
\end{align}
Left multiplication by $\WAtheta^*$ gives the compression identity.
\end{proof}

\begin{corollary}[Expectation and matrix-element equality]
\label{cor:paired-expectation-equality}
For $X\in\mathcal B(\actHP)$ and $\psi,\chi\in\actHP$,
\begin{equation}
  \boxed{
  \inner{\WAtheta\psi}{\iotaMaxA(X)\WAtheta\chi}_{\scrK}
  =\inner{\psi}{X\chi}_{\Psi}.}
  \label{eq:paired-matrix-element-equality}
\end{equation}
For normalized $\psi$,
\begin{equation}
  \inner{\WAtheta\psi}{\iotaMaxA(X)\WAtheta\psi}_{\scrK}
  =\inner{\psi}{X\psi}_{\Psi}
  =\inner{\polA\psi}{\alphaA(X)\polA\psi}_{\Phi}.
  \label{eq:three-expectation-equality}
\end{equation}
When $X\in\obsP$, $\iotaMaxA(X)=\iotaA(X)$.
\end{corollary}

This is the SSP replacement for a literal matter-or-mind measurement.  The
same proposition has one expectation value expressed through three
representations.

\subsection{Paired effects and the absence of a one-half loss}

If \(E\in\obsP\) is an effect, then
\begin{equation}
  \iotaA(E)=E\oplus\alphaA(E)
  \label{eq:paired-effect}
\end{equation}
is an effect on \(\actK\).

\begin{proposition}[Paired probability is not halved]
\label{prop:paired-probability-not-halved}
For normalized \(\psi\in\actHP\),
\begin{align}
  \inner{\WA\psi}{\iotaA(E)\WA\psi}_{\scrK}
  &=\frac12\inner{\psi}{E\psi}_{\Psi}
    +\frac12\inner{\polA\psi}{\alphaA(E)\polA\psi}_{\Phi}
  \label{eq:paired-two-halves}\\
  &=\inner{\psi}{E\psi}_{\Psi}.
  \label{eq:paired-full-probability}
\end{align}
The two equal structural contributions sum to the ordinary Born probability.
\end{proposition}

\begin{proof}
The first equality is block expansion.  The second is
\cref{thm:BRLC-effect-form}.
\end{proof}

The coefficients \(1/2\) arise from Hilbert normalization of the total graph
state.  They do not imply that the physical and phenomenal descriptions are
competing outcomes.  The paired effect contains both corresponding blocks,
so the total probability is the same probability as in either conventional
sectoral description.

\subsection{Spectral distributions}

\begin{corollary}[Matched spectral measures]
\label{cor:paired-spectral-measures}
Let \(X=X^*\in\obsP\), and let \(E_X\) be its spectral measure.  Then the
spectral measure of \(\alphaA(X)\) is
\begin{equation}
  E_{\alphaA(X)}(\Delta)=\alphaA(E_X(\Delta)),
  \label{eq:corresponding-spectral-measure}
\end{equation}
and the spectral distribution in the three corresponding states agrees:
\begin{equation}
  \inner{\psi}{E_X(\Delta)\psi}_{\Psi}
  =\inner{\polA\psi}
         {E_{\alphaA(X)}(\Delta)\polA\psi}_{\Phi}
  =\inner{\WA\psi}
         {\iotaMaxA(E_X(\Delta))\WA\psi}_{\scrK}.
  \label{eq:three-spectral-distributions}
\end{equation}
\end{corollary}

\begin{proof}
Unitary conjugation preserves the bounded Borel functional calculus.  Each
spectral projection $E_X(\Delta)$ belongs to $\mathcal B(\actHP)$, so the
maximal embedding $\iotaMaxA$ is defined even when that projection does not
belong to the selected $C^*$-algebra $\obsP$.  Apply
\cref{cor:paired-expectation-equality}.
\end{proof}

Thus the BRLC identifies complete outcome distributions, not only pairwise
state overlaps.

\subsection{Normal states and density operators}

The pure-state geometry extends immediately to normal mixed states on the
active supports.

\begin{proposition}[Mixed-state correspondence]
\label{prop:mixed-state-correspondence}
Let \(\rho\ge0\) be trace class on \(\actHP\) with \(\Tr\rho=1\).  Define
\begin{equation}
  \rho_{\Phi}:=\polA\rho\polA^*,
  \qquad
  \rho_{\mathrm{rel}}^{(\theta)}
  :=\WAtheta\rho\WAtheta^*.
  \label{eq:mixed-state-maps}
\end{equation}
Then both are positive trace-one operators, and for every \(X\in\obsP\),
\begin{equation}
  \boxed{
  \Tr(\rho X)
  =\Tr(\rho_{\Phi}\alphaA(X))
  =\Tr(\rho_{\mathrm{rel}}^{(\theta)}\iotaA(X)).}
  \label{eq:mixed-state-expectation-equality}
\end{equation}
Moreover,
\begin{equation}
  \rho_{\mathrm{rel}}^{(\theta)}
  =\frac12
   \begin{pmatrix}
     \rho&e^{-i\theta}\rho\polA^*\\
     e^{i\theta}\polA\rho&\polA\rho\polA^*
   \end{pmatrix}
  \label{eq:relational-density-block}
\end{equation}
on \(\actK\).
\end{proposition}

\begin{proof}
Unitary and isometric conjugation preserve positivity and trace.  The
expectation identities follow from cyclicity of the trace and
\eqref{eq:paired-compression-identity}.  The block formula follows from the
definition of \(\WAtheta\).
\end{proof}

The structural sector compressions are
\begin{equation}
  \PP\rho_{\mathrm{rel}}^{(\theta)}\PP
  =\frac12\rho,
  \qquad
  \PF\rho_{\mathrm{rel}}^{(\theta)}\PF
  =\frac12\rho_{\Phi}.
  \label{eq:mixed-sector-compressions}
\end{equation}
After renormalization, they recover the corresponding conventional density
operators.  The off-diagonal blocks encode the cross-sector relation and its
relative phase.

\begin{remark}[Direct sum, not partial trace]
\label{rem:direct-sum-not-partial-trace}
The carrier is presently a direct sum, not a bipartite tensor product of two
independent systems.  The operations in \eqref{eq:mixed-sector-compressions}
are sector compressions, not partial traces over independent subsystems.
\end{remark}

\subsection{The correspondence algebra is not yet the full phenomenal algebra}

The paired algebra answers only one question:
\begin{equation}
  \text{how is a selected physical proposition represented phenomenally and
  relationally?}
  \label{eq:paired-algebra-question}
\end{equation}
It does not yet encode the endogenous v\=asan\=a-conditioned operators that
generate differentiated \=ak\=aras, memory, conceptuality, affective
orientation, or grasper--grasped distortion.

\begin{modelchoice}[Correspondence versus generation]
\label{choice:correspondence-versus-generation}
The algebra \(\obsPair\) is treated as the cross-domain correspondence
algebra.  It is not identified with the full matter-generating algebra
\(\algP\), the full phenomenal determination algebra \(\algF\), or a final
total observable algebra.
\end{modelchoice}

This separation prevents the BRLC from overclaiming that every phenomenal
operation is merely a conjugate copy of a physical one.  The active
correspondence is exact where defined; phenomenal differentiation and excess
may contain additional structure.
% =============================================================================
% END INLINED FILE: sections/18_paired_observable_algebra.tex
% =============================================================================
% =============================================================================
% BEGIN INLINED FILE: sections/19_ssp_factorization_phase.tex
% =============================================================================
\section{SSP factorization, phase blindness, and structural observables}
\label{sec:ssp-factorization-phase}

The active supports are unitarily equivalent through \(\polA\).  This permits
a canonical factorization of the active carrier into a two-dimensional sector
coordinate and one copy of the active physical state space.  In that
factorization the SSP reading becomes completely explicit: paired observables
act on the state coordinate, while the Krein and relative-phase operators act
on the sector coordinate.

\subsection{Canonical sector tensorization}

Let \(\eplus,\eminus\) be the standard orthonormal basis of \(\CC^2\).
Define
\begin{equation}
  \boxed{
  \UA:\actK\longrightarrow\CC^2\otimes\actHP}
  \label{eq:UA-factorization-map}
\end{equation}
by
\begin{equation}
  \UA(\psi\oplus\phi)
  :=\eplus\otimes\psi
    +\eminus\otimes\polA^*\phi.
  \label{eq:UA-factorization-formula}
\end{equation}

\begin{theorem}[Active carrier factorization]
\label{thm:active-carrier-factorization}
The map \(\UA\) is unitary.  Under this unitary,
\begin{align}
  \UA\fundJ\UA^*
  &=\sigmaz\otimes\idop_{\actHP},
  \label{eq:J-factorized}\\
  \UA\RA\UA^*
  &=\sigmax\otimes\idop_{\actHP},
  \label{eq:R-factorized}\\
  \UA(-i\IA)\UA^*
  &=\sigmay\otimes\idop_{\actHP}.
  \label{eq:Q-factorized}
\end{align}
For every \(X\in\obsP\),
\begin{equation}
  \boxed{
  \UA\iotaA(X)\UA^*
  =\idop_{2}\otimes X.}
  \label{eq:paired-algebra-factorized}
\end{equation}
\end{theorem}

\begin{proof}
The norm identity
\begin{equation}
  \|\psi\oplus\phi\|_{\actK}^2
  =\|\psi\|^2+\|\polA^*\phi\|^2
\end{equation}
follows because \(\polA^*\) is unitary from \(\actHF\) to \(\actHP\).
Surjectivity is immediate.  The operator identities follow by evaluating the
block operators on \(\psi\oplus\phi\).  For the paired algebra,
\begin{align}
  \UA\iotaA(X)(\psi\oplus\phi)
  &=\eplus\otimes X\psi
    +\eminus\otimes\polA^*\polA X\polA^*\phi\\
  &=\eplus\otimes X\psi
    +\eminus\otimes X\polA^*\phi.
\end{align}
\end{proof}

The factor \(\CC^2\) is a structural sector coordinate.  It is not introduced
as a new material qubit or an independently measurable two-level system.

\subsection{The BRLC as a projective slice}

Define the equatorial sector vector
\begin{equation}
  \ket{\theta}_{\mathrm{sec}}
  :=\frac1{\sqrt2}
    \left(\eplus+e^{i\theta}\eminus\right).
  \label{eq:sector-equator-vector}
\end{equation}

\begin{proposition}[Factorization of phase-graph states]
\label{prop:phase-graph-factorization}
For every \(\psi\in\actHP\),
\begin{equation}
  \boxed{
  \UA\WAtheta\psi
  =\ket{\theta}_{\mathrm{sec}}\otimes\psi.}
  \label{eq:WA-factorized-state}
\end{equation}
Consequently,
\begin{equation}
  \UA\NAtheta
  =\ket{\theta}_{\mathrm{sec}}\otimes\actHP,
  \label{eq:NA-factorized-subspace}
\end{equation}
and
\begin{equation}
  \UA\PAtheta\UA^*
  =\proj{\theta}_{\mathrm{sec}}\otimes\idop_{\actHP}.
  \label{eq:graph-projection-factorized}
\end{equation}
\end{proposition}

\begin{proof}
Substitute \eqref{eq:WA-theta-definition} into
\eqref{eq:UA-factorization-formula}.  The projection formula follows by
conjugating \(\WAtheta\WAtheta^*\).
\end{proof}

The aligned BRLC locus is therefore
\begin{equation}
  \UA\LBR
  =\projective{
      \ket{0}_{\mathrm{sec}}\otimes\actHP
    }.
  \label{eq:BRLC-factorized-locus}
\end{equation}
It is one projective slice of the equatorial neutral family.  The Awareness
bridge fixes the sector vector \(\ket0_{\mathrm{sec}}\), while the Born
geometry resides in the \(\actHP\) factor.

\subsection{Phase blindness of the paired algebra}

\begin{theorem}[Paired observables are phase blind]
\label{thm:paired-observables-phase-blind}
For every normalized \(\psi\in\actHP\), every \(X\in\obsP\), and every
\(\theta\),
\begin{equation}
  \inner{\WAtheta\psi}{\iotaA(X)\WAtheta\psi}_{\scrK}
  =\inner{\psi}{X\psi}_{\Psi}.
  \label{eq:phase-blind-expectation}
\end{equation}
Thus no observable in \(\obsPair\) distinguishes the aligned graph, the
anti-aligned graph, or any other relative-phase graph.
\end{theorem}

\begin{proof}
This is \cref{cor:paired-expectation-equality}.  Equivalently, in the
factorized representation,
\begin{equation}
  \inner{\theta|\idop_2|\theta}_{\mathrm{sec}}
  \inner{\psi}{X\psi}_{\Psi}
  =\inner{\psi}{X\psi}_{\Psi}.
\end{equation}
\end{proof}

The phase blindness is not a defect.  It expresses the fact that the paired
algebra represents the same physical--phenomenal proposition independently
of how the two descriptions are coordinatized by a relative phase.

\subsection{The commutant in the maximal case}

Assume for this subsection that
\begin{equation}
  \obsP=\mathcal B(\actHP).
  \label{eq:maximal-physical-observable-algebra}
\end{equation}

\begin{theorem}[Structural commutant]
\label{thm:structural-commutant}
Under \cref{eq:maximal-physical-observable-algebra},
\begin{equation}
  \UA\obsPair\UA^*
  =\idop_2\otimes\mathcal B(\actHP),
  \label{eq:paired-full-factor}
\end{equation}
and its commutant is
\begin{equation}
  \boxed{
  (\UA\obsPair\UA^*)'
  =\mathcal B(\CC^2)\otimes\idop_{\actHP}.}
  \label{eq:paired-full-commutant}
\end{equation}
In particular, the transformed operators corresponding to
\(\fundJ,\RA,-i\IA\), and the structural sector projections belong to this
commutant.
\end{theorem}

\begin{proof}
The first identity is \eqref{eq:paired-algebra-factorized}.  The commutant of
\(\idop_2\otimes\mathcal B(\actHP)\) in
\(\mathcal B(\CC^2\otimes\actHP)\) is
\(\mathcal B(\CC^2)\otimes\idop_{\actHP}\).  The remaining claim follows
from \cref{eq:J-factorized,eq:R-factorized,eq:Q-factorized}.
\end{proof}

This theorem separates two operator layers:
\begin{equation}
  \boxed{
  \begin{aligned}
  \idop_2\otimes X
  &\quad\text{acts on corresponding state content},\\
  S\otimes\idop
  &\quad\text{acts on structural sector relation and phase}.
  \end{aligned}}
  \label{eq:content-sector-operator-separation}
\end{equation}
The two layers commute in the maximal correspondence representation.

\subsection{Structural sector projections are not paired effects}

Under \(\UA\), the active sector projections are
\begin{equation}
  \UA\PP\UA^*
  =\proj{\eplus}\otimes\idop,
  \qquad
  \UA\PF\UA^*
  =\proj{\eminus}\otimes\idop.
  \label{eq:sector-projections-factorized}
\end{equation}
They are not of the form \(\idop_2\otimes X\), except in the trivial zero
space.  Therefore they are not elements of the paired observable algebra.

By the structural-sector-projection model choice \cref{choice:ssp},
\(\PP\) and \(\PF\) remain structural compressions of the common carrier;
they are not promoted to physical effects representing mutually exclusive
outcomes ``matter'' and ``mind.''

Under this choice, the normalized active graph vector
\begin{equation}
  \WA\psi
  =2^{-1/2}(\psi\oplus\polA\psi)
  \label{eq:SSP-total-state}
\end{equation}
is one relational state.  Its two components are recovered by compression:
\begin{equation}
  \sqrt2\,\PP\WA\psi=\psi,
  \qquad
  \sqrt2\,\PF\WA\psi=\polA\psi.
  \label{eq:SSP-compressions}
\end{equation}
The factors \(\sqrt2\) merely undo total-state normalization.

\begin{remark}[What the rejected measurement reading would say]
\label{rem:rejected-sector-measurement-reading}
If \(\PP\) and \(\PF\) were instead admitted as a two-outcome PVM, every
normalized neutral graph state would yield
\begin{equation}
  \Pr(\Psi)=\Pr(\Phi)=\frac12.
\end{equation}
That would treat matter and mind as mutually exclusive measurement outcomes.
The SSP model choice rejects precisely this doctrinal reading.
\end{remark}

\subsection{Relative phase and cross-sector structure}

The paired algebra cannot distinguish \(\theta\).  The structural Pauli
operators can:
\begin{align}
  \inner{\WAtheta\psi}{\RA\WAtheta\psi}_{\scrK}
  &=\cos\theta,
  \label{eq:RA-phase-expectation-active}\\
  \inner{\WAtheta\psi}{(-i\IA)\WAtheta\psi}_{\scrK}
  &=\sin\theta,
  \label{eq:QA-phase-expectation-active}\\
  \inner{\WAtheta\psi}{\fundJ\WAtheta\psi}_{\scrK}
  &=0.
  \label{eq:J-neutral-expectation-active}
\end{align}
These values are independent of \(\psi\) because the operators act only on
the sector coordinate.

\begin{openquestion}[Operational status of structural phase]
\label{q:operational-structural-phase}
Are \(\RA\) and \(-i\IA\) operational observables, dynamical generators,
structural diagnostics, or only mathematical coordinatizations of the
cross-sector relation?  Their formal phase resolution does not by itself
assign a laboratory measurement protocol.
\end{openquestion}

The Awareness bridge selects the aligned graph through the channel itself.
Operational access to the selecting phase is therefore not required for the
BRLC to be mathematically well defined.

\subsection{Doctrinal cash value of SSP}

\begin{identification}[One event, two conventional compressions]
\label{id:one-event-two-compressions}
The aligned relational ray \([\WA\psi]\) is interpreted as one cognitive
event in the common carrier.  Its physical compression \([\psi]\) and
phenomenal compression \([\polA\psi]\) are distinct conventional
descriptions or aspects of that event, related by an exact preservation of
Born geometry.
\end{identification}

This has four consequences.
\begin{enumerate}[label=(\roman*)]
  \item Matter and phenomenal presentation remain mathematically
  distinguishable because the carrier is graded.
  \item They are not treated as independent substances because admissible
  relational states lie on a graph rather than in an unconstrained product
  of independent coordinates.
  \item They are not treated as mutually exclusive outcomes because the
  sector projections are structural, not measurement effects.
  \item Their corresponding probability assignments agree because the polar
  orientation and paired algebra are unitarily equivalent representations.
\end{enumerate}

The BRLC therefore formalizes nonduality neither as identity nor as a
fifty--fifty superposition between competing ontologies.  It formalizes an
isometric relation between distinguishable aspects of one state.
% =============================================================================
% END INLINED FILE: sections/19_ssp_factorization_phase.tex
% =============================================================================
% =============================================================================
% BEGIN INLINED FILE: sections/21_phenomenal_determination_gns.tex
% =============================================================================
\section{The phenomenal determination algebra and the appearance space}
\label{sec:phenomenal-determination-gns}

The luminous ground vector is cyclic for the selected phenomenal algebra:
\begin{equation}
  \overline{\piF(\algF)\ket{\Omega_0}}=\scrHF.
  \label{eq:full-phenomenal-cyclicity}
\end{equation}
The subalgebra below isolates the determinations used to generate
differentiated appearance and makes the corresponding quotient geometry
explicit.

\subsection{The selected determination algebra}

\begin{modelchoice}[Concrete determination layer]
\label{choice:determination-algebra}
Choose a unital $C^*$-subalgebra
\begin{equation}
  \detAlg\subseteq\algF
  \label{eq:determination-algebra-subalgebra}
\end{equation}
whose represented action contains the admissible phenomenal determinations
used in the cognitive realization.  Unless otherwise stated, the model assumes
that $\ket{\Omega_0}$ is cyclic for this selected algebra:
\begin{equation}
  \overline{\piF(\detAlg)\ket{\Omega_0}}=\scrHF.
  \label{eq:determination-algebra-cyclicity}
\end{equation}
\end{modelchoice}

The full algebra $\algF$ remains available.  The point of
\cref{choice:determination-algebra} is to avoid making the cyclicity claim
vacuous by automatically taking every bounded operator on $\scrHF$ as an
admissible phenomenal determination.

\subsection{The GNS null ideal}

Define
\begin{equation}
  \boxed{
  \nullIdeal
  :=\left\{
      X\in\detAlg:
      \omegaF(X^*X)=0
    \right\}.}
  \label{eq:phenomenal-null-ideal}
\end{equation}
Since
\begin{equation}
  \omegaF(X^*X)
  =\|\piF(X)\ket{\Omega_0}\|^2,
  \label{eq:null-ideal-vector-norm}
\end{equation}
we also have
\begin{equation}
  \nullIdeal
  =\left\{
      X\in\detAlg:
      \piF(X)\ket{\Omega_0}=0
    \right\}.
  \label{eq:null-ideal-vector-kernel}
\end{equation}

\begin{proposition}[Ground-null operators form a left ideal]
\label{prop:null-left-ideal}
The set $\nullIdeal$ is a norm-closed left ideal of $\detAlg$.
\end{proposition}

\begin{proof}
If $X\in\nullIdeal$ and $Y\in\detAlg$, then
\begin{equation}
  \omegaF((YX)^*(YX))
  =\omegaF(X^*Y^*YX)
  \leq\|Y\|^2\omegaF(X^*X)
  =0.
\end{equation}
Closedness follows from continuity of the state functional.
\end{proof}

The quotient $\detAlg/\nullIdeal$ carries the pre-inner product
\begin{equation}
  \inner{[X]}{[Y]}_0
  :=\omegaF(X^*Y).
  \label{eq:GNS-preinner}
\end{equation}
Its completion is the GNS space of the restricted state
$\omegaF\restr{}{\detAlg}$.

\begin{theorem}[Concrete GNS realization of phenomenal appearance]
\label{thm:concrete-GNS-appearance}
The map
\begin{equation}
  U_0:[X]\longmapsto\piF(X)\ket{\Omega_0}
  \label{eq:GNS-concrete-map}
\end{equation}
is an isometry from $\detAlg/\nullIdeal$ into $\scrHF$.  Under
\cref{eq:determination-algebra-cyclicity}, it extends uniquely to a unitary
onto $\scrHF$.
\end{theorem}

\begin{proof}
Equation \eqref{eq:GNS-preinner} gives
\begin{equation}
  \inner{U_0[X]}{U_0[Y]}_{\Phi}
  =\inner{\piF(X)\Omega_0}{\piF(Y)\Omega_0}_{\Phi}
  =\omegaF(X^*Y).
\end{equation}
The kernel is exactly $\nullIdeal$, and cyclicity gives dense range.
\end{proof}

Thus a phenomenal vector is not fundamentally an additional substance placed
next to luminosity.  It is represented by an equivalence class of
phenomenal determinations acting on the luminous ground.

\subsection{Centered determinations}

Define the centered linear subspace
\begin{equation}
  \boxed{
  \centDet
  :=\ker\omegaF
  =\left\{
      X\in\detAlg:\omegaF(X)=0
    \right\}.}
  \label{eq:centered-determination-space}
\end{equation}
For every $X\in\detAlg$, write
\begin{equation}
  X=\omegaF(X)\idop+X^{\circ},
  \qquad
  X^{\circ}:=X-\omegaF(X)\idop\in\centDet.
  \label{eq:centered-determination-decomposition}
\end{equation}
Then
\begin{equation}
  \piF(X)\ket{\Omega_0}
  =\omegaF(X)\ket{\Omega_0}
  +\piF(X^{\circ})\ket{\Omega_0},
  \label{eq:ground-content-vector-decomposition}
\end{equation}
with
\begin{equation}
  \braket{\Omega_0}{\piF(X^{\circ})\Omega_0}=0.
  \label{eq:centered-vector-orthogonality}
\end{equation}

\warninglabel The space $\centDet$ is generally neither a subalgebra nor an
ideal.  Centering is a linear decomposition relative to the ground state, not
a new multiplication law.

\begin{definition}[Differentiated appearance space]
\label{def:differentiated-appearance-space}
The differentiated appearance space is
\begin{equation}
  \boxed{
  \appSpace
  :=\overline{
      \left\{
        \piF(X^{\circ})\ket{\Omega_0}:
        X\in\detAlg
      \right\}
    }.}
  \label{eq:differentiated-appearance-space}
\end{equation}
\end{definition}

\begin{theorem}[Centered determinations generate the ground complement]
\label{thm:centered-generates-ground-complement}
Under \eqref{eq:determination-algebra-cyclicity},
\begin{equation}
  \boxed{
  \appSpace=\QzeroF\scrHF.}
  \label{eq:appearance-space-ground-complement}
\end{equation}
\end{theorem}

\begin{proof}
Every centered vector is orthogonal to $\ket{\Omega_0}$, so
$\appSpace\subseteq\QzeroF\scrHF$.  Conversely, let $\eta\perp\Omega_0$.
By cyclicity there is a sequence $X_n\in\detAlg$ such that
$\piF(X_n)\Omega_0\to\eta$.  Taking the inner product with $\Omega_0$ gives
$\omegaF(X_n)\to0$.  Hence
\begin{equation}
  \piF(X_n^{\circ})\Omega_0
  =\piF(X_n)\Omega_0-\omegaF(X_n)\Omega_0
  \longrightarrow\eta.
\end{equation}
\end{proof}

\subsection{The covariance kernel}

For $X,Y\in\detAlg$, define
\begin{equation}
  \boxed{
  \Cov_{0}^{\Phi}(X,Y)
  :=\omegaF(X^*Y)
   -\overline{\omegaF(X)}\,\omegaF(Y).}
  \label{eq:phenomenal-covariance-kernel}
\end{equation}

\begin{proposition}[Covariance is the inner product of differentiated content]
\label{prop:covariance-content-inner-product}
For all $X,Y\in\detAlg$,
\begin{equation}
  \Cov_{0}^{\Phi}(X,Y)
  =\inner{\piF(X^{\circ})\Omega_0}
          {\piF(Y^{\circ})\Omega_0}_{\Phi}.
  \label{eq:covariance-content-identity}
\end{equation}
Consequently $\Cov_{0}^{\Phi}$ is positive semidefinite, and
\begin{equation}
  \Cov_{0}^{\Phi}(X,X)
  =\|\QzeroF\piF(X)\Omega_0\|^2.
  \label{eq:covariance-appearance-norm}
\end{equation}
\end{proposition}

\begin{proof}
Expand $X^{\circ}$ and $Y^{\circ}$ using
\eqref{eq:centered-determination-decomposition}.
\end{proof}

For a finite family $X_1,\dots,X_n$, the matrix
\begin{equation}
  C_{jk}:=\Cov_{0}^{\Phi}(X_j,X_k)
  \label{eq:appearance-Gram-matrix}
\end{equation}
is the Gram matrix of the differentiated vectors
$\piF(X_j^{\circ})\Omega_0$.  Its rank is therefore the dimension of their
linear span.  This supplies a canonical local geometry of differentiated
appearance without identifying differentiation with a loss of luminosity.

\subsection{Normalized appearance states}

Whenever $\piF(X)\Omega_0\neq0$, define
\begin{equation}
  \ket{\phi_X}
  :=\frac{\piF(X)\ket{\Omega_0}}
          {\|\piF(X)\ket{\Omega_0}\|}.
  \label{eq:normalized-appearance-state}
\end{equation}
The ground fidelity is
\begin{equation}
  F_0^{\Phi}(\phi_X)
  =\frac{|\omegaF(X)|^2}
         {|\omegaF(X)|^2+\Cov_0^{\Phi}(X,X)}.
  \label{eq:appearance-ground-fidelity-covariance}
\end{equation}
A centered determination with nonzero action gives
\begin{equation}
  F_0^{\Phi}(\phi_X)=0.
  \label{eq:centered-appearance-orthogonal}
\end{equation}
Such a state is fully orthogonal to the undifferentiated ground while still
being generated from that ground.

\begin{identification}[All appearances are luminosity]
\label{id:all-appearances-luminosity}
The cyclic quotient representation
\begin{equation}
  \scrHF\simeq\overline{\detAlg/\nullIdeal}
\end{equation}
is interpreted as the statement that every phenomenal appearance is an
algebraic determination of luminosity.  Orthogonality to
$\ket{\Omega_0}$ measures absence of an undifferentiated ground component,
not absence of luminosity.
\end{identification}

This distinction follows the doctrinal analysis according to which all
cognitions remain reflexively luminous even when differentiated appearance
and dualistic error prevent luminosity from being ascertained as such
\cite[pp.~162--163]{Yiannopoulos2012}; see also
\cite[pp.~424--425]{Yiannopoulos2020}.
% =============================================================================
% END INLINED FILE: sections/21_phenomenal_determination_gns.tex
% =============================================================================
% =============================================================================
% BEGIN INLINED FILE: sections/22_vasana_generator_system.tex
% =============================================================================
\section{V\=asan\=a-conditioned generators and the word algebra of appearance}
\label{sec:vasana-generator-system}

The preceding section identified the complete differentiated appearance space
as the centered cyclic image of a selected determination algebra.  The next
question is which operators generate that algebra and how latent dispositions
condition their action.

\subsection{V\=asan\=a as conditioning data}

\begin{modelchoice}[Disposition profiles]
\label{choice:vasana-profiles}
Let $\vasSpace$ be a set of disposition profiles.  A profile
$\vas\in\vasSpace$ does not, by itself, denote a phenomenal state or a
single operator.  It indexes a family of admissible determination generators
\begin{equation}
  \genSet(\vas)
  =\left\{G_a(\vas):a\in\mathcal A\right\}
  \subseteq\detAlg.
  \label{eq:vasana-generator-family}
\end{equation}
The profile records which generators are active and how strongly or in what
form they act in the cognitive continuum under consideration.
\end{modelchoice}

This choice deliberately avoids equating a v\=asan\=a with a vector in
$\scrHF$.  A disposition is represented by the way it conditions generation
of appearances, not by appearing as one more phenomenal content.

Define the unital $*$-algebra of finite words
\begin{equation}
  \wordSet(\vas)
  :=\Alg^*\!\left(
      \idop,\genSet(\vas)
    \right)
  \subseteq\detAlg,
  \label{eq:vasana-word-algebra}
\end{equation}
and the profile-generated phenomenal subspace
\begin{equation}
  \boxed{
  \Hvas
  :=\overline{\piF(\wordSet(\vas))\ket{\Omega_0}}.}
  \label{eq:profile-generated-phenomenal-subspace}
\end{equation}

\begin{definition}[Phenomenally complete profile]
\label{def:phenomenally-complete-profile}
A profile $\vas$ is phenomenally complete when
\begin{equation}
  \Hvas=\scrHF.
  \label{eq:profile-completeness}
\end{equation}
\end{definition}

The full determination algebra can be generated by the union of all profile
families even when no single ordinary profile is phenomenally complete.

\subsection{Generator classes}

The realization distinguishes five functional classes:
\begin{equation}
  \boxed{
  \genSet(\vas)
  =\Gak(\vas)
   \cup\Gvik(\vas)
   \cup\Gsmr(\vas)
   \cup\Gaff(\vas)
   \cup\Gdv(\vas).}
  \label{eq:five-generator-classes}
\end{equation}
The classes need not be disjoint and no commutation relations are assumed.
Their intended roles are:
\begin{enumerate}[label=(\roman*)]
  \item $\Gak$: generators of differentiated cognitive images or
  \emph{\=ak\=aras}, including vivid nonconceptual appearance;
  \item $\Gvik$: conceptual determination, exclusion, convention, and
  superimposition (\emph{vikalpa}, \emph{apoha}, \emph{\=aropa});
  \item $\Gsmr$: memory, recognition, and temporal concatenation
  (\emph{sm\d{r}ti}, \emph{pratyabhij\~n\=a});
  \item $\Gaff$: affective and perspectival modulation, including pleasure,
  pain, desire, and related first-order phenomenal variation;
  \item $\Gdv$: the nonconceptual internal distortion producing the derived
  grasper--grasped frame.
\end{enumerate}

\begin{identification}[Doctrinal source of the generator classes]
\label{id:generator-doctrinal-source}
The five classes formalize the dissertation's reconstruction of ordinary
cognition as v\=asan\=a-conditioned, image-bearing, mnemonic, affectively
qualified, and necessarily structured by erroneous phenomenological duality
\cite[pp.~27--32, 64--70, 107--136, 147--156, 414--456]{Yiannopoulos2020}.
\end{identification}

The classification is functional rather than ontological.  A single operator
may, for example, both generate an image and modulate its affective character.
Likewise, the formalism does not yet assume that conceptual and
nonconceptual operations act on orthogonal Hilbert subspaces.

\subsection{Words as conditioned histories}

A word
\begin{equation}
  W_{\mathbf a}(\vas)
  :=G_{a_n}(\vas)\cdots G_{a_2}(\vas)G_{a_1}(\vas)
  \label{eq:conditioned-generator-word}
\end{equation}
encodes an ordered history of determination.  Its generated state is
\begin{equation}
  \ket{\phi_{\mathbf a,\vas}}
  \propto
  \piF(W_{\mathbf a}(\vas))\ket{\Omega_0}
  \label{eq:word-generated-appearance}
\end{equation}
whenever the vector is nonzero.

\begin{proposition}[Order sensitivity]
\label{prop:word-order-sensitivity}
For two generators $G_a(\vas),G_b(\vas)$,
\begin{equation}
  \piF(G_bG_a)\ket{\Omega_0}
  -\piF(G_aG_b)\ket{\Omega_0}
  =\piF([G_b,G_a])\ket{\Omega_0}.
  \label{eq:word-order-commutator}
\end{equation}
Hence noncommutativity records dependence of the phenomenal result on the
order of conditioning operations.
\end{proposition}

No claim is made that every cognitive process is literally a discrete word.
The word algebra is the minimal algebraic realization.  Continuous-time,
simultaneous, and completely positive dynamical versions may instead be
defined from the same generator families.

\subsection{Profile-relative differentiated appearance}

For $X\in\wordSet(\vas)$, define
\begin{equation}
  X^{\circ}=X-\omegaF(X)\idop.
\end{equation}
The profile-relative appearance space is
\begin{equation}
  \boxed{
  \Evas
  :=\overline{
      \left\{
        \piF(X^{\circ})\ket{\Omega_0}:
        X\in\wordSet(\vas)
      \right\}
    }.}
  \label{eq:profile-relative-appearance-space}
\end{equation}
It satisfies
\begin{equation}
  \Evas\subseteq\QzeroF\Hvas.
  \label{eq:profile-appearance-inclusion}
\end{equation}
If $\ket{\Omega_0}$ is cyclic for $\wordSet(\vas)$, then equality holds by
\cref{thm:centered-generates-ground-complement} applied to the generated
representation.

For a word $X$, define its differentiated production at the ground by
\begin{equation}
  \epsilon_0(X)
  :=\|\QzeroF\piF(X)\ket{\Omega_0}\|^2
  =\Cov_0^{\Phi}(X,X).
  \label{eq:differentiated-production-functional}
\end{equation}
This quantity measures how much ground-orthogonal appearance the word
produces from $\ket{\Omega_0}$.  It is not an amount of luminosity.

\subsection{The exhausted profile}

\begin{modelchoice}[Exhaustion acts trivially on the ground ray]
\label{choice:exhausted-profile-ground-triviality}
Distinguish an exhausted disposition profile
$\vasEmpty\in\vasSpace$.  Its active generators satisfy
\begin{equation}
  \boxed{
  \QzeroF\piF(G_a(\vasEmpty))\ket{\Omega_0}=0,
  \qquad
  \QzeroF\piF(G_a(\vasEmpty)^*)\ket{\Omega_0}=0
  \quad
  \text{for every active }a.}
  \label{eq:exhausted-generator-ground-triviality}
\end{equation}
Equivalently, the ground line reduces every active generator: both the generator and its adjoint preserve that line.
\end{modelchoice}

This condition does not assert that the universal determination algebra has
vanished, nor that each generator is the zero operator.  It says only that the
currently active v\=asan\=a-conditioned operations produce no differentiated
content from the realized ground.

\begin{proposition}[Ground preservation by exhausted words]
\label{prop:exhausted-words-ground-preserve}
Under \cref{choice:exhausted-profile-ground-triviality}, every word
$W\in\wordSet(\vasEmpty)$ satisfies
\begin{equation}
  \piF(W)\ket{\Omega_0}
  \in\CC\ket{\Omega_0}.
  \label{eq:exhausted-word-ground-line}
\end{equation}
Consequently
\begin{equation}
  \HvasEmpty=\CC\ket{\Omega_0}
  \label{eq:exhausted-profile-subspace}
\end{equation}
for the algebra generated only by the active exhausted-profile operations.
\end{proposition}

\begin{proof}
Each generator maps the one-dimensional ground line into itself, and the same
is true of adjoints and finite products.  Taking the closure gives the result.
\end{proof}

\warninglabel Equation \eqref{eq:exhausted-profile-subspace} concerns the
active profile-relative algebra, not the full latent algebra $\detAlg$.  The
latter remains cyclic and contains the possible differentiated determinations
that are inactive in the exhausted condition.
% =============================================================================
% END INLINED FILE: sections/22_vasana_generator_system.tex
% =============================================================================
% =============================================================================
% BEGIN INLINED FILE: sections/23_operator_stages_ordinary_appearance.tex
% =============================================================================
\section{Operator stages of ordinary appearance}
\label{sec:operator-stages-appearance}

The generator classification becomes useful only if it clarifies how ordinary
phenomenal content is organized.  The realization therefore records a minimal
word model while refusing to turn that model into an unnecessarily rigid
temporal sequence.

\subsection{Appearance before determination}

Choose a profile $\vas$ and an image-forming generator
\begin{equation}
  K\in\Gak(\vas).
\end{equation}
The corresponding primary appearance vector is
\begin{equation}
  \ket{\phi_K}
  \propto\piF(K)\ket{\Omega_0}.
  \label{eq:primary-akara-state}
\end{equation}
The vector is already reflexively manifest because it lies in the phenomenal
sector.  No second Awareness action is required.  Whether $K$ is perceptual,
vivid, erroneous, or instrumentally reliable is supplied by the realization
of the generator class; none of those properties is identified with Hilbert
norm alone.

\begin{identification}[\=Ak\=ara as differentiated luminosity]
\label{id:akara-differentiated-luminosity}
The cognitive image is represented by the ground-orthogonal component
\begin{equation}
  \QzeroF\piF(K)\ket{\Omega_0}
  =\piF(K^{\circ})\ket{\Omega_0}.
  \label{eq:akara-centered-component}
\end{equation}
It is a determination of luminosity, not an independently existing object
placed before awareness.
\end{identification}

This matches the doctrinal claim that awareness is ultimately aware only of
awareness, even when its internally generated content appears as an object
\cite[pp.~31--32]{Yiannopoulos2020}.

\subsection{Conceptual exclusion and superimposition}

The dissertation reconstructs conceptual determination as a two-sided
operation: causal or phenomenal features are excluded, and a conventionally
constructed unity is then projected or superimposed upon experience
\cite[pp.~64--70, 90--94]{Yiannopoulos2020}.

\begin{modelchoice}[Filter--superimposition realization of apoha]
\label{choice:apoha-filter-superimposition}
A conceptual determination indexed by a convention $\lambda$ may be modeled
by
\begin{equation}
  C_{\lambda}=S_{\lambda}F_{\lambda},
  \label{eq:apoha-two-stage-operator}
\end{equation}
where
\begin{equation}
  0\leq F_{\lambda}\leq\idop
  \label{eq:apoha-filter-effect}
\end{equation}
is a bounded effect representing selective exclusion and
$S_{\lambda}$ is a bounded superimposition operator representing the
application of the conventional determination.
\end{modelchoice}

Applied to a primary appearance,
\begin{equation}
  \ket{\phi_{\lambda,K}}
  \propto
  \piF(C_{\lambda}K)\ket{\Omega_0}.
  \label{eq:conceptually-determined-appearance}
\end{equation}
Neither $F_{\lambda}$ nor $S_{\lambda}$ is assumed to preserve norm; the
resulting nonzero vector is normalized only after the complete determination.
The decomposition is a model choice, not a claim that the historical term
\emph{apoha} already denotes a Hilbert-space projection.

\subsection{Memory, recognition, and concatenation}

Let
\begin{equation}
  M\in\Gsmr(\vas)
\end{equation}
be a mnemonic generator.  A history-sensitive determination can be encoded
by a word such as
\begin{equation}
  X_{\mathrm{rec}}
  =C_{\lambda}M K_n\cdots K_1.
  \label{eq:recognition-word}
\end{equation}
Its state is
\begin{equation}
  \ket{\phi_{\mathrm{rec}}}
  \propto
  \piF(X_{\mathrm{rec}})\ket{\Omega_0}.
  \label{eq:recognition-state}
\end{equation}
The order of the $K_j$ can matter by
\cref{prop:word-order-sensitivity}.  This provides a minimal algebraic
representation of mnemonic concatenation and object persistence without
assuming that a persisting object exists independently of the sequence of
momentary cognitions.

The model is compatible with the dissertation's account of ordinary object
persistence and recognition as conceptual, mnemonic concatenations of
momentary cognition \cite[pp.~107--128]{Yiannopoulos2020}.

\subsection{Conceptual and nonconceptual tracks}

Define the norm-closed operator systems generated by the nonconceptual and
conceptual classes:
\begin{align}
  \mathcal S_{\mathrm{nc}}(\vas)
  &:=\overline{\Span\bigl(
      \idop,\Gak(\vas),\Gaff(\vas),\Gdv(\vas),
      \text{adjoints}
    \bigr)}^{\|\cdot\|},
  \label{eq:nonconceptual-operator-system}\\
  \mathcal S_{\mathrm{c}}(\vas)
  &:=\overline{\Span\bigl(
      \idop,\Gvik(\vas),\Gsmr(\vas),
      \text{adjoints}
    \bigr)}^{\|\cdot\|}.
  \label{eq:conceptual-operator-system}
\end{align}
They are self-adjoint unital closed linear subspaces; no multiplicative
closure is asserted.

\begin{modelchoice}[Tracks are operator types, not Hilbert sectors]
\label{choice:tracks-not-sectors}
The conceptual and nonconceptual tracks are represented as classes of
operations acting on the same phenomenal Hilbert space.  No orthogonal
decomposition is assumed:
\begin{equation}
  \scrHF
  =\scrHF^{\mathrm{nc}}\oplus\scrHF^{\mathrm{c}}.
  \label{eq:rejected-conceptual-sector-split}
\end{equation}
\end{modelchoice}

This permits the simultaneous or rapidly interwoven operation of the two
tracks without reifying them into independent consciousnesses.  No general
commutation relation between
$\mathcal S_{\mathrm{nc}}(\vas)$ and
$\mathcal S_{\mathrm{c}}(\vas)$ is imposed.

\subsection{Affective and perspectival variation}

To model variation in the subjective or affective quality of an appearance,
choose two represented subalgebras
\begin{equation}
  \contentAlg,\affectAlg\subseteq\piF(\detAlg)''.
  \label{eq:content-affect-subalgebras}
\end{equation}
They need not commute and need not generate the full phenomenal algebra.

\begin{definition}[Content-equivalent states]
\label{def:content-equivalent-states}
Two phenomenal density operators $\rho_1,\rho_2$ are content-equivalent relative
to $\contentAlg$ when
\begin{equation}
  \Tr(\rho_1 X)=\Tr(\rho_2 X)
  \quad
  \text{for all }X\in\contentAlg.
  \label{eq:content-equivalence}
\end{equation}
\end{definition}

An affective generator may move a state within a content-equivalence class
while changing its restriction to $\affectAlg$.  Thus the same conventionally
identified object-content can occur with different pleasure, pain, sharpness,
dullness, or first-personal orientation.  This captures the doctrinal point
that subjective and affective variation belongs to the phenomenal event
itself, not to a separate substantial subject observing an otherwise
unchanged inner object \cite[pp.~414--456]{Yiannopoulos2020}.

\begin{openquestion}[Choice of phenomenal subalgebras]
\label{q:content-affect-algebra-choice}
Which represented subalgebras should count as object-content and affective
structure in a concrete realization?  The present construction defines the
relation but does not select a universal decomposition.
\end{openquestion}
% =============================================================================
% END INLINED FILE: sections/23_operator_stages_ordinary_appearance.tex
% =============================================================================
% =============================================================================
% BEGIN INLINED FILE: sections/24_derived_duality_frames.tex
% =============================================================================
\section{Derived grasper--grasped frames inside the phenomenal sector}
\label{sec:derived-duality-frames}

The fundamental matter--mind grading is carried by $\PP$ and $\PF$.  The
ordinary phenomenal split between grasper and grasped must therefore be a
second, derived structure entirely internal to $\scrHF$.

\subsection{A state-dependent internal frame}

\begin{definition}[Derived duality frame]
\label{def:derived-duality-frame}
Let $\ket\phi\in\scrHF$ be normalized with
$\QzeroF\phi\neq0$.  A derived duality frame for $\phi$ is a triple of
orthogonal projections
\begin{equation}
  (\Pgrasper,\Pgrasped,\Pdual)
  \label{eq:duality-frame-triple}
\end{equation}
satisfying
\begin{equation}
  \Pgrasper\Pgrasped=0,
  \qquad
  \Pgrasper+\Pgrasped=\Pdual,
  \qquad
  \Pdual\leq\QzeroF.
  \label{eq:duality-frame-relations}
\end{equation}
The two projections are interpreted, respectively, as the derived grasper
and grasped aspects of differentiated phenomenal content.
\end{definition}

The corresponding internal grading on the dualized support is
\begin{equation}
  \boxed{
  \Gammadual
  :=\Pgrasper-\Pgrasped.}
  \label{eq:derived-duality-grading}
\end{equation}
It satisfies
\begin{equation}
  (\Gammadual)^* =\Gammadual,
  \qquad
  (\Gammadual)^2=\Pdual.
  \label{eq:derived-duality-grading-properties}
\end{equation}

\begin{modelchoice}[V\=asan\=a dependence of the duality frame]
\label{choice:vasana-dependent-duality-frame}
The frame is not fixed as part of the universal carrier.  It is selected by
the active disposition profile and phenomenal state:
\begin{equation}
  (\Pgrasper,\Pgrasped)
  =\mathcal D(\vas,[\phi]).
  \label{eq:duality-frame-selection-map}
\end{equation}
No global decomposition of $\scrHF$ into permanent subject and object sectors
is assumed.
\end{modelchoice}

\subsection{The split is internal, not ontological}

For the differentiated component
\begin{equation}
  \eta_{\phi}:=\QzeroF\ket\phi,
  \label{eq:differentiated-component-duality}
\end{equation}
the frame yields
\begin{equation}
  \Pdual\eta_{\phi}
  =\Pgrasper\eta_{\phi}
   +\Pgrasped\eta_{\phi}.
  \label{eq:internal-grasper-grasped-split}
\end{equation}
Both components lie in $\scrHF$.  They do not occupy the physical and
phenomenal sectors of $\scrK$, and neither is an independently existing
subject or object.

\begin{identification}[Location of phenomenological duality]
\label{id:derived-duality-location}
The gr\=ahaka and gr\=ahya aspects are represented as a v\=asan\=a-conditioned
internal frame on differentiated phenomenal content.  Their apparent
opposition is therefore real as phenomenal organization but not as an
ultimate bifurcation of the carrier.
\end{identification}

This formalizes the dissertation's conclusion that cognition is
ontologically singular even though ordinary cognition appears with subjective
and objective aspects \cite[pp.~147--156, 220--228, 283--289]{Yiannopoulos2020}.

\subsection{No duality frame at the ground}

\begin{theorem}[The derived duality frame vanishes on the luminous ground]
\label{thm:duality-frame-vanishes-ground}
For every derived duality frame satisfying
\eqref{eq:duality-frame-relations},
\begin{equation}
  \Pgrasper\ket{\Omega_0}
  =\Pgrasped\ket{\Omega_0}
  =\Pdual\ket{\Omega_0}
  =0.
  \label{eq:duality-frame-annihilates-ground}
\end{equation}
\end{theorem}

\begin{proof}
Since $\Pdual\leq\QzeroF$ and
$\QzeroF\ket{\Omega_0}=0$, the support projection annihilates the ground.
The two subprojections do likewise.
\end{proof}

Thus the luminous ground is not a balanced superposition of an internal
subject and object.  The derived frame has no nonzero support there at all.
This matches the disappearance of differentiated marks and subject--object
structure at the nondual luminous ground
\cite[pp.~182--185]{Yiannopoulos2012}.

\subsection{Duality as distortion rather than conceptual judgment}

The frame is compatible with the nonconceptual status of phenomenological
duality.  It can be generated by $\Gdv(\vas)$ before or independently of
subsequent conceptual operators in $\Gvik(\vas)$.  Hence removing conceptual
determination alone need not eliminate the frame.

\begin{modelchoice}[Duality generator]
\label{choice:duality-generator}
A concrete realization may assign to each profile a bounded self-adjoint
operator
\begin{equation}
  D_{\vas}\in\Gdv(\vas)
  \label{eq:duality-generator}
\end{equation}
whose generated differentiated support determines
$\Pdual$ and whose spectral decomposition supplies
$\Pgrasper$ and $\Pgrasped$.
\end{modelchoice}

No such spectral realization is imposed here.  The abstract frame
records the minimum needed to prevent the grasper--grasped split from being
confused with the fundamental Krein grading.

\subsection{No substantial subject ket}

\begin{corollary}[No independent subject state follows from the frame]
\label{cor:no-independent-subject-ket}
The existence of $\Pgrasper\phi$ does not imply a normalized state that exists
independently of the total phenomenal event.  It is a compression of the
differentiated content of $\phi$ relative to a derived frame.
\end{corollary}

This is especially important for the doctrinal claim that reflexive awareness
is not an inner subject turning back upon itself.  The self-presenting
character belongs to the whole phenomenal event; the grasper image is one
adventitious organization within that event.

\begin{openquestion}[Projection frame or weaker effect frame]
\label{q:projection-or-effect-duality-frame}
Should a concrete realization use sharp projections as above, a pair of
unsharp effects, a state-dependent tensor factorization, or a more general
operator-valued decomposition?  Projection frames give the cleanest minimal
model used here, but they are not claimed to be uniquely forced.
\end{openquestion}
% =============================================================================
% END INLINED FILE: sections/24_derived_duality_frames.tex
% =============================================================================
% =============================================================================
% BEGIN INLINED FILE: sections/25_sems_exhaustion_ground_recovery.tex
% =============================================================================
\section{Sems, v\=asan\=a exhaustion, and exact ground recovery}
\label{sec:sems-exhaustion-ground-recovery}

The algebra of appearance now permits a sharper distinction among intrinsic
reflexivity, differentiated mind, and exact realization of the ground.

\subsection{Sems is a projective condition, not one linear sector}

Define
\begin{equation}
  \boxed{
  \semsSet
  :=\left\{
      [\phi]\in\mathbb P(\scrHF):
      \QzeroF\phi\neq0
    \right\}.}
  \label{eq:sems-projective-set}
\end{equation}
This set contains all phenomenal rays with differentiated content.  It is not
a linear subspace: sums and scalar representatives are governed by the
projective geometry, and the omitted ground ray is not the zero vector.

Define the realized ground ray by
\begin{equation}
  \boxed{
  \rigRay:=\left\{[\Omega_0]\right\}.}
  \label{eq:rigpa-ground-ray-set}
\end{equation}
Then
\begin{equation}
  \mathbb P(\scrHF)
  =\semsSet\,\dot\cup\,\rigRay.
  \label{eq:projective-sems-rigpa-partition}
\end{equation}

\begin{identification}[Sems and rig pa]
\label{id:sems-rigpa}
The projective condition $\QzeroF\phi\neq0$ is interpreted as
\emph{sems}: differentiated phenomenal mind.  The exact condition
$[\phi]=[\Omega_0]$ is interpreted as \emph{rig pa}: realization of the
ground.
\end{identification}

The distinction is exact.  A state with arbitrarily high but nonunit ground
fidelity remains in $\semsSet$.

\subsection{Intrinsic reflexivity is common to both}

The odd Awareness extension satisfies
\begin{equation}
  \Aodd\iF\ket\phi=0
  \qquad
  \text{for every }\phi\in\scrHF.
  \label{eq:all-phenomenal-terminal}
\end{equation}
Thus the no-regress law of reflexive awareness applies throughout
$\semsSet$ and at $\rigRay$.

\begin{identification}[Rang rig versus rig pa]
\label{id:rangrig-rigpa}
\emph{Rang rig} is the intrinsic self-presenting terminality of phenomenal
states in general.  \emph{Rig pa} is that same reflexive awareness under the
additional exact ground condition
\begin{equation}
  \QzeroF\phi=0.
  \label{eq:rigpa-additional-condition}
\end{equation}
\end{identification}

This distinction gives mathematical form to the traditional claim that
reflexive awareness is present even when the nondual nature of awareness is
not definitively realized \cite[pp.~424--425]{Yiannopoulos2020}.

\subsection{Ground fidelity as non-adulteration, not luminosity quantity}

For normalized $\phi$, define
\begin{equation}
  \delta_0(\phi)
  :=1-F_0^{\Phi}(\phi)
  =\|\QzeroF\phi\|^2.
  \label{eq:adulteration-functional}
\end{equation}
Then
\begin{equation}
  \delta_0(\phi)=0
  \quad\Longleftrightarrow\quad
  [\phi]=[\Omega_0].
  \label{eq:adulteration-zero-ground}
\end{equation}

\begin{modelchoice}[Terminology of adulteration]
\label{choice:adulteration-terminology}
The scalar $\delta_0$ may be called ground deviation or phenomenal
adulteration.  It is not called the absence or inverse amount of luminosity.
Every state in the cyclic phenomenal representation remains luminous.
\end{modelchoice}

Ordinary states may satisfy
\begin{equation}
  F_0^{\Phi}(\phi)=0,
  \qquad
  \delta_0(\phi)=1.
  \label{eq:fully-ground-orthogonal-sems}
\end{equation}
They are then fully orthogonal to the undifferentiated ground while still
belonging to the centered cyclic appearance space $\appSpace$.

\subsection{Exact and asymptotic exhaustion}

Let $X_n\in\detAlg$ generate nonzero normalized states
\begin{equation}
  \ket{\phi_n}
  :=\frac{\piF(X_n)\ket{\Omega_0}}
          {\|\piF(X_n)\ket{\Omega_0}\|}.
  \label{eq:deconditioning-state-sequence}
\end{equation}
Write
\begin{equation}
  \piF(X_n)\ket{\Omega_0}
  =c_n\ket{\Omega_0}+\ket{\eta_n},
  \qquad
  c_n=\omegaF(X_n),
  \qquad
  \eta_n\perp\Omega_0.
  \label{eq:deconditioning-ground-content-split}
\end{equation}

\begin{theorem}[Projective ground recovery]
\label{thm:projective-ground-recovery}
Assume $c_n\neq0$ eventually and
\begin{equation}
  \frac{\|\eta_n\|}{|c_n|}\longrightarrow0.
  \label{eq:ground-recovery-ratio}
\end{equation}
Then
\begin{equation}
  F_0^{\Phi}(\phi_n)\longrightarrow1,
  \label{eq:ground-recovery-fidelity}
\end{equation}
and the rays $[\phi_n]$ converge to $[\Omega_0]$ in the Fubini--Study metric.
\end{theorem}

\begin{proof}
By orthogonality,
\begin{equation}
  F_0^{\Phi}(\phi_n)
  =\frac{|c_n|^2}{|c_n|^2+\|\eta_n\|^2}
  =\frac{1}{1+\|\eta_n\|^2/|c_n|^2}.
\end{equation}
The fidelity therefore converges to one.  Fubini--Study convergence follows
from
$d_{\mathrm{FS}}([\phi_n],[\Omega_0])
=\arccos\sqrt{F_0^{\Phi}(\phi_n)}$.
\end{proof}

\begin{definition}[Exact exhaustion in a generated state]
\label{def:exact-exhaustion-generated-state}
A generated state is exactly exhausted when
\begin{equation}
  \eta_X
  :=\piF(X^{\circ})\ket{\Omega_0}=0
  \label{eq:exact-exhaustion-centered-zero}
\end{equation}
and $\omegaF(X)\neq0$.
\end{definition}

Exact exhaustion gives
\begin{equation}
  [\phi_X]=[\Omega_0].
  \label{eq:exact-exhaustion-ground-ray}
\end{equation}
Thus no second ket is needed to represent a state beyond the luminous ground.

\subsection{The Third and Fourth Yogas revisited}

The MA thesis distinguishes the Third Yoga, where only undifferentiated
luminosity remains, from the Fourth Yoga, where even luminosity is not
objectified as an appearance \cite[pp.~182--185]{Yiannopoulos2012}.  Both use
one projective ground condition:
\begin{equation}
  [\phi]=[\Omega_0].
  \label{eq:third-fourth-same-ground-ray}
\end{equation}
The difference is representational or praxeological, not a transition to a
second vector below the ground.

Exact realization is the state condition $[\phi]=[\Omega_0]$.  Recognition,
stabilization, and irreversible v\=asan\=a exhaustion are distinct dynamical
questions treated in
\cref{sec:reflexive-terminality-recognition,sec:deconditioning-dynamics,sec:momentary-stabilization}.
% =============================================================================
% END INLINED FILE: sections/25_sems_exhaustion_ground_recovery.tex
% =============================================================================
% =============================================================================
% BEGIN INLINED FILE: sections/27_reflexive_terminality_recognition.tex
% =============================================================================
\section{Reflexive terminality and exact realization}
\label{sec:reflexive-terminality-recognition}

The universal reflexive character of phenomenal states is distinct from exact
realization of the luminous ground.  This section sharpens that distinction
and extends it from pure vectors to general normal states.

\subsection{The odd terminality law}

The grading-odd Awareness orientation satisfies
\begin{equation}
  \Aodd\iF=0,
  \qquad
  \Aodd^2=0.
  \label{eq:odd-terminal-law}
\end{equation}
The first identity states that the forward Awareness orientation has no second
phenomenal-to-further-phenomenal leg.  The second is its global nilpotency on
the graded carrier.

\begin{proposition}[No iterated Awareness tower]
\label{prop:no-iterated-awareness-tower}
For every physical vector $\psi\in\Dom A$,
\begin{equation}
  \Aodd\iP\psi\in\iF\scrHF,
  \qquad
  \Aodd^2\iP\psi=0.
  \label{eq:no-awareness-tower}
\end{equation}
For every phenomenal vector $\phi\in\scrHF$,
\begin{equation}
  \Aodd\iF\phi=0.
  \label{eq:phenomenal-terminal-all}
\end{equation}
Thus no sequence of distinct states
\begin{equation}
  \psi\longmapsto A\psi\longmapsto A(A\psi)\longmapsto\cdots
\end{equation}
is generated by repeated application of one Awareness orientation.
\end{proposition}

\begin{proof}
These are immediate consequences of the block form of $\Aodd$ and its
nilpotency.
\end{proof}

\begin{remark}[Technical annihilation is not ontological destruction]
\label{rem:annihilation-not-destruction}
The equation
\begin{equation}
  \Aodd\iF\ket{\Omega_0}=0
  \label{eq:A-annihilates-ground-technically}
\end{equation}
does say, in standard operator language, that $\Aodd$ annihilates the embedded
ground vector.  It does not describe a time evolution in which the luminous
ground is destroyed.  The output is zero because the typed forward map is
inapplicable after the phenomenal sector has already been reached.  State
evolution and ground preservation are separate structures developed in
\cref{sec:ground-nonalteration}.
\end{remark}

\begin{identification}[Algebraic shadow of rang rig]
\label{id:rangrig-terminal-law}
The universal terminality law \eqref{eq:odd-terminal-law} is interpreted
as the algebraic shadow of \emph{rang rig}: phenomenal manifestation does not
require a second act of Awareness to make the first act manifest.  It is not
claimed that nilpotency exhausts the philosophical meaning of reflexive
awareness.
\end{identification}

This interpretation is consistent with the dissertation's reconstruction of
reflexive awareness as intrinsic to every cognition and as the result through
which cognition is immediately present, irrespective of the causal ontology
of its differentiated content \cite[pp.~31, 424--425]{Yiannopoulos2020}.

\subsection{Reflexivity does not imply realization}

The terminality law holds for every state in $\scrHF$, including ordinary
states in $\semsSet$.  It therefore cannot by itself select the realized
ground.  The exact realization condition remains
\begin{equation}
  \QzeroF\phi=0
  \quad\Longleftrightarrow\quad
  [\phi]=[\Omega_0]
  \label{eq:exact-rigpa-pure}
\end{equation}
for nonzero pure vectors.

To include mixed states, let
\begin{equation}
  \mathfrak S(\scrHF)
  :=\left\{
      \rho\in\mathcal T_1(\scrHF):
      \rho\geq0,\ \Tr\rho=1
    \right\}
  \label{eq:phenomenal-density-state-space}
\end{equation}
be the normal phenomenal state space.  Define
\begin{equation}
  F_0^{\Phi}(\rho)
  :=\Tr(\PzeroF\rho),
  \qquad
  \delta_0^{\Phi}(\rho)
  :=1-F_0^{\Phi}(\rho)
  =\Tr(\QzeroF\rho).
  \label{eq:mixed-ground-fidelity-deviation}
\end{equation}

\begin{theorem}[Exact realization for normal states]
\label{thm:exact-rigpa-density}
For $\rho\in\mathfrak S(\scrHF)$, the following are equivalent:
\begin{enumerate}[label=(\roman*)]
  \item $F_0^{\Phi}(\rho)=1$;
  \item $\delta_0^{\Phi}(\rho)=0$;
  \item $\rho=\PzeroF$.
\end{enumerate}
\end{theorem}

\begin{proof}
The equivalence of (i) and (ii) is the definition.  If (ii) holds, then
\begin{equation}
  0=\Tr(\QzeroF\rho)
  =\|\rho^{1/2}\QzeroF\|_2^2,
\end{equation}
so $\rho^{1/2}\QzeroF=0$.  Hence $\rho\QzeroF=\QzeroF\rho=0$ and
$\rho=\PzeroF\rho\PzeroF$.  Since $\PzeroF$ has rank one and $\Tr\rho=1$,
we obtain $\rho=\PzeroF$.  The converse is immediate.
\end{proof}

\begin{identification}[Rig pa as exact ground realization]
\label{id:rigpa-density}
For pure or mixed phenomenal states, \emph{rig pa} is identified with the
single state $\PzeroF$.  It cannot coexist with a nonzero differentiated
component.  High but nonunit ground fidelity remains a state of \emph{sems}.
\end{identification}

The distinction matches the doctrinal problem identified in the source
tradition: reflexive awareness may already apprehend cognition's nondual
nature while, under the continuing force of the seeds of error, no definitive
ascertainment of that nonduality follows \cite[p.~425]{Yiannopoulos2020}.

\subsection{No second recognition operator is assumed}

\begin{modelchoice}[Realization is first a state condition]
\label{choice:no-recognition-operator}
No additional operator $R$ is introduced whose action is supposed to make
reflexive awareness aware of itself.  Exact realization is first the state
condition
\begin{equation}
  \rho=\PzeroF.
  \label{eq:realization-state-not-second-operator}
\end{equation}
Additional epistemic records, reports, or conventional ascertainments may be
modeled, but none is required to constitute reflexive manifestness.
\end{modelchoice}

This choice prevents the mathematical model from rebuilding the regress that
$\Aodd^2=0$ was introduced to block.  It also keeps distinct:
\begin{equation}
  \boxed{
  \begin{aligned}
  \Aodd\iF=0
  &\quad\text{intrinsic reflexive terminality},\\
  \rho=\PzeroF
  &\quad\text{exact realization of the luminous ground},\\
  \text{a later judgment or report}
  &\quad\text{a conventional cognitive product, if modeled at all}.
  \end{aligned}}
  \label{eq:three-level-reflexivity-realization-report}
\end{equation}
% =============================================================================
% END INLINED FILE: sections/27_reflexive_terminality_recognition.tex
% =============================================================================
% =============================================================================
% BEGIN INLINED FILE: sections/28_ground_nonalteration.tex
% =============================================================================
\section{Ground non-alteration and stabilizer structure}
\label{sec:ground-nonalteration}

The contemplative instruction not to alter or adulterate the natural state
cannot be represented by $\Aodd\ket{\Omega_0}=0$ alone.  Nilpotent terminality
states that Awareness has no further forward leg.  Non-alteration concerns
which operations or dynamics preserve the luminous ground once it is given.

\subsection{Projective preservation by operators}

\begin{definition}[Ground-line preserving operator]
\label{def:ground-line-preserving-operator}
A bounded operator $X\in\mathcal B(\scrHF)$ preserves the luminous ground line
when
\begin{equation}
  X\ket{\Omega_0}\in\CC\ket{\Omega_0},
  \label{eq:ground-line-preserving-vector}
\end{equation}
or equivalently
\begin{equation}
  \QzeroF X\PzeroF=0.
  \label{eq:ground-line-preserving-block}
\end{equation}
It preserves the ground projectively when the scalar in
\eqref{eq:ground-line-preserving-vector} is nonzero.
\end{definition}

\begin{proposition}[One-sided preservation and reduction]
\label{prop:ground-preservation-reduction}
Let $X\in\mathcal B(\scrHF)$.
\begin{enumerate}[label=(\roman*)]
  \item $X$ preserves the ground line if and only if
  $\QzeroF X\PzeroF=0$.
  \item Both $X$ and $X^*$ preserve the ground line if and only if
  \begin{equation}
    [X,\PzeroF]=0.
    \label{eq:ground-reducing-commutator}
  \end{equation}
  In that case the ground line reduces $X$.
\end{enumerate}
\end{proposition}

\begin{proof}
The first claim is the definition of invariance of $\Ran\PzeroF$.  If both
$X$ and $X^*$ preserve it, then the orthogonal complement is also invariant,
so both off-diagonal blocks vanish and $X$ commutes with $\PzeroF$.  The
converse is immediate.
\end{proof}

For a normalized state, multiplication by a nonzero scalar followed by
normalization returns the same ray.  Thus projective non-alteration is weaker
than literal vector fixation but is the invariant notion appropriate to
quantum state space.

\subsection{Unitary non-alteration}

Let $U(t)$ be a strongly continuous unitary group on $\scrHF$.

\begin{theorem}[Equivalent unitary ground-preservation conditions]
\label{thm:unitary-ground-preservation}
The following are equivalent:
\begin{enumerate}[label=(\roman*)]
  \item $U(t)\PzeroF U(t)^*=\PzeroF$ for every $t\in\RR$;
  \item $U(t)\ket{\Omega_0}=e^{-itg_0}\ket{\Omega_0}$ for some
  $g_0\in\RR$ and every $t$;
  \item if $G$ is the self-adjoint generator of $U(t)$, then
  \begin{equation}
    G\ket{\Omega_0}=g_0\ket{\Omega_0}.
    \label{eq:ground-generator-eigenvector}
  \end{equation}
\end{enumerate}
\end{theorem}

\begin{proof}
A one-dimensional invariant subspace of a unitary group carries a strongly
continuous one-dimensional unitary representation of $\RR$, hence a character
$t\mapsto e^{-itg_0}$.  Stone's theorem gives the generator statement.  The
remaining implications are immediate.
\end{proof}

In density-operator language the phase disappears:
\begin{equation}
  U(t)\PzeroF U(t)^*=\PzeroF.
  \label{eq:density-ground-unitary-fixed}
\end{equation}
This is the cleanest mathematical form of leaving the ground unchanged under
closed reversible evolution.

\subsection{Ground-fixing quantum channels}

\begin{definition}[Ground-nonaltering channel]
\label{def:ground-nonaltering-channel}
A completely positive trace-preserving map
\begin{equation}
  \mathcal E:\mathcal T_1(\scrHF)\to\mathcal T_1(\scrHF)
\end{equation}
is ground-nonaltering when
\begin{equation}
  \boxed{\mathcal E(\PzeroF)=\PzeroF.}
  \label{eq:channel-ground-fixed}
\end{equation}
\end{definition}

A family $(\mathcal E_t)_{t\geq0}$ is ground-nonaltering when every member
fixes the ground.  This condition states only that practice cannot drive an
already realized ground into differentiated appearance.  It does not yet say
that non-ground states approach the ground.

\begin{identification}[Do not alter the natural state]
\label{id:nonalteration-doctrinal}
The contemplative instruction not to modify, fabricate, or adulterate the
natural state is represented by invariance of the ground ray or ground density
operator:
\begin{equation}
  U(t)\PzeroF U(t)^*=\PzeroF
  \quad\text{or}\quad
  \mathcal E_t(\PzeroF)=\PzeroF.
  \label{eq:nonalteration-summary}
\end{equation}
It is not represented by treating the zero output of $\Aodd$ as destruction
of the ground.
\end{identification}

The MA thesis describes the terminal contemplative condition as realized
without exertion and without conditioning, after differentiated marks and the
subject--object structure have vanished \cite[pp.~184--185]{Yiannopoulos2012}.
The fixed-ground condition is the minimal dynamical structure compatible with
that description.

\subsection{Logical independence from terminality}

\begin{proposition}[Terminality and non-alteration are distinct]
\label{prop:terminality-nonalteration-distinct}
The identity $\Aodd\iF=0$ neither implies nor is implied by
$\mathcal E(\PzeroF)=\PzeroF$ for an independently supplied state-evolution
channel $\mathcal E$.
\end{proposition}

\begin{proof}
The first identity follows from the grading and holds before any dynamical
channel is chosen.  A channel may preserve or fail to preserve $\PzeroF$
without changing the block definition of $\Aodd$.  Conversely, a channel can
fix $\PzeroF$ on a Hilbert space with no odd Awareness architecture at all.
\end{proof}

Thus the meditation reading requires both structures when both claims are
intended:
\begin{equation}
  \boxed{
  \begin{aligned}
  \Aodd\iF=0
  &\quad\text{nothing further must illuminate awareness},\\
  \mathcal E_t(\PzeroF)=\PzeroF
  &\quad\text{nothing in the practice dynamics adulterates the ground}.
  \end{aligned}}
  \label{eq:terminality-versus-nonalteration-box}
\end{equation}
% =============================================================================
% END INLINED FILE: sections/28_ground_nonalteration.tex
% =============================================================================
% =============================================================================
% BEGIN INLINED FILE: sections/29_deconditioning_dynamics.tex
% =============================================================================
\section{Recognition, deconditioning, and ground-attractive dynamics}
\label{sec:deconditioning-dynamics}

Non-alteration alone leaves every non-ground state where it is.  A model of
contemplative transformation must also specify whether practice changes the
phenomenal state, changes only its conventional ascertainment, or combines
both levels.

\subsection{Recognition and deconditioning are not identical}

The doctrinal sources distinguish the ever-present nondual reflexivity of
awareness from its definitive ascertainment.  All cognitions are luminous,
but luminosity is not ordinarily ascertained; the seeds of error can make the
nondual nature already apprehended by reflexive awareness function as though
it had not been apprehended \cite[pp.~162--163]{Yiannopoulos2012}
\cite[p.~425]{Yiannopoulos2020}.

Accordingly, the realization distinguishes:
\begin{enumerate}[label=(\roman*)]
  \item \emph{structural presence}: every phenomenal state belongs to the
  cyclic luminous representation and satisfies reflexive terminality;
  \item \emph{deconditioning}: the differentiated component of the state is
  reduced;
  \item \emph{exact realization}: the state reaches $\PzeroF$;
  \item \emph{conventional ascertainment}: any later judgment, memory, or
  report concerning that realization.
\end{enumerate}

A recognition-only model can describe a case in which the state is already
$\PzeroF$ but has not been conventionally classified.  It cannot by itself
move a state with $\delta_0^{\Phi}>0$ into the ground.

\subsection{A no-go result for ground-preserving unitaries}

\begin{theorem}[No unitary deconditioning under exact non-alteration]
\label{thm:no-unitary-deconditioning}
Let $U(t)$ be a unitary family satisfying
\begin{equation}
  U(t)\PzeroF U(t)^*=\PzeroF.
\end{equation}
Then for every normal state $\rho$,
\begin{equation}
  F_0^{\Phi}(U(t)\rho U(t)^*)=F_0^{\Phi}(\rho),
  \qquad
  \delta_0^{\Phi}(U(t)\rho U(t)^*)=\delta_0^{\Phi}(\rho).
  \label{eq:unitary-fidelity-invariant}
\end{equation}
Consequently no such unitary dynamics can attract a non-ground state to the
luminous ground.
\end{theorem}

\begin{proof}
Using cyclicity of the trace and invariance of $\PzeroF$,
\begin{equation}
  \Tr\!\left(\PzeroF U(t)\rho U(t)^*\right)
  =\Tr\!\left(U(t)^*\PzeroF U(t)\rho\right)
  =\Tr(\PzeroF\rho).
\end{equation}
\end{proof}

This result does not say that unitary dynamics can never map one selected
vector to the ground.  It says that a single reversible dynamics cannot both
leave the ground unchanged and contract the ground deviation of other states.
Ground-attractive practice therefore requires reduced nonunitary dynamics, a
nonlinear normalized flow, or an enlarged unitary dilation whose reduced
phenomenal dynamics is nonunitary.

\subsection{Ground-attractive channels}

\begin{definition}[Ground-attractive deconditioning family]
\label{def:ground-attractive-family}
A family of completely positive trace-preserving maps
$(\mathcal D_t)_{t\geq0}$ is ground-attractive when
\begin{equation}
  \mathcal D_t(\PzeroF)=\PzeroF
  \label{eq:deconditioning-ground-fixed}
\end{equation}
and
\begin{equation}
  \delta_0^{\Phi}(\mathcal D_t(\rho))
  \longrightarrow0
  \qquad(t\to\infty)
  \label{eq:deconditioning-attraction}
\end{equation}
for every $\rho\in\mathfrak S(\scrHF)$ in the intended basin.
\end{definition}

\begin{theorem}[Canonical amplitude-deconditioning semigroup]
\label{thm:canonical-deconditioning-semigroup}
Fix $\gamma>0$ and set $r_t=e^{-\gamma t}$.  Define, for
$\rho\in\mathcal T_1(\scrHF)$,
\begin{align}
  \mathcal D_t(\rho)
  :={}&\PzeroF\rho\PzeroF
  +r_t^{1/2}
   \left(\PzeroF\rho\QzeroF+\QzeroF\rho\PzeroF\right)
  \notag\\
  &+r_t\QzeroF\rho\QzeroF
  +(1-r_t)\Tr(\QzeroF\rho)\PzeroF.
  \label{eq:canonical-deconditioning-channel}
\end{align}
Then $(\mathcal D_t)_{t\geq0}$ is a strongly continuous completely positive
trace-preserving semigroup with:
\begin{enumerate}[label=(\roman*)]
  \item $\mathcal D_t(\PzeroF)=\PzeroF$;
  \item
  \begin{equation}
    \delta_0^{\Phi}(\mathcal D_t(\rho))
    =e^{-\gamma t}\delta_0^{\Phi}(\rho);
    \label{eq:deconditioning-exact-contraction}
  \end{equation}
  \item $\PzeroF$ is the unique fixed normal state for every $t>0$;
  \item
  \begin{equation}
    \|\mathcal D_t(\rho)-\PzeroF\|_1
    \leq
    2e^{-\gamma t/2}\sqrt{\delta_0^{\Phi}(\rho)}.
    \label{eq:deconditioning-trace-distance-bound}
  \end{equation}
\end{enumerate}
Its generator is
\begin{equation}
  \boxed{
  \mathcal L_0(\rho)
  =\gamma\left(
      \Tr(\QzeroF\rho)\PzeroF
      -\frac12\{\QzeroF,\rho\}
    \right).}
  \label{eq:deconditioning-generator}
\end{equation}
\end{theorem}

\begin{proof}
Let $(e_j)$ be an orthonormal basis of $\QzeroF\scrHF$ and define
\begin{equation}
  K_0(t)=\PzeroF+r_t^{1/2}\QzeroF,
  \qquad
  K_j(t)=(1-r_t)^{1/2}\ket{\Omega_0}\bra{e_j}.
\end{equation}
Then the series
\begin{equation}
  \mathcal D_t(\rho)
  =K_0(t)\rho K_0(t)^*
   +\sum_j K_j(t)\rho K_j(t)^*
\end{equation}
converges in trace norm and gives
\eqref{eq:canonical-deconditioning-channel}.  Moreover,
\begin{equation}
  K_0(t)^*K_0(t)+\sum_jK_j(t)^*K_j(t)
  =\PzeroF+r_t\QzeroF+(1-r_t)\QzeroF=\idop,
\end{equation}
so the map is completely positive and trace preserving.  The block formula
shows $\mathcal D_t\mathcal D_s=\mathcal D_{t+s}$ and strong continuity.
Taking the trace against $\QzeroF$ proves
\eqref{eq:deconditioning-exact-contraction}.  If a density operator is fixed
for $t>0$, its ground deviation vanishes and
\cref{thm:exact-rigpa-density} gives $\rho=\PzeroF$.  Finally, the
Fuchs--van de Graaf bound for fidelity with a pure state gives
\eqref{eq:deconditioning-trace-distance-bound}.  Differentiating at $t=0$
gives \eqref{eq:deconditioning-generator}.
\end{proof}

The semigroup is a standard pure-ground amplitude-damping construction in the
language of completely positive dynamics
\cite{GKS1976,Lindblad1976,FuchsVandeGraaf1999}.  Its novelty here is only its
realization role.

\begin{modelchoice}[Canonical deconditioning model]
\label{choice:canonical-deconditioning-model}
The family \eqref{eq:canonical-deconditioning-channel} is retained as a
minimal explicit model of v\=asan\=a exhaustion.  It is not asserted to be the
unique or biologically literal dynamics of contemplative practice.
\end{modelchoice}

\subsection{Ground deviation as a Lyapunov functional}

For the canonical semigroup,
\begin{equation}
  \frac{d}{dt}\delta_0^{\Phi}(\mathcal D_t(\rho))
  =-\gamma\delta_0^{\Phi}(\mathcal D_t(\rho)).
  \label{eq:ground-deviation-lyapunov}
\end{equation}
Thus $\delta_0^{\Phi}$ is a strict Lyapunov functional off the ground.  The
luminous ground is both fixed and globally attractive, while exact realization
at any finite time occurs only for states already at the ground unless a
finite-time reset or projection-like channel is separately chosen.

\begin{remark}[Reduced irreversibility]
\label{rem:reduced-irreversibility}
A completely positive contraction need not be fundamental nonunitarity of the
universal carrier.  It may arise as the reduced dynamics of a larger unitary
system after environmental or unrecorded degrees of freedom are discarded.
The reduced model does not specify which degrees of freedom provide such a
dilation.
\end{remark}
% =============================================================================
% END INLINED FILE: sections/29_deconditioning_dynamics.tex
% =============================================================================
% =============================================================================
% BEGIN INLINED FILE: sections/30_momentary_stabilization.tex
% =============================================================================
\section{Momentary cognitive continua and stabilization}
\label{sec:momentary-stabilization}

A Buddhist cognitive continuum is not naturally represented as one enduring
substance carrying numerically identical states through time.  The minimum
formal requirement is a sequence of momentary states connected by causal
channels.

\subsection{A channel system of moments}

\begin{definition}[Momentary phenomenal realization system]
\label{def:momentary-phenomenal-system}
A momentary phenomenal realization system consists of:
\begin{equation}
  \left(
    \mathscr H_{\Phi,n},P_{0,n}^{\Phi},\rho_n,\mathcal T_n
  \right)_{n\geq0},
  \label{eq:momentary-system-data}
\end{equation}
where:
\begin{enumerate}[label=(\roman*)]
  \item $\mathscr H_{\Phi,n}$ is the phenomenal Hilbert space at moment $n$;
  \item $P_{0,n}^{\Phi}$ is a rank-one luminous ground projection on $\mathscr H_{\Phi,n}$;
  \item $\rho_n\in\mathfrak S(\mathscr H_{\Phi,n})$;
  \item
  \begin{equation}
    \mathcal T_n:\mathcal T_1(\mathscr H_{\Phi,n})
    \longrightarrow\mathcal T_1(\mathscr H_{\Phi,n+1})
  \end{equation}
  is completely positive and trace preserving, with
  $\rho_{n+1}=\mathcal T_n(\rho_n)$.
\end{enumerate}
\end{definition}

The spaces may all be identified with one fixed representation, but the
fibered notation avoids building temporal numerical identity into the model.
Define
\begin{equation}
  \delta_n
  :=1-\Tr(P_{0,n}^{\Phi}\rho_n).
  \label{eq:momentary-ground-deviation}
\end{equation}

\begin{definition}[Ground covariance and contraction]
\label{def:momentary-ground-covariance}
The system is ground-covariant when
\begin{equation}
  \mathcal T_n(P_{0,n}^{\Phi})=P_{0,n+1}^{\Phi}.
  \label{eq:momentary-ground-covariance}
\end{equation}
It is contractive with factors $q_n\in[0,1]$ when
\begin{equation}
  \delta_{n+1}\leq q_n\delta_n.
  \label{eq:momentary-ground-contraction}
\end{equation}
\end{definition}

\begin{proposition}[Exact realization is absorbing]
\label{prop:exact-realization-absorbing}
In a ground-covariant system, if $\rho_n=P_{0,n}^{\Phi}$ at some moment $n$, then
\begin{equation}
  \rho_{n+1}=P_{0,n+1}^{\Phi}.
  \label{eq:rigpa-absorbing-next-moment}
\end{equation}
If ground covariance continues, exact realization persists at all later
moments.
\end{proposition}

\begin{proof}
Apply \eqref{eq:momentary-ground-covariance} to the equality
$\rho_n=P_{0,n}^{\Phi}$ and iterate.
\end{proof}

This is the momentary version of non-alteration: the ground need not be one
enduring vector, but its canonical state type is carried into the next moment.

\subsection{Stabilization theorem}

\begin{theorem}[Product contraction implies ground stabilization]
\label{thm:momentary-ground-stabilization}
Suppose the momentary system is ground-covariant and satisfies
\eqref{eq:momentary-ground-contraction}.  Then
\begin{equation}
  \delta_n
  \leq
  \left(\prod_{k=0}^{n-1}q_k\right)\delta_0.
  \label{eq:momentary-product-bound}
\end{equation}
If
\begin{equation}
  \prod_{k=0}^{\infty}q_k=0,
  \label{eq:momentary-product-zero}
\end{equation}
then
\begin{equation}
  \Tr(P_{0,n}^{\Phi}\rho_n)\longrightarrow1
  \label{eq:momentary-fidelity-convergence}
\end{equation}
and
\begin{equation}
  \|\rho_n-P_{0,n}^{\Phi}\|_1
  \leq2\sqrt{\delta_n}
  \longrightarrow0.
  \label{eq:momentary-trace-convergence}
\end{equation}
\end{theorem}

\begin{proof}
Iteration gives \eqref{eq:momentary-product-bound}.  The product hypothesis
implies $\delta_n\to0$.  The trace-norm estimate is the pure-state
Fuchs--van de Graaf inequality in each fiber.
\end{proof}

The theorem allows intervals with $q_n=1$, during which no progress occurs,
provided the accumulated contraction over the continuum is sufficient.  It
also permits nonuniform practice and moment-dependent conditioning.

\subsection{Recognition, event, and stabilization}

\begin{definition}[Realization event and stabilization]
\label{def:realization-event-stabilization}
A realization event occurs at moment $n$ when
\begin{equation}
  \rho_n=P_{0,n}^{\Phi}.
\end{equation}
A continuum is asymptotically stabilized in realization when
\begin{equation}
  \|\rho_n-P_{0,n}^{\Phi}\|_1\longrightarrow0.
\end{equation}
\end{definition}

A single exact event and stable realization are therefore distinct notions.
Ground covariance makes an exact event absorbing; product contraction gives a
sufficient condition for asymptotic stabilization without positing an
enduring self.

\begin{identification}[Continuity without substance]
\label{id:continuity-without-substance}
The sequence of ground projectors $P_{0,n}^{\Phi}$ represents one canonical
luminous state type instantiated moment by moment.  The causal continuity of
reflexive awareness belongs to the channel system, not to a numerically
identical substantial subject persisting through time.
\end{identification}

This is compatible with the dissertation's report that the Yog\=ac\=ara model
treats reflexive continuity as unbroken even where ordinary differentiated
content or a subsequent conceptual ascertainment is unavailable, including
states such as deep sleep and coma \cite[p.~425]{Yiannopoulos2020}.  Variation
across distinct sentient continua requires family-level realization data
beyond the present construction, which addresses only temporal succession
within one continuum.
% =============================================================================
% END INLINED FILE: sections/30_momentary_stabilization.tex
% =============================================================================
% =============================================================================
% BEGIN INLINED FILE: sections/31_third_fourth_yoga_same_ground.tex
% =============================================================================
\section{The Third and Fourth Yogas as one ground realization}
\label{sec:third-fourth-yoga-same-ground}

The Third and Fourth Yogas share one ground ray.  The algebraic distinction
between luminosity-only and nonappearance is supplied without introducing a
second ket or a state below the ground.

\subsection{The exhausted profile representation is one-dimensional}

Let
\begin{equation}
  \detAlgEmpty
  :=\overline{\wordSet(\vasEmpty)}^{\|\cdot\|}
  \subseteq\detAlg
  \label{eq:exhausted-profile-Cstar-algebra}
\end{equation}
be the $C^*$-algebra generated by the operations active under the exhausted
profile.

\begin{theorem}[Ground character of the exhausted profile]
\label{thm:exhausted-profile-character}
Under \cref{choice:exhausted-profile-ground-triviality}, there exists a
character
\begin{equation}
  \chi_0^{\Phi}:\detAlgEmpty\to\CC
  \label{eq:exhausted-ground-character}
\end{equation}
such that
\begin{equation}
  \piF(X)\ket{\Omega_0}
  =\chi_0^{\Phi}(X)\ket{\Omega_0}
  \qquad
  (X\in\detAlgEmpty).
  \label{eq:exhausted-character-action}
\end{equation}
Moreover,
\begin{equation}
  \chi_0^{\Phi}(X)=\omegaF(X),
  \label{eq:exhausted-character-state}
\end{equation}
and every centered exhausted-profile determination acts trivially on the
ground:
\begin{equation}
  \omegaF(X)=0
  \quad\Longrightarrow\quad
  \piF(X)\ket{\Omega_0}=0.
  \label{eq:centered-exhausted-annihilates-ground}
\end{equation}
The cyclic representation generated from $\ket{\Omega_0}$ is therefore
one-dimensional.
\end{theorem}

\begin{proof}
By \cref{prop:exhausted-words-ground-preserve}, every word acts by a scalar on
the ground.  Continuity extends this to $\detAlgEmpty$.  If
$X\Omega_0=\chi(X)\Omega_0$ and $Y\Omega_0=\chi(Y)\Omega_0$, then
$XY\Omega_0=\chi(X)\chi(Y)\Omega_0$; adjoints give the $*$-property and the
identity gives unitality.  Taking the inner product with $\Omega_0$ proves
\eqref{eq:exhausted-character-state}.  The centered claim and
one-dimensionality follow.
\end{proof}

This is a precise sense in which appearance has ceased while the state has
not become zero.  The active determination algebra retains only scalar action
on the luminous ground; all centered appearance-generating action is absent.

\subsection{One state, two praxeological descriptions}

\begin{identification}[Third Yoga]
\label{id:third-yoga}
The Third Yoga is represented by exact realization of the luminous ground:
\begin{equation}
  \rho=\PzeroF.
  \label{eq:third-yoga-ground-state}
\end{equation}
The emphasis falls on undifferentiated luminosity as the sole remaining
phenomenal condition.
\end{identification}

\begin{identification}[Fourth Yoga]
\label{id:fourth-yoga}
The Fourth Yoga is represented by the same state
\begin{equation}
  \rho=\PzeroF,
  \label{eq:fourth-yoga-same-state}
\end{equation}
together with the exhausted-profile fact that the active centered appearance
algebra acts trivially on that ground.  The emphasis falls on nonappearance
and the nonobjectification of luminosity as one further phenomenal item.
\end{identification}

Thus
\begin{equation}
  \boxed{
  \rho_{\mathrm{III}}
  =\rho_{\mathrm{IV}}
  =\PzeroF.}
  \label{eq:third-fourth-identical-density}
\end{equation}
The distinction is praxeological and representational, not a spectral passage
to $\ket{\Omega_{-1}}$.

The MA thesis explicitly describes the Third Yoga as the disappearance of all
differentiated appearance into nondual luminosity, and the Fourth Yoga as
nonappearance even of luminosity insofar as it is treated as an appearance
\cite[pp.~182--185]{Yiannopoulos2012}.  It also insists that the Fourth Yoga
is not a blank void.  The one-state construction preserves all three points.

\subsection{The precise meaning of ``contentless''}

The ground is contentless only in the restricted sense
\begin{equation}
  \QzeroF\rho\QzeroF=0,
  \label{eq:no-differentiated-content-ground}
\end{equation}
that is, no differentiated \=ak\=ara, grasper--grasped frame, mnemonic
concatenation, or conceptual determination is present.  It is not the zero
operator, the zero vector, or the absence of phenomenality.

\begin{proposition}[No phenomenal spectral state below the ground]
\label{prop:no-state-below-ground}
If $H_{\Phi}$ is self-adjoint and bounded below with
$E_0^{\Phi}=\inf\spec H_{\Phi}$, then
\begin{equation}
  \spec(H_{\Phi})\cap(-\infty,E_0^{\Phi})=\varnothing.
  \label{eq:no-spectrum-below-ground}
\end{equation}
Equivalently, its closed quadratic form satisfies
\begin{equation}
  q_{H_{\Phi}}[\phi]
  \geq E_0^{\Phi}\|\phi\|^2
  \qquad
  (\phi\in\Dom q_{H_{\Phi}}).
  \label{eq:ground-form-lower-bound}
\end{equation}
If the ground projection is uniquely $\PzeroF$, that lower endpoint selects
one phenomenal ground ray.  The Fourth Yoga therefore cannot be represented
by a lower spectral label $\ket{\Omega_{-1}}$.
\end{proposition}

\begin{proof}
The spectral statement is the definition of the spectral infimum, and the
form inequality follows from the spectral theorem \cite{Kato,ReedSimon}.
Ground-ray uniqueness is additional to lower boundedness and identifies the
one-dimensional eigenspace at $E_0^{\Phi}$.
\end{proof}

\subsection{Recognition without fabrication}

The exhausted character and ground-fixing dynamics together give the formal
structure suggested by ``without exertion and without conditioning'':
\begin{equation}
  \boxed{
  \begin{aligned}
  \mathcal E_t(\PzeroF)&=\PzeroF,
  \\
  \piF(X^{\circ})\ket{\Omega_0}&=0
  \quad(X\in\detAlgEmpty).
  \end{aligned}}
  \label{eq:unconditioned-ground-two-conditions}
\end{equation}
The first line says the realized ground is not altered.  The second says the
active exhausted profile produces no differentiated appearance from it.  No
additional operation is required to make the ground reflexively manifest.
% =============================================================================
% END INLINED FILE: sections/31_third_fourth_yoga_same_ground.tex
% =============================================================================
% =============================================================================
% BEGIN INLINED FILE: sections/33_volition_return_orientation.tex
% =============================================================================
\section{Volition as adjoint return and normalized orientation}
\label{sec:volition-return-orientation}

The physical-to-phenomenal orientation has a canonical reverse direction.  The
central point is that Volition is neither a second independent substance nor
an arbitrary force added after Awareness.  At the typed level it is the
Hilbert adjoint
\begin{equation}
  \boxed{
  V:=A^*:\Dom A^*\subseteq\scrHF\longrightarrow\scrHP,}
  \label{eq:typed-volition}
\end{equation}
while at the normalized support level it is the reverse polar orientation
\begin{equation}
  \boxed{
  \restr{\polA^*}{\actHF}:\actHF\longrightarrow\actHP.}
  \label{eq:polar-volition}
\end{equation}
These two maps have related but different roles.  The raw adjoint $A^*$
carries transition magnitude and may be unbounded; the partial isometry
$u_A^*$ is bounded on all of $\scrHF$ and carries only the reverse orientation.

\subsection{Polar decomposition of the adjoint}

For the closed densely defined channel
\begin{equation}
  A=\polA\absA,
\end{equation}
the adjoint has polar form
\begin{equation}
  \boxed{A^*=\absA\polA^*}
  \label{eq:adjoint-polar-form}
\end{equation}
on its natural domain \cite{Kato,ReedSimon}.  Thus a raw Volition return consists of two
mathematical stages:
\begin{equation}
  \phi
  \xmapsto{\ \polA^*\ }
  \polA^*\phi
  \xmapsto{\ \absA\ }
  \absA\polA^*\phi.
  \label{eq:volition-two-stage}
\end{equation}
The first stage identifies the corresponding physical direction.  The second
weights that direction by the positive transition magnitude.

\begin{proposition}[Normalized Volition is the inverse orientation]
\label{prop:normalized-volition-inverse}
On the active supports,
\begin{equation}
  \polA^*\polA=I_{\actHP},
  \qquad
  \polA\polA^*=I_{\actHF}.
  \label{eq:active-inverse-orientations}
\end{equation}
Consequently,
\begin{equation}
  \polA^*(\polA\psi)=\psi
  \quad(\psi\in\actHP),
  \qquad
  \polA(\polA^*\phi)=\phi
  \quad(\phi\in\actHF).
  \label{eq:active-round-trips}
\end{equation}
\end{proposition}

\begin{proof}
This is the unitary equivalence established by polar decomposition between
$(\Ker A)^\perp$ and $\overline{\Ran A}$.
\end{proof}

\statuslabel{Imported theorem} The inverse relation belongs to the polar
orientations on active support.  It does not imply $A^*A=I$ or $AA^*=I$ for
the raw channel.

\subsection{Restricted ground return versus full polar return}

The canonical rank-one ground orientation is
\begin{equation}
  \uzero=\ket{\Omega_0}\!\bra0,
  \qquad
  \uzero^*=\ket0\!\bra{\Omega_0}.
  \label{eq:ground-volition-orientation}
\end{equation}
Therefore the luminous ground always has the exact restricted return
\begin{equation}
  \boxed{\uzero^*\ket{\Omega_0}=\ket0.}
  \label{eq:restricted-ground-return}
\end{equation}
This follows from the weak ground bridge alone.

The full polar return is subtler.  Since
\begin{equation}
  A\ket0=a_0\ket{\Omega_0}=\polA\absA\ket0,
\end{equation}
application of $\polA^*$ gives
\begin{equation}
  \boxed{
  \polA^*\ket{\Omega_0}
  =\frac{\absA\ket0}{a_0}.}
  \label{eq:full-polar-ground-preimage}
\end{equation}
The vector on the right has norm one.  Define
\begin{equation}
  \ket{0_A}:=\polA^*\ket{\Omega_0}
  =a_0^{-1}\absA\ket0.
  \label{eq:polar-ground-preimage-def}
\end{equation}

\begin{proposition}[Ground-return compatibility criterion]
\label{prop:ground-return-compatibility}
The following are equivalent:
\begin{enumerate}[label=(\roman*)]
  \item $\polA^*\ket{\Omega_0}=\ket0$;
  \item $\polA\ket0=\ket{\Omega_0}$;
  \item $\absA\ket0=a_0\ket0$;
  \item $A^*A\ket0=a_0^2\ket0$.
\end{enumerate}
Thus exact agreement between restricted and full normalized ground return is
precisely the strong polar ground condition of
\cref{choice:strong-polar-ground,thm:strong-ground-equivalences}.
\end{proposition}

\begin{proof}
The equivalence of (iii) and (iv) follows by positive functional calculus.
Equation \eqref{eq:full-polar-ground-preimage} gives (i)$\Leftrightarrow$(iii),
and the unitary inverse relation on active support gives
(i)$\Leftrightarrow$(ii).
\end{proof}

\begin{modelchoice}[Full-channel ground reciprocity]
\label{choice:full-ground-reciprocity}
One may impose the equivalent conditions of
\cref{prop:ground-return-compatibility}.  This condition is not required for
the existence of the exact ground bridge, because the restricted return
$\uzero^*$ is already canonical.  Full-channel reciprocity is required only
if the complete polar orientation is to identify the physical vacuum and
luminous ground directly.
\end{modelchoice}

\subsection{Raw ground return}

The weak bridge does not guarantee that $\ket{\Omega_0}\in\Dom A^*$.  The
bounded map $\uzero^*$ is therefore more primitive than the raw expression
$A^*\ket{\Omega_0}$.  When the latter exists,
\begin{equation}
  A^*\ket{\Omega_0}
  =\absA\ket{0_A}.
  \label{eq:raw-ground-return-general}
\end{equation}
Under full-channel ground reciprocity this reduces to
\begin{equation}
  A^*\ket{\Omega_0}=a_0\ket0.
  \label{eq:raw-ground-return-strong}
\end{equation}

\warninglabel An adjoint is not automatically an everywhere-defined state
map.  Projective return is supplied globally on the representable support by
$u_A^*$; the raw return $A^*$ remains subject to its operator domain.

\subsection{No additional agentive functioning}

The term \emph{mental causation} can misleadingly suggest that a substantial
mind performs an intermediate action upon matter.  The present architecture
requires no such operation.  Volition is the reverse orientation already
forced by the adjoint structure of the one transition channel.  In this
sense, phenomenal-to-physical return is structural rather than an additional
agentive \emph{vy\=ap\=ara} inserted between a cognition and its result.  This is
compatible with the Dharmak\=irtian analysis according to which cognition is
an arising from causes rather than a substantial agent transitively
apprehending or manipulating an object.

\begin{identification}[Volition as reverse orientation]
\label{id:volition-reverse-orientation}
The typed adjoint $V=A^*$ and normalized orientation $u_A^*$ are interpreted
as the two levels of phenomenal-to-physical return:
\begin{equation}
  \boxed{
  \begin{aligned}
  u_A^*&:\ \text{reverse orientation and representability},\\
  A^*&:\ \text{reverse orientation together with channel magnitude}.
  \end{aligned}}
  \label{eq:volition-two-level-reading}
\end{equation}
Neither is a second ontological primitive independent of Awareness.
\end{identification}
% =============================================================================
% END INLINED FILE: sections/33_volition_return_orientation.tex
% =============================================================================
% =============================================================================
% BEGIN INLINED FILE: sections/34_asymmetric_presentability.tex
% =============================================================================
\section{Asymmetric presentability and the one-sided inverse law}
\label{sec:asymmetric-presentability}

The polar architecture imposes two simultaneous claims: every physical state
admits a phenomenal presentation, while the phenomenal domain contains
possibilities with no physical counterpart.  These claims require a one-sided,
not a two-sided, inverse law.

\subsection{Universal phenomenal presentability of physical states}

\begin{axiom}[Physical presentability]
\label{ax:physical-presentability}
No nonzero physical direction in the domain of the typed Awareness channel is
annihilated:
\begin{equation}
  \boxed{\Ker A=\{0\}.}
  \label{eq:physical-presentability}
\end{equation}
Equivalently,
\begin{equation}
  \actHP=\scrHP,
  \qquad
  \suppProjP=\polA^*\polA=I_{\scrHP}.
  \label{eq:physical-presentability-polar}
\end{equation}
\end{axiom}

\statuslabel{Realization axiom} This is the formal claim that every physical
ray has a normalized phenomenal orientation under $u_A$.  It does not require
the raw unbounded channel $A$ to act on every vector; $u_A$ extends as a
bounded isometry on the whole physical Hilbert space.

\subsection{The final support and its phenomenal defect}

The final support projection and its orthogonal complement are
\begin{equation}
  \boxed{
  \suppProjF:=\polA\polA^*,
  \qquad
  \defectProjF:=I_{\scrHF}-\suppProjF.}
  \label{eq:rep-excess-projections}
\end{equation}
Their ranges are
\begin{equation}
  \Ran\suppProjF=\overline{\Ran A}=\actHF,
  \qquad
  \Ran\defectProjF=(\actHF)^{\perp}=\Ker A^*.
  \label{eq:rep-excess-ranges}
\end{equation}
The second space is the standard final-defect space of the polar orientation.
Its cognitive interpretation as \emph{phenomenal excess} is additional to the
operator-theoretic fact.

\begin{axiom}[Strict phenomenal excess]
\label{ax:strict-phenomenal-excess}
The phenomenal final support has a nontrivial orthogonal defect:
\begin{equation}
  \boxed{
  \defectHF
  =\Ker A^*
  \neq\{0\}.}
  \label{eq:strict-phenomenal-excess}
\end{equation}
Equivalently,
\begin{equation}
  \suppProjF<I_{\scrHF}.
\end{equation}
\end{axiom}

Here ``strictly larger'' means that the final phenomenal support is a proper
closed subspace.  In infinite-dimensional separable Hilbert spaces it does
not necessarily mean a larger cardinal Hilbert dimension: a Hilbert space may
be unitarily isomorphic to one of its proper closed subspaces.  The invariant
content is the nontrivial defect space $\Ker A^*$.

\subsection{The one-sided inverse}

\begin{theorem}[Asymmetric normalized inverse law]
\label{thm:asymmetric-inverse-law}
Under \cref{ax:physical-presentability,ax:strict-phenomenal-excess},
\begin{equation}
  \boxed{
  \polA^*\polA=I_{\scrHP},
  \qquad
  \polA\polA^*=\suppProjF<I_{\scrHF}.}
  \label{eq:mature-one-sided-inverse}
\end{equation}
Thus $u_A^*$ is a left inverse of $u_A$, while the reverse ordered product is
the projection onto the final phenomenal support.
\end{theorem}

\begin{proof}
Physical presentability makes the initial support of $u_A$ equal to all of
$\scrHP$.  Polar decomposition gives the final support
$\overline{\Ran A}$ and the associated range projection $\suppProjF$.  Strict
phenomenal excess makes this projection proper.
\end{proof}

In normalized orientation form,
\begin{equation}
  \boxed{
  V_{\mathrm{or}}A_{\mathrm{or}}=I_{\Psi},
  \qquad
  A_{\mathrm{or}}V_{\mathrm{or}}=\suppProjF,}
  \label{eq:orientation-AV-law}
\end{equation}
with
\begin{equation}
  A_{\mathrm{or}}=u_A,
  \qquad
  V_{\mathrm{or}}=u_A^*.
\end{equation}

\begin{identification}[Asymmetric nonduality]
\label{id:asymmetric-nonduality}
Equation \eqref{eq:mature-one-sided-inverse} is interpreted as follows:
\begin{equation}
  \boxed{
  \begin{aligned}
  &\text{every physical state admits a norm-preserving phenomenal presentation,}\\
  &\text{but not every phenomenal state admits a normalized polar return.}
  \end{aligned}}
  \label{eq:asymmetric-nonduality-reading}
\end{equation}
The asymmetry is internal to one transition architecture and does not
reintroduce two independently self-sufficient substances.
\end{identification}

\begin{remark}[Polar antecedents and raw antecedents]
\label{rem:polar-vs-raw-antecedent}
Every vector in $\actHF=\overline{\Ran A}$ has a unique normalized polar
antecedent under $u_A^*$.  If $\Ran A$ is not closed, however, a vector in
$\actHF\setminus\Ran A$ need not possess a literal raw antecedent
$\psi\in\Dom A$ satisfying $A\psi=\phi$.  The phrases ``presentation'' and
``return'' refer to the supportwise polar correspondence unless raw range
membership is stated explicitly.  In the Fredholm regime the range is closed,
so the distinction disappears.
\end{remark}

\subsection{Final-support and defect components}

Every phenomenal vector has the unique orthogonal decomposition
\begin{equation}
  \boxed{
  \phi=\phi_A+\phi_A^{\perp},
  \qquad
  \phi_A:=\suppProjF\phi,
  \qquad
  \phi_A^{\perp}:=\defectProjF\phi.}
  \label{eq:phenomenal-rep-exc-decomp}
\end{equation}
The normalized return sees only the final-support component:
\begin{equation}
  \boxed{
  \polA^*\phi
  =\polA^*\phi_A,
  \qquad
  \polA^*\phi_A^{\perp}=0.}
  \label{eq:return-sees-representable}
\end{equation}
Re-presentation after return yields
\begin{equation}
  \polA\polA^*\phi=\suppProjF\phi=\phi_A.
  \label{eq:phen-return-cycle-projection}
\end{equation}
Thus the normalized phenomenal--physical--phenomenal cycle is not identity on
the whole phenomenal space.  It is the projection onto the final support of
the channel.

\subsection{Canonical example}

The unilateral-shift realization in \cref{sec:fredholm-index-excess} shows
that the asymmetry is possible even when the sector Hilbert spaces are
abstractly isomorphic.
% =============================================================================
% END INLINED FILE: sections/34_asymmetric_presentability.tex
% =============================================================================
% =============================================================================
% BEGIN INLINED FILE: sections/35_mental_causation_actualization.tex
% =============================================================================
\section{Mental causation as phenomenal-to-physical actualization}
\label{sec:mental-causation-actualization}

The phrase \emph{mental causation} is used here in a narrow structural sense:
a phenomenal state or determination can contribute to physical actuality
only through the Volition return channel.  This does not mean that every
phenomenal vector is automatically a physically realized state.  The
one-sided inverse law makes actualization partial on phenomenal projective
space.

\subsection{Presentable support weight}

For a normalized phenomenal vector $\phi\in\scrHF$, define its final-support
weight by
\begin{equation}
  \boxed{
  \presentWeight(\phi)
  :=\|\polA^*\phi\|^2
  =\inner{\phi}{\suppProjF\phi}_{\Phi}.}
  \label{eq:pure-actualization-weight}
\end{equation}
Then
\begin{equation}
  0\leq\presentWeight(\phi)\leq1,
\end{equation}
with
\begin{equation}
  \presentWeight(\phi)=1
  \iff \phi\in\actHF,
  \qquad
  \presentWeight(\phi)=0
  \iff \phi\in\defectHF.
  \label{eq:actualization-weight-extremes}
\end{equation}

The notation $w_A$ is deliberately neutral.  It is a structural support
weight.  It becomes a success probability only if the return operation is
embedded as a branch of a separately supplied quantum instrument.

\begin{definition}[Conditional projective actualization]
\label{def:conditional-projective-actualization}
For $\presentWeight(\phi)>0$, define
\begin{equation}
  \boxed{
  \mathfrak v_A([\phi])
  :=\left[
  \frac{\polA^*\phi}{\|\polA^*\phi\|}
  \right]
  \in\mathbb P(\scrHP).}
  \label{eq:projective-actualization-map}
\end{equation}
This partial projective map is undefined on
$\mathbb P(\defectHF)$ because Volition returns those vectors to the zero
vector, which has no projective ray.
\end{definition}

The phenomenal state recovered after actualization and re-presentation is
\begin{equation}
  \frac{\suppProjF\phi}{\|\suppProjF\phi\|}.
  \label{eq:re-presented-actualized-state}
\end{equation}
Thus actualization does not select an arbitrary physical state.  It returns
the physical correlate of the state's final-support component.

\subsection{A trace-nonincreasing quantum operation}

For a normal phenomenal state $\rho$, define the subnormalized return
operation
\begin{equation}
  \boxed{
  \returnOp(\rho):=\polA^*\rho\polA.}
  \label{eq:subnormalized-return-operation}
\end{equation}
This is a completely positive map from trace-class operators on $\scrHF$ to
trace-class operators on $\scrHP$.

\begin{proposition}[Trace and conditional return]
\label{prop:return-operation-trace}
For every normal phenomenal state $\rho$,
\begin{equation}
  \Tr\returnOp(\rho)
  =\Tr(\suppProjF\rho)
  =:\presentWeight(\rho).
  \label{eq:mixed-actualization-weight}
\end{equation}
Hence $\returnOp$ is trace-nonincreasing.  If $\presentWeight(\rho)>0$, the
conditional physical state is
\begin{equation}
  \boxed{
  \returnOpCond(\rho)
  :=\frac{\polA^*\rho\polA}
  {\Tr(\suppProjF\rho)}.}
  \label{eq:conditional-return-state}
\end{equation}
\end{proposition}

\begin{proof}
Cyclicity of the trace gives
\begin{equation}
  \Tr(\polA^*\rho\polA)
  =\Tr(\rho\polA\polA^*)
  =\Tr(\rho\suppProjF).
\end{equation}
Positivity and complete positivity follow from the one-Kraus form.
\end{proof}

\statuslabel{Derived proposition} The quantity
$\presentWeight(\rho)$ is a structural support weight.  If
$\returnOp$ is implemented as one branch of a quantum instrument, it becomes
the success probability of that branch.  The SSP architecture does not by
itself require such an operational measurement interpretation.

\subsection{Raw Volition and spectral reweighting}

When $A$ is bounded, the raw return operation on vectors is
\begin{equation}
  V\phi=A^*\phi=\absA\polA^*\phi.
  \label{eq:raw-volition-vector-return}
\end{equation}
For unbounded $A$, this remains valid on $\Dom A^*$.  The polar return
$u_A^*$ supplies the canonical supportwise correspondence.  The positive
factor $|A|$ then supplies mode-dependent spectral reweighting.  Unless the
returned vector lies in a single singular subspace of $|A|$, that reweighting
may change its projective ray rather than merely multiply one fixed direction
by a scalar.  The normalized actualization map therefore isolates the
projective correspondence from, but does not erase, the raw channel dynamics.

\subsection{Volition is not arbitrary selection}

Volition selects a physical correlate from phenomenal possibility only in the
following channel-constrained sense:
\begin{equation}
  \phi
  \longmapsto
  \suppProjF\phi
  \longmapsto
  \polA^*\phi.
  \label{eq:constrained-selection}
\end{equation}
The first arrow removes the component with no physical partner.  The second
identifies the corresponding normalized physical direction.  There is no
freedom to map one phenomenal state to an unrelated physical ray unless
additional control operators are supplied.

\begin{identification}[Mental causation]
\label{id:mental-causation}
Mental causation is the phenomenal-to-physical orientation of the common
transition channel:
\begin{equation}
  \boxed{
  \begin{aligned}
  \text{phenomenal possibility}
  &\xrightarrow{\ \suppProjF\ }
  \text{final-support phenomenal content}\\[-1mm]
  &\xrightarrow{\ \polA^*\ }
  \text{physical actuality}.
  \end{aligned}}
  \label{eq:mental-causation-slogan}
\end{equation}
The process is not a second force external to physical law.  It is the return
orientation by which phenomenal determination is translated into the
physical sector.
\end{identification}

\subsection{Ground actualization}

Because the ground bridge gives
\begin{equation}
  \ket{\Omega_0}\in\Ran A\subseteq\actHF,
\end{equation}
one has
\begin{equation}
  \presentWeight(\Omega_0)=1.
  \label{eq:ground-fully-actualizable}
\end{equation}
Thus the luminous ground is never part of the phenomenal defect space.

The canonical restricted ground operation is
\begin{equation}
  \boxed{
  \returnGroundOp(\rho):=\uzero^*\rho\uzero.}
  \label{eq:restricted-ground-return-operation}
\end{equation}
Since $\uzero=\ket{\Omega_0}\!\bra0$,
\begin{equation}
  \returnGroundOp(\rho)
  =\Tr(\PzeroF\rho)\,\PzeroP.
  \label{eq:restricted-ground-return-formula}
\end{equation}
Whenever $\Tr(\PzeroF\rho)>0$, conditional normalization therefore returns
exactly $\PzeroP$.

At the vector level,
\begin{equation}
  \uzero^*\ket{\Omega_0}=\ket0.
\end{equation}
Under full polar ground reciprocity, the global conditional map also returns
the luminous ground to the vacuum ray:
\begin{equation}
  \mathfrak v_A([\ket{\Omega_0}])=[\ket0].
  \label{eq:ground-global-actualization}
\end{equation}
Without that stronger compatibility, the global return yields the full polar
preimage $\ket{0_A}$ rather than necessarily $\ket0$.
% =============================================================================
% END INLINED FILE: sections/35_mental_causation_actualization.tex
% =============================================================================
% =============================================================================
% BEGIN INLINED FILE: sections/36_phenomenal_excess_deconditioning.tex
% =============================================================================
\section{Phenomenal excess, differentiated appearance, and deconditioning}
\label{sec:phenomenal-excess-deconditioning}

Phenomenal excess must be distinguished from differentiated phenomenal
appearance generally.  A state can be far from the luminous ground while
remaining completely inside the final support of Awareness.  Conversely,
every final-defect direction is necessarily ground-orthogonal.

\subsection{The phenomenal defect is contained in differentiated appearance}

The weak ground bridge gives
\begin{equation}
  \ket{\Omega_0}=a_0^{-1}A\ket0\in\Ran A.
  \label{eq:ground-in-representable-range}
\end{equation}
Hence
\begin{equation}
  \defectHF=\Ker A^*\perp\ket{\Omega_0}.
  \label{eq:excess-ground-orthogonal}
\end{equation}
The orthogonal range--adjoint-kernel decomposition used here is standard for
dense closed operators; see \cite{Kato,ReedSimon}.

\begin{proposition}[Three-part phenomenal decomposition]
\label{prop:three-part-phenomenal-decomposition}
Under the ground bridge,
\begin{equation}
  \boxed{
  \scrHF
  =\GzeroF
  \oplus
  \actAppF
  \oplus
  \defectHF,}
  \label{eq:three-part-phenomenal-decomp}
\end{equation}
where
\begin{equation}
  \actAppF:=\QzeroF\actHF,
  \qquad
  \defectHF=\Ker A^*.
  \label{eq:rep-excess-appearance-spaces}
\end{equation}
Moreover,
\begin{equation}
  \QzeroF\scrHF
  =\actAppF\oplus\defectHF.
  \label{eq:differentiated-split}
\end{equation}
\end{proposition}

\begin{proof}
Since $\GzeroF\subseteq\actHF$ and
$\scrHF=\actHF\oplus\defectHF$, the final phenomenal support splits as
$\actHF=\GzeroF\oplus(\QzeroF\actHF)$.  The defect is orthogonal to the entire
final support and hence to the ground.
\end{proof}

The interpretation is therefore nested:
\begin{equation}
  \boxed{
  \text{phenomenal excess}
  \subset
  \text{differentiated phenomenal appearance}.}
  \label{eq:excess-subset-appearance}
\end{equation}
Not every sems-state is unactualizable.  Ordinary presentable appearances
occupy $\GzeroF\oplus\actAppF$; only the $\defectHF$ component lacks a
normalized physical partner.

\subsection{Cyclicity and phenomenal defect}

Phenomenal cyclicity remains
\begin{equation}
  \overline{\piF(\algF)\ket{\Omega_0}}=\scrHF.
  \label{eq:cyclicity-reaches-excess}
\end{equation}
Therefore the phenomenal determination algebra can densely generate both
final-support and defect directions from the same luminous ground.  The
defect is not a second dark phenomenal substance.  It is an unpaired
direction inside the one cyclic phenomenal representation.

Define the defect component generated by a determination $X$ by
\begin{equation}
  \eta_X^{\perp,A}
  :=\defectProjF\piF(X)\ket{\Omega_0}.
  \label{eq:excess-part-determination}
\end{equation}
The determination is fully contained in the final support on the ground when
\begin{equation}
  \eta_X^{\perp,A}=0.
  \label{eq:determination-fully-representable}
\end{equation}
The detailed classification of which v\=asan\=a-conditioned generators produce
defect components remains open.

\begin{openquestion}[Status of phenomenal excess]
\label{oq:status-phenomenal-excess}
Which phenomenal determinations densely generate vectors in $\Ker A^*$?
Candidate readings include unrealized imagination, counterfactual or
chimerical phenomenal possibility, and mathematically admissible but
cognitively inaccessible directions.  Only the abstract claim that these
directions have no normalized physical partner under the selected channel is
adopted here.
\end{openquestion}

\subsection{Final-support and defect weights for normal states}

For a normal phenomenal state $\rho$, define
\begin{equation}
  \boxed{
  \defectWeight(\rho)
  :=\Tr(\defectProjF\rho).}
  \label{eq:excess-weight-density}
\end{equation}
Then
\begin{equation}
  \boxed{
  \presentWeight(\rho)
  =1-\defectWeight(\rho).}
  \label{eq:actualization-excess-complement}
\end{equation}
The state is wholly in the final support iff $\defectWeight(\rho)=0$, and
wholly in the phenomenal defect space iff $\defectWeight(\rho)=1$.

\subsection{A finite pair-transfer realization}
\label{sec:finite-pair-realization}

The three-part decomposition can be realized in a finite quantum system
whose ground is selected by its Hamiltonian.  In a canonical matching
calculation, Singh assumes four fermion modes with ordinary positive
Hilbert-space adjoints and gives an attractive pair-transfer Hamiltonian
and its in-phase lower state
\cite[p.~13, \S3.4, Eq.~(3.11)]{Singh2026PairTransfer}.
That passage summarizes the calculation attributed there to Paper~8.
We verify below that this state is the unique ground of the entire
six-dimensional two-particle sector.  We use the example to check ground
compatibility together with access to the final defect, specifying the
paired input, fixed orbital rotation and occupation readout.  The prepared
states and their evolution determine the minimal output carrier.
All calculations take place in positive Hilbert spaces.  The two typed
carriers can then be placed in the graded direct sum of
\cref{sec:typed-repair}.  Deriving the assumed canonical Hilbert structure
and the available interventions from microscopic generalized trace
dynamics is a separate problem.

The cited Hamiltonian is a canonical component.  If it is to represent
the whole result of positive Gaussian exchange, its coupling matrix must
also have nonnegative self-channels with $G_{cc}G_{dd}\geq j^2$
\cite[Eq.~(3.12)]{Singh2026PairTransfer}.  Singh's minimal symmetric
completion makes the two excitation lines coincide and removes the
dynamically accessible defect for the preparation below.  The bare
component, and symmetric positive completions strictly above that
minimum, retain distinct lines.  We compute this distinction below;
the orbital model subsequently
retains its full projected interaction and uses an explicit field to
split its excitation doublet.

\statuslabel{Mathematical realization}
The construction verifies the mathematical conditions of
\cref{sec:axiomatic-spine} and satisfies strong polar ground compatibility.
The assignment of a phenomenal meaning to its output carrier is a separate
interpretive choice.

\subsubsection*{Hamiltonian, ground, and preparations}
Let $c_1,c_2,d_1,d_2$ obey the canonical anticommutation relations
$\{a_r,a_s^*\}=\delta_{rs}\idop$ and $\{a_r,a_s\}=0$, and let
$\ket{\mathrm{vac}}$ be the empty Fock state.  Work in the six-dimensional
two-particle space $\mathcal F_2$.  We write $j$ for Singh's coupling
$J=4\kappa>0$ in his scalar-trace convention.  The canonical component is
\begin{equation}
  P_c=c_2c_1,\qquad P_d=d_2d_1,\qquad
  H=-j(P_c^*P_d+P_d^*P_c),\qquad j>0 .
  \label{eq:pair-parent}
\end{equation}
Define $\ket C=P_c^*\ket{\mathrm{vac}}$,
$\ket D=P_d^*\ket{\mathrm{vac}}$ and
\begin{equation}
  g=\frac{\ket C+\ket D}{\sqrt2},\qquad
  e=\frac{\ket C-\ket D}{\sqrt2},\qquad
  m=\frac{c_1^*d_2^*+d_1^*c_2^*}{\sqrt2}\ket{\mathrm{vac}} .
  \label{eq:pair-energy-vectors}
\end{equation}
These vectors are orthonormal, and
$Hg=-jg$, $He=je$, $Hm=0$.  The remaining three two-particle states
also have energy zero.  Thus $g$ is the unique ground in the full
two-particle sector.

Take $\scrHP$ to be a typed copy of the paired plane
$\Span\{\ket C,\ket D\}$, with corresponding vectors $g_{\Psi},e_{\Psi}$.
Its ground generator $H_{\Psi}$ is the restriction of $H$ to that plane.
Fix an apparatus angle $0<\theta<\pi/4$, and define rotated orbital modes
\begin{equation}
  f_i^*=\cos\theta\,c_i^*+\sin\theta\,d_i^*,\qquad
  h_i^*=-\sin\theta\,c_i^*+\cos\theta\,d_i^*
  \quad(i=1,2).
  \label{eq:pair-orbital-rotation}
\end{equation}
The preparation channel is specified on the typed input basis by
\begin{equation}
 A_\theta\ket{C_\Psi}=f_1^*f_2^*\ket{\mathrm{vac}},\qquad
 A_\theta\ket{D_\Psi}=h_1^*h_2^*\ket{\mathrm{vac}} .
 \label{eq:pair-preparation-map}
\end{equation}
The notation $\ket0=g_\Psi$ designates the two-particle ground of the
input Hamiltonian.  The empty state $\ket{\mathrm{vac}}$ serves to
construct the pair basis.  The output evolves under
\eqref{eq:pair-parent}, with the observable
\begin{equation}
  Z_{\theta}=\tfrac12(N_f-N_h),\qquad
  N_f=\sum_i f_i^*f_i,\qquad N_h=\sum_i h_i^*h_i .
  \label{eq:pair-imbalance}
\end{equation}
Both $H$ and $Z_{\theta}$ preserve
$\scrHF:=\Span\{g,e,m\}$.  Put
$H_{\Phi}=\restr{H}{\scrHF}$, $u=\cos2\theta$,
$v=\sin2\theta$ and
\begin{equation}
  b=ue+vm,\qquad d=-ve+um .
  \label{eq:pair-bright-defect}
\end{equation}
Expansion of the rotated pair states gives
\begin{equation}
  A_{\theta}g_{\Psi}=g,\qquad A_{\theta}e_{\Psi}=b,
  \qquad A_{\theta}^*A_{\theta}=I_{\scrHP}.
  \label{eq:pair-channel}
\end{equation}
In the remainder of this example, $A=A_{\theta}$ has codomain $\scrHF$.
Write $\widetilde A=\iota A$ for the same preparation into $\mathcal F_2$,
where $\iota$ is the inclusion.  One can recover $\scrHF$ from the
preparation in $\mathcal F_2$.  Let $E_\lambda$ be the spectral
projection of $H$ there, and put $Q_g=\proj{g_\Psi}$, $Q_e=\proj{e_\Psi}$.
The weights $\widetilde A^*E_\lambda\widetilde A$ of the preparation
response are
\begin{equation}
 R_{-j}=Q_g,\qquad R_0=v^2Q_e,\qquad R_{+j}=u^2Q_e .
 \label{eq:pair-port-weights}
\end{equation}
Indeed, $E_{-j}g=g$, $E_0 b=vm$, and $E_{+j}b=ue$.  Since $uv\ne0$,
the three spaces $\Ran(E_\lambda\widetilde A)$ are the lines through
$g,m,e$, respectively.  Their span is the smallest reducing subspace
of $H$ containing the preparations.  It has dimension three, while
the direct preparation range has dimension two.

\begin{proposition}[Ground correspondence and a cyclic defect]
\label{prop:pair-ground-defect}
With $\ket0=g_{\Psi}$ and $\ket{\Omega_0}=g$, the channel satisfies
the strong ground bridge with $a_0=1$, is injective, and has
$\Ker A^*=\CC d$.  The grounds are cyclic for
$\algP=C^*(I_{\scrHP},H_{\Psi},Z_0)$ and
$\algF=C^*(I_{\scrHF},H_{\Phi},Z_{\theta})$ in their defining
representations.  Here $Z_0$ is the input occupation imbalance restricted
to the paired plane.
\end{proposition}

\begin{proof}
Equation \eqref{eq:pair-channel} gives $\absA=I_{\scrHP}$,
$\polA=A$, $A\ket0=g$ and $A^*g=\ket0$.  Its range is
$\Span\{g,b\}$, whose orthogonal complement in $\scrHF$ is $\CC d$.
The input imbalance interchanges $g_{\Psi}$ and $e_{\Psi}$.
On the output,
\begin{equation}
  Z_{\theta}g=b,\qquad Z_{\theta}b=g,\qquad Z_{\theta}d=0 .
  \label{eq:pair-readout-action}
\end{equation}
Since $H_{\Phi}b=ju e$ and $uv\ne0$, the vectors
$g,Z_{\theta}g,H_{\Phi}Z_{\theta}g$ span $\scrHF$.
This proves both cyclicity statements.
\end{proof}

In fact $\algP$ and $\algF$ are the full matrix algebras on their
respective carriers.  Each Hamiltonian has simple spectrum there.
On the input the imbalance joins its two eigenvectors.  On the output
the spectral projections isolate the two nonzero matrix elements
$\inner{e}{Z_\theta g}=u$ and $\inner{m}{Z_\theta g}=v$.
These matrix units, their adjoints and their products give all matrix
units.  Every nonzero vector is therefore cyclic.  The ground rays
are selected by $H_\Psi$ and $H_\Phi$.

\subsubsection*{Evolution and an outcome measurement of defect weight}
Starting from the prepared excitation $b$, set
$\rho_t=\proj{e^{-itH_{\Phi}/\hbar}b}$.  Direct evolution gives
\begin{equation}
  \defectWeight(\rho_t)
  =|\inner{d}{e^{-itH_{\Phi}/\hbar}b}|^2
  =\sin^2(4\theta)\sin^2\!\left(\frac{jt}{2\hbar}\right).
  \label{eq:pair-defect-evolution}
\end{equation}
At $\theta=\pi/8$ and $t=\pi\hbar/j$, the weight is one.
Ground preservation and departure from the direct range therefore coexist.

For a normalized state $\rho$ supported on $\scrHF$,
\eqref{eq:pair-readout-action} gives
\begin{equation}
  Z_{\theta}^{\,2}=AA^*=\suppProjF,\qquad
  \defectProjF=I_{\scrHF}-Z_{\theta}^{\,2},\qquad
  \defectWeight(\rho)=\Pr_{\rho}(Z_{\theta}=0)
                     =1-\Tr(\rho Z_\theta^2).
  \label{eq:pair-defect-measurement}
\end{equation}
The outcomes are $+1,0,-1$.  Zero records one fermion in each rotated
orbital.  Its frequency or, equivalently, the second moment of the
imbalance determines the defect weight.  The first moment alone can
vanish both for a state on the defect line and for a state with equal
probabilities of $+1$ and $-1$.

In $\mathcal F_2$, put $W=\mathcal F_2\ominus\scrHF$ and
$P_d=\proj{d}$.  The $+1$ and $-1$ eigenspaces are precisely the two
rotated pair lines in \eqref{eq:pair-preparation-map}.  Thus
$Z_\theta^2=\widetilde A\widetilde A^*$ and
$I_{\mathcal F_2}-Z_\theta^2=P_d+P_W$ on $\mathcal F_2$.  For an arbitrary
normalized state $\sigma$ on $\mathcal F_2$,
\begin{equation}
  \Pr_\sigma(Z_\theta=0)
  =\underbrace{\Tr(\sigma P_d)}_{\text{weight on the LG defect line}}
   +\underbrace{\Tr(\sigma P_W)}_{\text{weight outside the carrier}} .
  \label{eq:pair-ambient-zero}
\end{equation}
The ideal protocol preserves $\scrHF$.  If a preparation puts at most
$\varepsilon$ weight in $W$, the zero-outcome probability overestimates
the line weight by at most $\varepsilon$.

The centered ground correlation is
\begin{equation}
  \inner{g}{Z_{\theta}(t)Z_{\theta}g}
  =u^2e^{-2ijt/\hbar}+v^2e^{-ijt/\hbar},\qquad
  Z_{\theta}(t)=e^{itH_{\Phi}/\hbar}Z_{\theta}e^{-itH_{\Phi}/\hbar}.
  \label{eq:pair-two-frequency-response}
\end{equation}
Indeed, $Z_{\theta}g=b$ and $\inner{g}{Z_{\theta}g}=0$.
This scalar ground correlation measures the two excitation lines.
In a Hilbert-space realization with a self-adjoint Hamiltonian and a
stationary ground vector, distinct positive-weight atoms at the two
excitation energies require two orthogonal eigendirections orthogonal
to the ground.  This also gives a lower bound of three dimensions, which
$\scrHF$ attains.  At $\theta=0$ only the line through $e$ survives,
and at $\theta=\pi/4$ only the line through $m$ survives.
In each case the minimal output is the span of the ground and the
surviving direction.  The channel is onto this space, so its final defect
vanishes.  Keeping the original three-dimensional codomain at either
endpoint would retain the unused line $\CC d$.

\subsubsection*{Dependence on the preparation interface}
The channel's direct range belongs to the fixed preparation protocol.
For example, independent coherent preparations at $\theta=0$ and
$\theta=\pi/4$ reach $e$ and $m$ separately.  Their combined span together
with $g$ fills $\scrHF$.  Explicitly, the coherent map on a doubled
input, $L(x,y)=A_0 x+A_{\pi/4}y$, has range $\scrHF$ and hence zero
final defect.  Both preparations send $g_\Psi$ to $g$, so
$\Ker L=\CC(g_\Psi,-g_\Psi)$ by rank--nullity.  On this doubled input
the map fails both injectivity (R4) and strict excess (R5).

The rotation in \eqref{eq:pair-orbital-rotation} is compatible with the
ground for an algebraic reason.  Define
$T=\sum_i c_i^*d_i$,
$X=(T+T^*)/2$, $Y=(T-T^*)/(2i)$ and
$Z=(N_c-N_d)/2$.  The anticommutation relations give
\begin{equation}
  T^2=2P_c^*P_d,\qquad Xg=m,\qquad Yg=0,\qquad Zg=e .
  \label{eq:pair-mixer-selection}
\end{equation}
With these signs $U_\theta=\exp(-2i\theta Y)$ conjugates $c_i^*,d_i^*$
to $f_i^*,h_i^*$, respectively, and fixes $\ket{\mathrm{vac}}$.
Its restriction to the paired input plane gives
\eqref{eq:pair-preparation-map}, and $Yg=0$ gives $U_\theta g=g$.
Up to a scalar on this particle-number sector, a number-conserving
one-body orbital drive acting identically on both internal labels is
$\alpha X+\beta Y+\gamma Z$, with real coefficients.  Its ground ray is invariant throughout
a continuous pulse exactly when $\alpha=\gamma=0$.  The surviving generator $Y$
produces the real orbital rotations.  Access to that drive and a fixed
preparation angle are apparatus assumptions.  Singh's canonical example
supplies \eqref{eq:pair-parent} and its ground; it leaves those
interventions to a further physical model.  The calculations above
establish the realization for the stated protocol and give the defect
weight an outcome interpretation within it.

\subsubsection*{Effect of exchange self-channels}

With pair annihilators $P_c=c_2c_1$ and $P_d=d_2d_1$, define symmetric
self-channels in the same two-particle sector by
\begin{equation}
 H_s=-s(P_c^*P_c+P_d^*P_d)
       -j(P_c^*P_d+P_d^*P_c),\qquad j>0,\quad s\geq0.
 \label{eq:pair-completed-hamiltonian}
\end{equation}
The exchange matrix is positive semidefinite exactly when $s\geq j$.
This is the condition for the stated positive-exchange completion;
the bare finite Hamiltonian at $s=0$ is still a bounded self-adjoint
operator on a positive Hilbert space.
Direct application to the vectors already defined gives
\begin{equation}
 H_sg=-(s+j)g,\qquad H_se=(j-s)e,\qquad
 H_sm=0,\qquad H_s|_W=0.
 \label{eq:pair-completed-spectrum}
\end{equation}
Thus $g$ remains the unique ground.  For the same preparation channel
$A$, with $b=ue+vm$ and $d=-ve+um$, the correlation and defect weight are
\begin{align}
 \langle g,Z_\theta(t)Z_\theta g\rangle
 &=u^2e^{-2ijt/\hbar}+v^2e^{-i(s+j)t/\hbar},\\
 \defectWeight\bigl(\proj{e^{-itH_s/\hbar}b}\bigr)
 &=\sin^2(4\theta)\sin^2\!\left(\frac{(j-s)t}{2\hbar}\right),
 \label{eq:pair-completed-response}
\end{align}
where $Z_\theta(t)$ here uses $H_s$ for the evolution.  These formulas
follow by evolving the two orthogonal excitation components of $b$.
For $uv\ne0$ and $s\ne j$, their two frequencies are distinct and
the preparation generates $\Span\{g,e,m\}$.

At the minimal positive completion $s=j$,
$H_s=-j(P_c+P_d)^*(P_c+P_d)$ and the six-dimensional spectrum is
$\{-2j,0,0,0,0,0\}$.  The range $\Span\{g,b\}$ is invariant.
If the output carrier is kept equal to $\Span\{g,e,m\}$, its algebraic
final defect $\Ker A^*=\CC d$ persists but is dynamically inaccessible.
The ground is then not cyclic on that carrier for $C^*(I,H_s,Z_\theta)$.
On the response-minimal carrier $\Span\{g,b\}$, the channel is onto
and its final defect is zero.  Within symmetric positive completions,
a dynamically accessible defect therefore requires $s>j$.
In the displayed basis $[Y,H_s]e=i(j-s)m$; the minimal completion
commutes with $Y$, whereas the bare term and the non-minimal symmetric
completions break this orbital rotation symmetry.

For unequal self-channels $s_c,s_d\geq0$ with $s_cs_d\geq j^2$,
the paired block is $-\left(\begin{smallmatrix}s_c&j\\j&s_d\end{smallmatrix}\right)$.
Writing its ground as $\alpha C+\beta D$, one has
$Y(\alpha C+\beta D)=i(\alpha-\beta)m/\sqrt2$; $\alpha=\beta$ is possible exactly when
$s_c=s_d$.  Thus asymmetric completion also removes preservation of
the ground by the specified continuous rotation.  This does not exclude
a different preparation derived from additional data.

\subsubsection*{Which operations the pair component supplies}

The pair component does not by itself supply the preparation mixer.
Let $\Pi_c=(-1)^{N_c}$.  The pair Hamiltonian, both Hermitian
quadratures of $P_c^*P_d$, the imbalance $Z$, and the self-channels
$P_c^*P_c,P_d^*P_d$ commute with $\Pi_c$ and preserve
$\Span\{C,D\}$.  The vector $m$ has the opposite $\Pi_c$ eigenvalue.
No sequence generated by these controls therefore reaches $m$ from
the paired plane.  The one-body mixer $Y$ supplies an additional
operation.  This is a restriction on the stated control family,
not a prohibition on a larger system with further sources.

There is also an exact internal-symmetry test.  Suppose the columns
$c=(c_1,c_2)^T$ and $d=(d_1,d_2)^T$ transform in closed two-dimensional
unitary representations $R_c(h),R_d(h)$ of a group (the annihilator
convention is $c\mapsto R_c(h)c$, $d\mapsto R_d(h)d$).  For a fixed
numerical matrix $K$, put
\begin{equation}
 T_K=c^*Kd,\qquad Y_K=(T_K-T_K^*)/(2i).
 \label{eq:pair-general-mixer}
\end{equation}
Here $c^*=(c_1^*,c_2^*)$.  The mixer is invariant precisely when
$R_c(h)K=KR_d(h)$ for every $h$.  In the pair-phase convention
of \eqref{eq:pair-energy-vectors}, the CAR give
\begin{equation}
 \|Y_Kg\|^2
 =\frac14\left(
 |K_{11}-\overline{K_{22}}|^2+
 |K_{12}+\overline{K_{21}}|^2\right).
 \label{eq:pair-mixer-ground-certificate}
\end{equation}
The vector $Y_Kg$ is orthogonal to $g$, so preservation of the ground
ray by this mixer is equivalent to $Y_Kg=0$.

\begin{proposition}[Invariant ground-preserving mixer]
\label{prop:pair-invariant-mixer}
A nonzero invariant mixer of the form \eqref{eq:pair-general-mixer}
annihilating $g$ exists if and only if $R_c$ and $R_d$ are unitarily
equivalent.  The original aligned mixer $K=I_2$ is invariant exactly
when the two representation matrices agree in the displayed bases.
\end{proposition}
\begin{proof}
Equation \eqref{eq:pair-mixer-ground-certificate} vanishes exactly for
\begin{equation}
 K=\begin{pmatrix}z&w\\-\overline w&\overline z\end{pmatrix}.
\end{equation}
For $K\ne0$, this gives $K=rU$, where
$r=\sqrt{|z|^2+|w|^2}>0$ and $U\in SU(2)$.  Invariance then makes
$U$ a unitary intertwiner.  Conversely, multiplying any unitary
intertwiner by a scalar phase makes its determinant one, without
changing its intertwining property; the resulting $SU(2)$ matrix
has the displayed form.
\end{proof}

Pair transfer requires only equality of the determinant characters
$\det R_c=\det R_d$, a weaker condition.  For example, diagonal
$U(1)$ charges $(+1,-1)$ on $c$ and $(+2,-2)$ on $d$ permit the pair
term but no invariant one-particle mixer between the blocks.
A charged mediator or reference can permit a joint invariant
interaction; it must then be retained in the preparation and
measurement model.  Even an allowed intertwiner does not establish
that a source supplies its required Hermitian quadrature.

For the symmetric family with $s\ne j$, the unused space $W$ in
\eqref{eq:pair-ambient-zero} remains distinct from the cyclic defect
line $\CC d$.  The common-label orbital
currents $X,Y,Z$ annihilate $W$, but particle-number conservation
alone does not protect it: with
$w_0=(c_1^*d_2^*-d_1^*c_2^*)\ket{\mathrm{vac}}/\sqrt2$,
\begin{equation}
 Q:=n_{c_1}-n_{c_2}-n_{d_1}+n_{d_2},\qquad
 Qw_0=2m,\quad Qm=2w_0,\quad Qg=Qe=0.
\end{equation}
For every $s$, $m$ and $W$ have the same zero energy under $H_s$.
A perturbation proportional to $Q$ therefore mixes $m$ and $w_0$
without an energy gap suppressing that mixing; in particular,
$Qd=2u w_0$.  The leakage magnitude still depends on the perturbation
strength and duration.  Thus an independently specified current family is needed for any
protection claim.  Adding $W$ to the output carrier while retaining
only the original orbital-current algebra would also destroy
cyclicity of $g$ on that enlarged carrier.


\subsection{An orbital source preserving the full ground}
\label{sec:orbital-ground-source}

The pair-transfer example specifies its admissible mixer separately from
its Hamiltonian.  A different, independently specified orbital model
supplies both a ground and a ground-preserving rotation source.  This
construction does not identify the full local interaction with the
pair-transfer term alone, and does not derive the model from microscopic
generalized trace dynamics.

The excitation splitting has a different symmetry origin here.
In the pair family, $s\ne j$ breaks the orbital rotations generated
by $Y$.  The orbital parent below commutes with total angular momentum
even at nonzero field: it preserves those rotations, while the field
breaks zero-field time reversal and separates the opposite angular-momentum
branches.  Both constructions retain their ground-preserving operation
in the parameter ranges stated.

\statuslabel{Mathematical realization within a specified positive parent}
Work on the antisymmetric four-particle Hilbert space
\(\bigwedge^4(L^2(\mathbb R^2)\otimes\CC^2)\).  Both internal labels
have mass \(M>0\) and the same nonzero gauge charge \(q\).  Consider
\begin{align}
 H_{\lambda,\mathcal B}
 &=\sum_{i=1}^4\left[
 \frac{(\mathbf p_i-q\mathbf A_{\mathcal B}(\mathbf r_i))^2}{2M}
 +\frac{M\omega_0^2r_i^2}{2}\right]
 -\frac{\lambda}{2\pi a^2}\sum_{i<j}
 e^{-|\mathbf r_i-\mathbf r_j|^2/(2a^2)},
 \label{eq:orbital-parent}\\
 \mathbf A_{\mathcal B}(x,y)&=\frac{\mathcal B}{2}(-y,x),\qquad
 \lambda>0,\quad a>0,\quad
 \Omega_{\mathrm B}=\frac{q\mathcal B}{2M},\quad
 \omega^2=\omega_0^2+\Omega_{\mathrm B}^2.
 \label{eq:orbital-bias-parameters}
\end{align}
Only the orbital field coupling is included.  Spin Zeeman coupling, if
present in an implementation, must be included in the ground-level
crossing test.  Expanding minimal coupling gives the exact identity
\begin{equation}
 H_{\lambda,\mathcal B}
 =H_{\lambda,0}^{(\omega)}-\Omega_{\mathrm B}L_{\mathrm{tot}},
 \qquad L_{\mathrm{tot}}=\sum_i(x_ip_{yi}-y_ip_{xi}).
 \label{eq:orbital-minimal-coupling}
\end{equation}
Keeping the same \(H_{\lambda,0}^{(\omega)}\) as the field varies
requires the trap compensation
\(\omega_0(\mathcal B)^2=\omega^2-\Omega_{\mathrm B}(\mathcal B)^2>0\).
The source is a physical field; changing coordinates to a rotating frame
alone would not change the laboratory response.

\subsubsection*{Ground selection including the filled core}
Set \(\ell=\sqrt{\hbar/(M\omega)}\), and use the oscillator orbitals
\begin{equation}
 \phi_s=\frac{e^{-r^2/(2\ell^2)}}{\sqrt\pi\ell},\qquad
 \phi_x=\frac{\sqrt2x}{\ell}\phi_s,\qquad
 \phi_y=\frac{\sqrt2y}{\ell}\phi_s.
 \label{eq:orbital-three-orbitals}
\end{equation}
At zero interaction and bias, two fermions fill \(\phi_s\).  The remaining
two occupy the four \(p\)-shell spin modes.  This six-dimensional
\(s^2p^2\) space is the entire ground eigenspace, at energy \(6\hbar\omega\),
separated from other configurations by \(\hbar\omega\).  Use the vectors
\(g,e,m\) of \cref{eq:pair-energy-vectors}, with
\(\ket{\mathrm{vac}}\) there replaced by the filled-core state and the
two internal labels identified with the two spin labels.  The three
remaining states \(W\) are the spin triplet.  Direct orbital differentiation
and the canonical anticommutation relations give
\begin{equation}
 L_{\mathrm{tot}}=2\hbar Y\quad\hbox{on }s^2p^2,
 \qquad L_{\mathrm{tot}}g=0,\qquad
 L_{\mathrm{tot}}e=2i\hbar m,\qquad
 L_{\mathrm{tot}}m=-2i\hbar e.
 \label{eq:orbital-derived-current}
\end{equation}
For \(\xi=a^2/\ell^2\) and \(v=-\lambda/(8\pi\ell^2)\), Gaussian
integration gives the projected interaction energies
\begin{align}
 E_g^{(1)}&=E_{\mathrm{core}}^{(1)}
       +\frac{4v(\xi^2+\xi+1)}{(1+\xi)^3},\\
 E_e^{(1)}=E_m^{(1)}&=E_{\mathrm{core}}^{(1)}
       +\frac{2v(2\xi^2+2\xi+1)}{(1+\xi)^3},\\
 E_W^{(1)}&=E_{\mathrm{core}}^{(1)}
       +\frac{4v\xi}{(1+\xi)^2},\qquad
 E_{\mathrm{core}}^{(1)}=\frac{4v(5\xi+2)}{(1+\xi)^2}.
 \label{eq:orbital-projected-energies}
\end{align}
These are the exact entries of the first-order degenerate-perturbation
matrix, including core contractions and exchange.  They are not
all-orders energies of the continuum parent.  In particular, writing
\(d_a=[4\pi\ell^2(1+\xi)^3]^{-1}\), the singlet and triplet gaps
above \(g\) are \(\lambda d_a\) and \(2\lambda d_a\).

\begin{proposition}[Full ground preservation and split excitations]
\label{prop:orbital-full-ground}
Fix \(a,\omega>0\) and \(0<|\chi|<1\).  For all sufficiently small
\(\lambda>0\), choose
\(\Omega_{\mathrm B}=\chi\lambda d_a/(2\hbar)\) and compensate the trap
as above.  The full fixed-particle-number Hamiltonian has a unique ground
\(G_\lambda\) of energy \(E_0(\lambda)\), independent of this
compensated bias, with
\(L_{\mathrm{tot}}G_\lambda=0\).  There are normalized excited vectors
\(F_{\pm,\lambda}\) such that
\begin{align}
 L_{\mathrm{tot}}F_{\pm,\lambda}&=\pm2\hbar F_{\pm,\lambda},\\
 H_{\lambda,\mathcal B}F_{\pm,\lambda}
   &=(E_2(\lambda)\mp2\hbar\Omega_{\mathrm B})F_{\pm,\lambda}.
\end{align}
Their energy gaps above the ground are positive and distinct.
\end{proposition}

\begin{proof}
At fixed \(a>0\), the interaction has operator norm at most
\(6\lambda/(2\pi a^2)\).  Analytic perturbation theory for the isolated
sixfold oscillator cluster applies \cite{Kato,ReedSimon}.
The simple minimum in \cref{eq:orbital-projected-energies} selects a unique
exact ground branch.  Rotation and spin symmetry keep it in the
orbital-angular-momentum-zero spin-singlet sector.  The branches continuing
\((e\pm im)/\sqrt2\) have angular momenta \(\pm2\hbar\);
time reversal at zero bias makes their common energy \(E_2(\lambda)\).
Equation \eqref{eq:orbital-minimal-coupling} then gives the stated exact
field shifts without changing the eigenvectors.  The singlet gaps are
\(\lambda d_a(1\mp\chi)+O(\lambda^2)\), and the triplet gap is
\(2\lambda d_a+O(\lambda^2)\).  The strict margin \(|\chi|<1\) keeps
these positive for small \(\lambda\).  To control the higher spectrum despite the unbounded magnetic
term, write the attraction as $-\lambda W_a$, where
$0\leq W_a\leq 6/(2\pi a^2)$, and let $H_{\mathrm{ho}}$ be the
four-particle oscillator at frequency $\omega$.  In their common
oscillator basis,
\begin{equation}
 |L_{\mathrm{tot}}|\leq H_{\mathrm{ho}}/\omega-4\hbar I.
\end{equation}
Consequently, as quadratic forms,
\begin{equation}
 H_{\lambda,\mathcal B}\geq
 (1-|\Omega_{\mathrm B}|/\omega)H_{\mathrm{ho}}
 +4\hbar|\Omega_{\mathrm B}|I-\lambda\|W_a\|I.
\end{equation}
The min--max principle keeps the seventh eigenvalue at least
$7\hbar\omega-3\hbar|\Omega_{\mathrm B}|-\lambda\|W_a\|$.
It therefore remains separated from the six branches near
$6\hbar\omega$ when $\Omega_{\mathrm B}=O(\lambda)$.  This proves
the global fixed-particle-number ground claim.  The small-coupling
statement is for fixed $a,\omega,\chi$; it is not uniform in the
contact limit $a\to0$.
\end{proof}

\subsubsection*{An exact selected carrier and its physical quadrupole}
The spatial observable
\begin{equation}
 Q_x=\sum_i\frac{x_i^2-y_i^2}{2\ell^2},\qquad
 Q_y=\sum_i\frac{x_iy_i}{\ell^2}
 \label{eq:orbital-quadrupoles}
\end{equation}
has valence-shell compressions \(Z\) and \(X\) from
\cref{eq:pair-mixer-selection}.  In particular its low transition matrix
element from \(g\) to \(e\) is nonzero.  This persists for sufficiently
small \(\lambda\): $Q_x$ is $H_{\mathrm{ho}}$-bounded and the
isolated eigenbranches are continuous in the oscillator graph norm.  Choose the phases of \(F_{\pm,\lambda}\) by these
equal-magnitude quadrupole matrix elements, and put
\begin{equation}
 e_\lambda=\frac{F_{+,\lambda}+F_{-,\lambda}}{\sqrt2},\qquad
 m_\lambda=\frac{F_{+,\lambda}-F_{-,\lambda}}{i\sqrt2},\qquad
 \scrHF=\Span\{G_\lambda,e_\lambda,m_\lambda\}.
 \label{eq:orbital-low-carrier}
\end{equation}
This is an exact reducing subspace of the full Hamiltonian.  Let
\(\scrHP=\CC^2\) with orthonormal basis \(\ket0,\ket1\), and specify
\begin{equation}
 A\ket0=G_\lambda,\qquad A\ket1=e_\lambda,\qquad
 H_\Phi=H_{\lambda,\mathcal B}|_{\scrHF},\qquad
 H_\Psi=A^*H_\Phi A=\diag(E_0(\lambda),E_2(\lambda)).
 \label{eq:orbital-typed-preparation}
\end{equation}
The input Hamiltonian is a compression, not a restriction to an invariant
prepared plane.  Here \(A^*A=I_{\scrHP}\), the ground bridge is strong with
\(a_0=1\), and \(\Ker A^*=\CC m_\lambda\).  With \(P_{\mathrm{low}}\)
the projection onto \(\scrHF\), the normalized filtered quadrupole is
\begin{equation}
 Z_\lambda=
 \frac{P_{\mathrm{low}}Q_xP_{\mathrm{low}}}
 {\|P_{\mathrm{low}}Q_xG_\lambda\|}
 =\ket{G_\lambda}\!\bra{e_\lambda}
  +\ket{e_\lambda}\!\bra{G_\lambda},
 \qquad Z_\lambda^2=AA^*.
 \label{eq:orbital-filtered-readout}
\end{equation}
Define $Z_\Psi=A^*Z_\lambda A$.  The grounds are cyclic for
$C^*(I,H_\Psi,Z_\Psi)$ and $C^*(I,H_\Phi,Z_\lambda)$ in the
defining representations.  As in the preceding
finite example these are full matrix algebras; cyclicity does not select
the grounds.  For $\rho_t=\proj{e^{-itH_\Phi/\hbar}e_\lambda}$
and $Z_\lambda(t)=e^{itH_\Phi/\hbar}Z_\lambda e^{-itH_\Phi/\hbar}$,
one obtains
\begin{align}
 \defectWeight(\rho_t)&=\sin^2(2\Omega_{\mathrm B}t),
 &t_*&=\frac{\pi}{4|\Omega_{\mathrm B}|},\\
 \inner{G_\lambda}{Z_\lambda(t)Z_\lambda G_\lambda}
 &=e^{-i(E_2-E_0)t/\hbar}\cos(2\Omega_{\mathrm B}t).
 \label{eq:orbital-low-response}
\end{align}
For a normalized state $\rho$ supported on this carrier,
$\defectWeight(\rho)=1-\Tr(\rho Z_\lambda^2)=\Pr_\rho(Z_\lambda=0)$.
Both response atoms have weight \(1/2\).  Their distinct frequencies
force both excited directions into the minimal carrier of this port.
The identity for \(Z_\lambda^2\) holds on the selected carrier.  Spectral
selection, the normalized preparation and a detector implementing the
filtered readout remain specified operational requirements.

\subsubsection*{The unfiltered observable need not isolate one defect line}
Define the normalized vector
\(e_Q=Q_xG_\lambda/\|Q_xG_\lambda\|\), and choose the direct preparation
range \(\Span\{G_\lambda,e_Q\}\).  Let \(C_0(t)\) denote the normalized
zero-bias ground correlation of \(Q_x\), at the same compensated
oscillator frequency \(\omega\).  Rotational covariance and zero-bias
time reversal pair the full angular-momentum \(\pm2\hbar\) spectral
measures.  For evolution starting in $e_Q$, the exact unfiltered correlation
and full defect weight are
\begin{equation}
 C_{\mathcal B}(t)=C_0(t)\cos(2\Omega_{\mathrm B}t),\qquad
 p_{\mathrm{defect}}(t)=1-|C_0(t)|^2\cos^2(2\Omega_{\mathrm B}t).
 \label{eq:orbital-unfiltered-response}
\end{equation}
The evolving state remains in the minimal reducing carrier generated by
this port, so the second identity measures its departure from the direct
range within that carrier.  Complete departure at \(t_*\) does not imply
transfer into a specified one-dimensional defect, nor ensure a later
complete return.  The minimal output of the unfiltered port may be larger
than three dimensions.  
This example identifies a physical rotation current and constrains its
coefficient within a specified positive orbital parent.  The interaction,
trap compensation, preparation restriction and detector model remain
realization data.  The output carrier may be given a phenomenal interpretation only
through additional identification.  These calculations do not derive
a gravitational coupling or a microscopic GTD reconstruction.  Also, \([H,L_{\mathrm{tot}}]=0\) does not prevent the
excursion: the prepared real excitation combines distinct chiral energy
eigenstates, and its line is not invariant under their relative phase
evolution.


\subsection{Canonical deconditioning removes defect weight}

The canonical deconditioning semigroup $\Decond_t$ of
\cref{thm:canonical-deconditioning-semigroup} attracts every normal state to
$\PzeroF$.  Because
\begin{equation}
  \defectProjF\leq\QzeroF,
  \label{eq:excess-under-ground-complement}
\end{equation}
its action on defect weight is exact.

\begin{theorem}[Exponential contraction of phenomenal excess]
\label{thm:deconditioning-contracts-excess}
For this canonical deconditioning semigroup,
\begin{equation}
  \boxed{
  \defectWeight(\Decond_t(\rho))
  =e^{-\gamma t}\defectWeight(\rho).}
  \label{eq:excess-exponential-decay}
\end{equation}
Equivalently,
\begin{equation}
  \boxed{
  \presentWeight(\Decond_t(\rho))
  =1-e^{-\gamma t}
    \bigl(1-\presentWeight(\rho)\bigr).}
  \label{eq:actualizability-increases}
\end{equation}
\end{theorem}

\begin{proof}
Write $P=\PzeroF$, $Q=\QzeroF$, and $E=\defectProjF$.  Since $E\leq Q$ and
$EP=PE=0$, all ground and cross terms in the formula for $\Decond_t$ vanish
when traced against $E$.  The $Q\rho Q$ term contributes
$e^{-\gamma t}\Tr(E\rho)$.
\end{proof}

\statuslabel{Derived proposition within the canonical model} The theorem does
not establish that all contemplative practice follows this semigroup.  It
shows that the chosen minimal deconditioning dynamics simultaneously removes
ground deviation and unpaired phenomenal excess.

\subsection{Return after deconditioning}

Under deconditioning,
\begin{equation}
  \presentWeight(\Decond_t\rho)\longrightarrow1,
  \qquad
  \Decond_t(\rho)\longrightarrow\PzeroF
  \quad\text{in trace norm}.
  \label{eq:deconditioning-return-limits}
\end{equation}
There are two distinct return statements.

First, the restricted ground operation satisfies
\begin{equation}
  \returnGroundOp(\Decond_t\rho)
  =\Tr(\PzeroF\Decond_t\rho)\,\PzeroP.
  \label{eq:restricted-return-after-deconditioning}
\end{equation}
For every $t>0$ the coefficient is positive, and conditional normalization
returns exactly the physical vacuum state $\PzeroP$.

Second, the global conditional return satisfies
\begin{equation}
  \boxed{
  \returnOpCond(\Decond_t\rho)
  \longrightarrow
  \polA^*\PzeroF\polA
  =\proj{0_A}.}
  \label{eq:global-return-after-deconditioning}
\end{equation}
This limit equals $\PzeroP$ exactly under full polar ground reciprocity,
namely when $\ket{0_A}=\ket0$.  The two claims should not be conflated.

\subsection{No equation of phenomenal defect with non-luminosity}

A vector in $\defectHF$ is still phenomenal and therefore luminous in the
structural rang-rig sense of \cref{id:rangrig-terminal-law}.  What it
lacks is physical presentability under the selected Awareness--Volition
channel, not phenomenal manifestness.  Accordingly,
\begin{equation}
  \boxed{
  \phi\in\Ker A^*
  \not\Longrightarrow
  \text{absence of luminosity}.}
  \label{eq:excess-not-nonluminosity}
\end{equation}
It means only that the normalized Volition orientation sends that direction
to zero.
% =============================================================================
% END INLINED FILE: sections/36_phenomenal_excess_deconditioning.tex
% =============================================================================
% =============================================================================
% BEGIN INLINED FILE: sections/37_fredholm_index_excess.tex
% =============================================================================
\section{Fredholm index as sectoral imbalance}
\label{sec:fredholm-index-excess}

The kernel and cokernel data of the typed Awareness channel measure the
unpaired directions on the two sides of the architecture.  When these data
are finite and the range is closed, their difference is the Fredholm index
\cite{Kato,ReedSimon}.

\subsection{Physical null directions and phenomenal defect}

Define
\begin{equation}
  \mathscr Z_{\Psi}(A):=\Ker A,
  \qquad
  \mathscr Z_{\Phi}(A):=\Ker A^*.
  \label{eq:sectoral-zero-spaces}
\end{equation}
The first consists of physical directions in $\Dom A$ with no Awareness
image.  The second consists of phenomenal directions with no normalized
Volition return.  Physical presentability sets the first space to zero;
strict phenomenal excess makes the second nonzero.

For a closed densely defined operator with closed range, the cokernel is
canonically represented by
\begin{equation}
  \coker A
  :=\scrHF/\Ran A
  \simeq\Ker A^*.
  \label{eq:cokernel-kernel-adjoint}
\end{equation}

\begin{definition}[Fredholm Awareness channel]
\label{def:fredholm-awareness}
The closed channel $A$ is Fredholm when
\begin{equation}
  \dim\Ker A<\infty,
  \qquad
  \Ran A\ \text{is closed},
  \qquad
  \dim\Ker A^*<\infty.
  \label{eq:fredholm-conditions}
\end{equation}
Its index is
\begin{equation}
  \boxed{
  \ind A
  :=\dim\Ker A-\dim\Ker A^*.}
  \label{eq:fredholm-index-awareness}
\end{equation}
\end{definition}

\begin{theorem}[Index reading under asymmetric presentability]
\label{thm:index-as-phenomenal-excess}
Assume \cref{ax:physical-presentability} and suppose $A$ is Fredholm.  Then
\begin{equation}
  \boxed{
  \ind A=-\dim\Ker A^*.}
  \label{eq:index-negative-excess}
\end{equation}
Consequently, strict finite-dimensional phenomenal excess is equivalent to
\begin{equation}
  \ind A<0.
  \label{eq:negative-index-excess}
\end{equation}
\end{theorem}

\begin{proof}
Physical presentability gives $\dim\Ker A=0$.  Substitute into
\eqref{eq:fredholm-index-awareness}.
\end{proof}

\begin{identification}[Index as unpaired sector count]
\label{id:index-unpaired-count}
For the orientation $A:\Dom A\subseteq\scrHP\to\scrHF$,
\begin{equation}
  \boxed{
  \ind A
  =
  \#(\text{unpaired physical directions})
  -
  \#(\text{unpaired phenomenal directions})}
  \label{eq:index-doctrinal-reading}
\end{equation}
when the Fredholm hypotheses make both counts finite.  A negative index
therefore records phenomenal surplus; reversing the orientation reverses the
sign:
\begin{equation}
  \ind A^*=-\ind A.
  \label{eq:adjoint-index-sign}
\end{equation}
\end{identification}

\subsection{Two-term closed Hilbert complex}

Equip $\Dom A$ with the graph norm
\begin{equation}
  \|\psi\|_A^2:=\|\psi\|_{\Psi}^2+\|A\psi\|_{\Phi}^2.
\end{equation}
Because $A$ is closed, $(\Dom A,\|\cdot\|_A)$ is a Hilbert space and the
typed channel defines the two-term closed Hilbert complex
\begin{equation}
  0\longrightarrow(\Dom A,\|\cdot\|_A)
  \xrightarrow{\ A\ }
  \scrHF
  \longrightarrow0.
  \label{eq:two-term-complex}
\end{equation}
For closed range, define
\begin{equation}
  \cohZero:=\Ker A,
  \qquad
  \cohOne:=\scrHF/\Ran A\simeq\Ker A^*.
  \label{eq:transition-cohomology}
\end{equation}
Then
\begin{equation}
  \boxed{
  \ind A=\dim\cohZero-\dim\cohOne.}
  \label{eq:index-euler-characteristic}
\end{equation}
Thus phenomenal excess is the degree-one cohomology of the transition
complex when physical null directions have been removed.

\subsection{Stability through the bounded transform}

For an unbounded closed Fredholm channel define the bounded transform
\begin{equation}
  F_A:=A(I+A^*A)^{-1/2}.
  \label{eq:bounded-transform-awareness}
\end{equation}
It has the same kernel, adjoint kernel, and index as $A$.

The stable statement is made at this bounded level: the index is constant
under norm-continuous paths of bounded Fredholm transforms $F_{A_s}$ and is
unchanged by compact perturbations $F_A\mapsto F_A+K$ that remain Fredholm.
An unbounded family may instead be treated in a separately specified gap or
graph topology under hypotheses for which Fredholmness and index are known to
be locally constant.  No equivalence of these topologies is asserted here, and
no claim is made that an arbitrary compact-looking change of the raw unbounded
formula preserves its domain or Fredholm character.

Hence a nonzero finite phenomenal surplus cannot be removed by such an
admissible small compact repair without crossing a non-Fredholm point.

\statuslabel{Imported theorem} Fredholm index stability and the bounded
transform are standard operator theory \cite{Kato,ReedSimon}.  Their
interpretation as sectoral imbalance belongs to the cognitive realization.

\subsection{Finite and infinite phenomenal defect}

Without assuming Fredholmness,
\begin{equation}
  \boxed{
  \dim\Ker A^*\in\{0,1,2,\ldots,\infty\}.}
  \label{eq:phenomenal-excess-multiplicity}
\end{equation}
If $\dim\Ker A^*=\infty$, the ordinary Fredholm index is not a finite integer.
The existence of a strictly larger phenomenal domain therefore does not by
itself imply that $A$ is Fredholm.  The finite index is a special robust
regime, not the universal measure of every possible surplus.

\subsection{The index-minus-one shift model}

Let
\begin{equation}
  \scrHP=\ell^2(\NN_0)=\overline{\Span}\{f_0,f_1,\ldots\},
  \qquad
  \scrHF=\ell^2(\NN_0)=\overline{\Span}\{g_0,g_1,\ldots\},
\end{equation}
and define the unilateral shift
\begin{equation}
  u_A f_n=g_{n+1},
  \label{eq:unilateral-shift-example}
\end{equation}
one has
\begin{equation}
  \Ker u_A=\{0\},
  \qquad
  \Ker u_A^*=\CC g_0.
\end{equation}
Therefore
\begin{equation}
  \boxed{\ind u_A=-1.}
  \label{eq:shift-index-minus-one}
\end{equation}
This is the canonical elementary model of one robust unpaired phenomenal
direction.  Choosing $\ket0=f_0$ and $\ket{\Omega_0}=g_1$ places the luminous
ground in the final support while $g_0$ spans the one-dimensional phenomenal
defect.

\subsection{Index and the luminous ground}

The ground bridge guarantees
\begin{equation}
  \ket0\notin\Ker A,
  \qquad
  \ket{\Omega_0}\notin\Ker A^*.
  \label{eq:grounds-not-index-zero-modes}
\end{equation}
Therefore the distinguished cyclic grounds do not themselves contribute to
the index imbalance.  The index counts unpaired directions beyond the
bridge-selected ground relation.
% =============================================================================
% END INLINED FILE: sections/37_fredholm_index_excess.tex
% =============================================================================
% =============================================================================
% BEGIN INLINED FILE: sections/38_current_architecture.tex
% =============================================================================
\section{Current architecture, assumptions, and open problems}
\label{sec:current-architecture}

The cognitive realization distinguishes three levels of the transition
channel:
\begin{equation}
  \boxed{
  \begin{array}{c|c|c}
  \text{level} & \text{Awareness} & \text{Volition}\\
  \hline
  \text{typed raw channel} & A & A^*\\
  \text{normalized active orientation} & u_A & u_A^*\\
  \text{restricted ground orientation} & u_0 & u_0^*
  \end{array}}
  \label{eq:three-level-AV-table}
\end{equation}
The raw channel retains domain and spectral-magnitude information.  Its polar
orientation carries the canonical projective and Born geometry on active
support.  The restricted orientation records the ground bridge independently
of whether the physical vacuum is a singular vector of the full channel.

\subsection{Current architecture}

The resulting structure is as follows.
\begin{enumerate}[label=(\roman*)]
  \item The physical and phenomenal sectors carry unique cyclic ground rays
  represented by $\ket0$ and $\ket{\Omega_0}$, joined by the realization
  axiom $A\ket0=a_0\ket{\Omega_0}$ with $a_0>0$.
  \item The restricted ground return is exact,
  \begin{equation}
    u_0^*\ket{\Omega_0}=\ket0,
  \end{equation}
  whereas the full polar return satisfies
  \begin{equation}
    u_A^*\ket{\Omega_0}=a_0^{-1}|A|\ket0
  \end{equation}
  and equals $\ket0$ precisely under strong polar ground compatibility.
  \item On active support, $u_A$ is unitary onto its final support.  Under
  physical presentability and strict phenomenal excess,
  \begin{equation}
    u_A^*u_A=I_{\scrHP},
    \qquad
    u_Au_A^*=\suppProjF<I_{\scrHF}.
  \end{equation}
  \item The graph of $u_A$ is hypermaximal neutral in the active Krein
  carrier.  Its projectivization is the Born Rule Locus, and matched effects,
  probability distributions, and normal states are transported through the
  paired observable representation.
  \item The luminous ground generates the selected phenomenal determination
  algebra cyclically.  Centered determinations generate differentiated
  appearance, while exact ground realization is the projective condition
  $[\phi]=[\Omega_0]$ or, for normal states, $\rho=\PzeroF$.
  \item Ground non-alteration, reflexive terminality, and deconditioning are
  distinct structures.  The canonical deconditioning semigroup fixes
  $\PzeroF$ and contracts ground deviation and final-defect weight by
  $e^{-\gamma t}$.
  \item Phenomenal-to-physical actualization is the completely positive,
  trace-nonincreasing operation
  \begin{equation}
    \returnOp(\rho)=u_A^*\rho u_A,
    \qquad
    \Tr\!\left(\returnOp(\rho)\right)=\presentWeight(\rho).
  \end{equation}
  \item The phenomenal sector decomposes as
  \begin{equation}
    \scrHF=\GzeroF\oplus\actAppF\oplus\defectHF,
    \qquad
    \defectHF=\Ker A^*.
  \end{equation}
  If $A$ is Fredholm and $\Ker A=\{0\}$, then
  \begin{equation}
    \ind A=-\dim\Ker A^*.
  \end{equation}
\end{enumerate}

\subsection{Adopted assumptions and status}

In addition to the two cyclic-ground axioms and their bridge, the realization
adopts
\begin{equation}
  \boxed{
  \Ker A=\{0\},
  \qquad
  \Ker A^*\neq\{0\}.}
  \label{eq:adopted-asymmetry}
\end{equation}
The first condition makes every physical direction phenomenally presentable
under the normalized orientation.  The second supplies a nontrivial
phenomenal final-defect space.  Neither condition follows from the bare
closed-channel algebra.

\begin{center}
\small
\begin{tabularx}{\textwidth}{>{\raggedright\arraybackslash}p{0.22\textwidth}
>{\raggedright\arraybackslash}X}
\toprule
\textbf{Status} & \textbf{Current content} \\
\midrule
Imported mathematics & Closed-operator adjoints and polar decomposition;
Krein graph geometry; cyclic/GNS representation theory; completely positive
maps and semigroups; bounded transforms and Fredholm index theory. \\
\addlinespace
Realization axioms & Unique cyclic physical and phenomenal grounds; the
ground bridge; physical presentability; strict phenomenal excess. \\
\addlinespace
Model choices & Structural rather than measurement-theoretic sector
projections; selected correspondence and determination algebras; chosen
v\=asan\=a generator classes; a canonical, but not unique, deconditioning
semigroup. \\
\addlinespace
Derived structure & Restricted and active neutral graphs; Born Rule Locus
Correspondence; paired observables; centered appearance geometry; partial
actualization; the final-support decomposition; the conditional index formula. \\
\addlinespace
Doctrinal identifications & Emptiness and the physical ground; luminosity and
the phenomenal ground; reflexive terminality; differentiated \emph{sems};
exact \emph{rig pa}; Volition as the return orientation. \\
\bottomrule
\end{tabularx}
\end{center}

\subsection{Doctrinal synthesis}

The current interpretive map is
\begin{equation}
  \boxed{
  \begin{aligned}
  \ket0
  &\longleftrightarrow
  \text{cyclic physical vacuum and representative of emptiness},\\
  \ket{\Omega_0}
  &\longleftrightarrow
  \text{unique luminous phenomenal ground},\\
  u_A
  &\longleftrightarrow
  \text{norm-preserving phenomenal presentation},\\
  u_A^*
  &\longleftrightarrow
  \text{phenomenal-to-physical actualization},\\
  \suppProjF
  &\longleftrightarrow
  \text{physically presentable phenomenal support},\\
  \Ker A^*
  &\longleftrightarrow
  \text{phenomenal possibility without normalized physical partner},\\
  \ind A<0\ \text{(Fredholm case)}
  &\longleftrightarrow
  \text{finite robust phenomenal surplus}.
  \end{aligned}}
  \label{eq:current-doctrinal-map}
\end{equation}
Here ``mental causation'' denotes neither a breach of physical law nor an act
performed by a separate mental substance.  It denotes the return orientation
internal to the common transition channel.  The operator-theoretic statements
and their doctrinal interpretations retain different logical statuses.

\subsection{Open problems and conditional extensions}

The present architecture leaves the following questions genuinely open.
\begin{enumerate}[label=(\roman*)]
  \item \emph{Ground compatibility and ground-state structure.}  Must the full
  polar return map $\ket{\Omega_0}$ to $\ket0$, or is exact reciprocity only a
  property of the restricted ground channel?  On which selected algebras are
  the two ground states pure, faithful, or separating?
  \item \emph{Admissible algebras and generators.}  Which matter,
  correspondence, and phenomenal determination algebras are physically or
  cognitively admissible?  How does $\obsF$ sit inside, or act upon, the larger
  determination algebra, and which concrete v\=asan\=a profiles select its
  generators?
  \item \emph{Phenomenal organization.}  Which represented subalgebras encode
  object-content and affective variation, and should derived
  grasper--grasped structure use sharp projection frames or more general
  effects?
  \item \emph{Operational meaning.}  Do the structural relative-phase
  operators or the final-support weight $\presentWeight$ admit a direct
  measurement protocol, or are they only internal coordinates of the
  realization?  In \cref{sec:finite-pair-realization}, the occupation
  readout measures final-support and defect weights.  Extending that
  identification to a cognitive realization requires its own instrument.
  \item \emph{Realization dynamics.}  What concrete dynamics relates exact
  ground realization to dynamical arrival, stabilization, and report?  Is the
  canonical reduced semigroup realized by an environment or unitary dilation,
  and is any physical fallback completion of the partial return warranted?
  \item \emph{Phenomenal excess.}  Which determinations generate components in
  $\Ker A^*$?  Is the defect finite, infinite, or sector-dependent, and when do
  Fredholm hypotheses provide a stable finite index?
  \Cref{sec:finite-pair-realization} exhibits one dynamically accessible
  defect under a specified preparation.  Which independently justified
  preparation restrictions preserve such a defect?
  The orbital model in \cref{sec:orbital-ground-source} supplies a
  ground-preserving source within a specified positive parent.  Its
  three-dimensional realization still requires selection of the low
  excitation doublet and a corresponding readout.  Which independently
  justified operations implement those restrictions?
  \item \emph{Physical implementation.}  Does a concrete gauge-selection or
  control model realize the abstract Volition orientation without adding an
  independent agentive mechanism?
  \item \emph{Families of realizations.}  Given a parameter space $\mathcal S$
  of cognitive continua or realization contexts, when do the physical and
  phenomenal ground rays form line bundles with fiberwise orientations
  \begin{equation}
    \mathscr L_0^{\Psi}
    \xrightarrow{\ u_0\ }
    \mathscr L_0^{\Phi},
    \qquad
    \mathscr N_0=\Graph(u_0)\longrightarrow\mathcal S?
    \label{eq:ground-line-bundle-target}
  \end{equation}
  No nontrivial bundle topology or Grassmann curvature is asserted until a
  parameter space and a sufficiently regular projection field are supplied.
\end{enumerate}
% =============================================================================
% END INLINED FILE: sections/38_current_architecture.tex
% =============================================================================

\clearpage
\appendix
% =============================================================================
% BEGIN INLINED FILE: appendices/A_notation.tex
% =============================================================================
\section{Notation charter}
\label{app:notation}

\begin{longtable}{>{\raggedright\arraybackslash}p{0.20\textwidth}
>{\raggedright\arraybackslash}p{0.31\textwidth}
>{\raggedright\arraybackslash}p{0.39\textwidth}}
\toprule
\textbf{Symbol} & \textbf{Mathematical type} & \textbf{Reserved role} \\
\midrule
\endfirsthead
\toprule
\textbf{Symbol} & \textbf{Mathematical type} & \textbf{Reserved role} \\
\midrule
\endhead
$\scrHP$ & Hilbert space & Physical or matter sector; \emph{r\=upadh\=atu}. \\
$\scrHF$ & Hilbert space & Phenomenal sector; \emph{dharmadh\=atu}. \\
$\scrK=\scrHP\oplus\scrHF$ & Hilbert direct sum with Krein form & Common graded carrier. \\
$\fundJ$ & Bounded self-adjoint involution & Fundamental symmetry distinguishing the sectors. \\
$\PP,\PF$ & Orthogonal projections & Structural sector projections under the SSP reading. \\
$A$ & Closed typed operator $\scrHP\to\scrHF$ & Awareness channel. \\
$V=A^*$ & Closed typed operator $\scrHF\to\scrHP$ & Hilbert-adjoint return channel interpreted as Volition. \\
$\Aodd,\Vodd$ & Grading-odd operators on $\scrK$ & TAUM orientations; nilpotent and Krein-adjoint. \\
$|A|=(A^*A)^{1/2}$ & Positive self-adjoint operator & Transition magnitude. \\
$u_A$ & Polar partial isometry & Canonical norm-preserving orientation on active support. \\
$\algP,\piP$ & Unital $C^*$-algebra and representation & Matter-generating operations and their action. \\
$\algF,\piF$ & Unital $C^*$-algebra and representation & Phenomenal determinations and their action. \\
$H_{\Psi},E_0^{\Psi}$ & Self-adjoint ground generator and spectral infimum & Select the physical vacuum ray. \\
$H_{\Phi},E_0^{\Phi}$ & Self-adjoint ground generator and spectral infimum & Select the luminous phenomenal ground ray. \\
$\ket0$ & Normalized vector in $\scrHP$ & Unique cyclic physical vacuum; emptiness representative. \\
$\ket\Omega=A\ket0$ & Nonzero vector in $\scrHF$ & Relative luminosity; awareness of emptiness. \\
$\ket{\Omega_0}$ & Normalized vector in $\scrHF$ & Unique cyclic intrinsic luminous ground. \\
$\PzeroP,\QzeroP$ & Complementary projections on $\scrHP$ & Physical ground ray and ground-orthogonal sector. \\
$\PzeroF,\QzeroF$ & Complementary projections on $\scrHF$ & Luminous ground ray and differentiated sector. \\
$\GzeroP,\GzeroF$ & One-dimensional Hilbert subspaces & Physical and phenomenal ground rays as linear spaces. \\
$\Kzero=\GzeroP\oplus\GzeroF$ & Two-dimensional Krein subspace & Ground plane of signature $(1,1)$. \\
$\Pizero$ & Rank-two orthogonal projection on $\scrK$ & Support projection of the ground plane. \\
$\omegaP$ & State on $\algP$ & Physical vacuum state functional. \\
$\omegaF$ & State on $\algF$ & Luminous phenomenal ground-state functional. \\
$F_0^{\Psi},F_0^{\Phi}$ & Numbers in $[0,1]$ & Ground fidelities; not amounts of matter or luminosity. \\
$a_0=\|A\ket0\|$ & Positive scalar & Raw ground-transition magnitude after phase fixing. \\
$\Azero=A\PzeroP$ & Bounded rank-one operator & Restricted ground channel, even when $A$ is unbounded. \\
$\uzero=a_0^{-1}\Azero$ & Rank-one partial isometry & Canonical bridge-selected orientation from $\GzeroP$ to $\GzeroF$. \\
$\NzeroPlus,\NzeroMinus$ & One-dimensional neutral graph subspaces & Aligned and anti-aligned ground relations. \\
$\ket{\Xi_0^+},\ket{\Xi_0^-}$ & Hilbert-normalized vectors in $\Kzero$ & Normalized generators of the two canonical graph lines. \\
$\Rzero$ & Self-adjoint partial involution & Ground reflection; $\pm1$ eigenspaces are $\NzeroPlus$ and $\NzeroMinus$. \\
$\Izero=\fundJ\Rzero$ & Skew-adjoint partial complex structure & Ground quarter-turn with $\Izero^2=-\Pizero$. \\
$\Qzero=-i\Izero$ & Self-adjoint partial involution & Relative-phase quadrature on the ground plane. \\
$\RA,\IA$ & Full supportwise polar operators & Active-support reflection and quarter-turn derived from $u_A$. \\
$R_{\Psi}=A^*A$, $R_{\Phi}=AA^*$ & Positive structural returns & Not automatically identical with $H_{\Psi}$ and $H_{\Phi}$. \\
$\suppProjP,\suppProjF$ & Orthogonal support projections & Initial and final supports of the polar orientation; TAUM crosswalk $\pi_{A,+},\pi_{A,-}$. \\
$\actHP,\actHF$ & Closed Hilbert subspaces & Initial physical support $(\ker A)^\perp$ and final phenomenal support $\overline{\operatorname{Ran}A}$. \\
$\actK$ & Active Krein subspace & $\actHP\oplus\actHF=\operatorname{Ran}\PiA$. \\
$\WAtheta$ & Isometric embedding $\actHP\to\actK$ & Normalized phase-graph embedding. \\
$\NAtheta$ & Hypermaximal neutral subspace & Graph of $e^{i\theta}u_A$; $\NAplus$ is aligned. \\
$\PAtheta$ & Hilbert-orthogonal projection & Projection onto $\NAtheta$. \\
$\LBR$ & Projective subspace & Born Rule Locus $\mathbb P(\NAplus)$. \\
$\BRprob$ & Function on pairs of rays & Pure-state Born transition probability. \\
$\alphaA$ & Unital normal $*$-isomorphism & Conjugates active physical operators into corresponding phenomenal operators. \\
$\obsP,\obsF$ & Selected unital $C^*$-algebras & Active physical correspondence algebra and its phenomenal image. \\
$\iotaMaxA$ & Unital normal $*$-homomorphism & Maximal embedding on $\mathcal B(\actHP)$. \\
$\iotaA$ & Faithful $*$-homomorphism & Restriction of $\iotaMaxA$ to $\obsP$. \\
$\obsPair$ & Paired $C^*$-algebra & Cross-domain correspondence algebra on $\actK$. \\
$\UA$ & Unitary factorization map & Identifies $\actK$ with $\mathbb C^2\otimes\actHP$. \\
\bottomrule
\end{longtable}

\subsection{Typographical discipline}

The document uses the bare italic $A$ for the typed channel and the
sans-serif $\mathsf A$ for its odd extension.  This distinction is mandatory.
Likewise, the bare zero vector $0$, the vacuum ket $\ket0$, and the ground
label in $\ket{\Omega_0}$ are never interchangeable.

The notation
\begin{equation}
  [\psi]
\end{equation}
means the projective ray of a nonzero vector $\psi$.  A phase change does not
produce a different ray or vector functional:
\begin{equation}
  [e^{i\theta}\psi]=[\psi].
\end{equation}

The word \emph{null} or \emph{neutral} always refers to the Krein form:
\begin{equation}
  \kInner{\xi}{\xi}=0.
\end{equation}
It does not imply $\xi=0$ or $\|\xi\|_{\scrK}=0$.  The vectors
$\ket{\Xi_0^\pm}$ are Hilbert-normalized and Krein-null.

The symbols $\omegaP$ and $\omegaF$ denote positive normalized functionals,
not vectors.  The correspondences
\begin{equation}
  \omegaP\leftrightarrow\ket0,
  \qquad
  \omegaF\leftrightarrow\ket{\Omega_0}
\end{equation}
are representation correspondences inside cyclic pointed representations,
not literal equalities of mathematical types.

\subsection{Appearance and determination notation}

\begin{longtable}{>{\raggedright\arraybackslash}p{0.20\textwidth}
>{\raggedright\arraybackslash}p{0.31\textwidth}
>{\raggedright\arraybackslash}p{0.39\textwidth}}
\toprule
\textbf{Symbol} & \textbf{Mathematical type} & \textbf{Reserved role} \\
\midrule
\endfirsthead
\toprule
\textbf{Symbol} & \textbf{Mathematical type} & \textbf{Reserved role} \\
\midrule
\endhead
$\detAlg$ & Selected unital $C^*$-subalgebra of $\algF$ & Admissible phenomenal determinations. \\
$\nullIdeal$ & Closed left ideal in $\detAlg$ & Operators whose action on $\ket{\Omega_0}$ vanishes. \\
$\centDet=\ker\omegaF$ & Centered linear subspace & Ground-zero-mean determinations; generally not an algebra or ideal. \\
$\appSpace$ & Closed Hilbert subspace & Differentiated appearance space $\QzeroF\scrHF$. \\
$\Cov_0^{\Phi}$ & Positive semidefinite sesquilinear kernel & Gram geometry of centered phenomenal content. \\
$\vasSpace$ & Parameter or profile set & Possible v\=asan\=a-conditioning data. \\
$\genSet(\vas)$ & Family of bounded determination generators & Active generators selected by a disposition profile. \\
$\wordSet(\vas)$ & Unital $*$-algebra of finite words & Ordered histories of active determination operations. \\
$\Hvas$ & Closed cyclic subspace & Phenomenal subspace generated under profile $\vas$. \\
$\Evas$ & Closed centered subspace & Profile-relative differentiated appearance. \\
Functional generator families & Subsets of $\genSet(\vas)$ & $\Gak$ (image), $\Gvik$ (conceptual), $\Gsmr$ (mnemonic), $\Gaff$ (affective), and $\Gdv$ (duality-producing). \\
$\contentAlg,\affectAlg$ & Represented subalgebras & Chosen object-content and affective comparison algebras. \\
$\Pgrasper,\Pgrasped$ & State-dependent orthogonal projections & Derived grasper and grasped aspects inside $\scrHF$. \\
$\Pdual$ & Internal support projection & Sum of the two derived duality projections; satisfies $\Pdual\leq\QzeroF$. \\
$\Gammadual$ & Self-adjoint partial involution & Internal derived-duality grading on $\operatorname{Ran}\Pdual$. \\
$\semsSet$ & Subset of projective phenomenal space & Rays with nonzero differentiated component. \\
$\rigRay$ & Singleton projective set & Exact luminous ground ray representing rig pa. \\
$\delta_0(\phi)$ & Number in $[0,1]$ & Ground deviation or adulteration $1-F_0^{\Phi}(\phi)$; not inverse luminosity. \\
\bottomrule
\end{longtable}

\subsection{Dynamics and realization notation}

\begin{longtable}{>{\raggedright\arraybackslash}p{0.20\textwidth}
>{\raggedright\arraybackslash}p{0.31\textwidth}
>{\raggedright\arraybackslash}p{0.39\textwidth}}
\toprule
\textbf{Symbol} & \textbf{Mathematical type} & \textbf{Reserved role} \\
\midrule
\endfirsthead
\toprule
\textbf{Symbol} & \textbf{Mathematical type} & \textbf{Reserved role} \\
\midrule
\endhead
$\mathfrak S(\scrHF)$ & Convex set of positive trace-one operators & Normal phenomenal state space. \\
$F_0^{\Phi}(\rho)$ & Number in $[0,1]$ & Mixed-state ground fidelity $\Tr(\PzeroF\rho)$. \\
$\delta_0^{\Phi}(\rho)$ & Number in $[0,1]$ & Mixed-state ground deviation $1-\Tr(\PzeroF\rho)$. \\
$\mathcal E$ & Completely positive trace-preserving map & Generic phenomenal state-evolution channel. \\
$\mathcal D_t$ & CPTP semigroup & Canonical ground-attractive deconditioning model. \\
$\mathcal L_0$ & Lindblad-type generator & Generator of the canonical deconditioning semigroup. \\
$(\mathscr H_{\Phi,n},P_{0,n}^{\Phi},\rho_n,\mathcal T_n)$ & Sequential channel system & Momentary cognitive continuum with transported ground type. \\
$\detAlgEmpty$ & Profile-relative $C^*$-algebra & Norm closure of operations active under the exhausted v\=asan\=a profile. \\
$\chi_0^{\Phi}$ & Character on $\detAlgEmpty$ & Scalar action of the exhausted-profile algebra on the luminous ground. \\
\bottomrule
\end{longtable}

\subsection{Return, defect, and index notation}

\begin{longtable}{>{\raggedright\arraybackslash}p{0.20\textwidth}
>{\raggedright\arraybackslash}p{0.31\textwidth}
>{\raggedright\arraybackslash}p{0.39\textwidth}}
\toprule
\textbf{Symbol} & \textbf{Mathematical type} & \textbf{Reserved role} \\
\midrule
\endfirsthead
\toprule
\textbf{Symbol} & \textbf{Mathematical type} & \textbf{Reserved role} \\
\midrule
\endhead
$\ket{0_A}=u_A^*\ket{\Omega_0}$ & Unit vector in $\actHP$ & Full polar preimage of the luminous ground; equals $\ket0$ only under strong compatibility. \\
$\suppProjF=u_Au_A^*$ & Orthogonal projection on $\scrHF$ & Final phenomenal support projection. \\
$\defectProjF=I_{\Phi}-\suppProjF$ & Orthogonal projection on $\scrHF$ & Projection onto the final-defect space. \\
$\defectHF=\Ker A^*$ & Closed Hilbert subspace & Final-defect space, interpreted as phenomenal excess. \\
$\actAppF=\QzeroF\actHF$ & Closed Hilbert subspace & Differentiated phenomenal appearance inside the final support. \\
$\presentWeight(\phi)$ & Number in $[0,1]$ & Pure-state final-support weight $\|u_A^*\phi\|^2$. \\
$\presentWeight(\rho)$ & Number in $[0,1]$ & Normal-state final-support weight $\Tr(\suppProjF\rho)$. \\
$\defectWeight(\rho)$ & Number in $[0,1]$ & Defect/excess weight $\Tr(\defectProjF\rho)=1-\presentWeight(\rho)$. \\
$\returnOp(\rho)$ & CP trace-nonincreasing map & Subnormalized global physical return $u_A^*\rho u_A$. \\
$\returnOpCond(\rho)$ & Normal physical state & Conditional global return when $\presentWeight(\rho)>0$. \\
$\returnGroundOp(\rho)$ & CP trace-nonincreasing map & Restricted ground branch $u_0^*\rho u_0$. \\
$\cohZero,\cohOne$ & Finite-dimensional cohomology spaces under Fredholm hypotheses & Physical null directions and phenomenal cokernel directions. \\
$\dim\Ker A^*$ & Nonnegative integer or $\infty$ & Phenomenal defect/excess multiplicity. \\
$\ind A$ & Integer for Fredholm $A$ & Unpaired physical count minus unpaired phenomenal count. \\
\bottomrule
\end{longtable}
% =============================================================================
% END INLINED FILE: appendices/A_notation.tex
% =============================================================================
% =============================================================================
% BEGIN INLINED FILE: appendices/B_provenance.tex
% =============================================================================
\section{Scholarly provenance and source status}
\label{app:provenance}

This appendix records the source status of the mathematical and interpretive
layers used in the cognitive realization.  It is a provenance note, not an
additional proof layer.

\subsection{Mathematical chassis and standard results}

The interpretation-neutral closed-channel architecture---the graded carrier,
nilpotent orientations, positive returns, polar geometry, support
projections, and status discipline---is imported from \emph{Toward a
Universal Model} \cite{TAUM2026}.  The present document adds realization data
and cognitive interpretation; it does not treat those additions as theorems of
the bare chassis.

The Krein-space language of fundamental symmetries, neutral graph subspaces,
and indefinite inner-product geometry is standard; see Bogn\'ar
\cite{Bognar}.  Cyclic representations, vector states, GNS quotients,
faithfulness, and purity use standard operator-algebraic theory as presented,
for example, by Bratteli and Robinson \cite{BratteliRobinson}.  Closed
operators, polar decomposition, form domains, bounded transforms, and
Fredholm stability follow the standard functional-analytic framework of Kato
and Reed--Simon \cite{Kato,ReedSimon}.

The completely positive dynamical language and the canonical semigroup model
use the Gorini--Kossakowski--Sudarshan--Lindblad framework
\cite{GKS1976,Lindblad1976}.  The trace-distance estimate from ground fidelity
uses the Fuchs--van de Graaf inequality \cite{FuchsVandeGraaf1999}.  No novelty
is claimed for these general results themselves.

The finite pair-transfer realization in \cref{sec:finite-pair-realization}
uses the canonical Hamiltonian and in-phase lower state reported by
Singh in Paper~10, following Paper~8 \cite{Singh2026PairTransfer}.
The canonical modes and positive adjoints are assumed in that calculation.
The full-sector uniqueness check, preparation, readout, ground
compatibility and current-selection tests are supplied here, as is the
response calculation showing that Singh's minimal symmetric positive
completion removes the accessible defect of this preparation.  The bare
component is therefore not presented as a complete positive-exchange
derivation of that defect.
The orbital parent in \cref{sec:orbital-ground-source} is a separate
model specified and calculated in this manuscript.  Its full-space
ground argument uses standard analytic perturbation theory and the
min--max principle \cite{Kato,ReedSimon}.  Neither the orbital parent
nor a microscopic derivation of its operational interface is attributed
to Singh.  Cognitive and doctrinal assignments remain additional
identifications.

\subsection{Interpretive sources}

The doctrinal distinction between luminosity as intrinsic to cognition and
the exact realization of undifferentiated luminosity is grounded in the
author's study of reflexive awareness and the Third and Fourth Yogas
\cite{Yiannopoulos2012}.  The account of v\=asan\=a-conditioned appearance,
conceptual and nonconceptual determination, mnemonic concatenation,
grasper--grasped distortion, and the difference between intrinsic reflexivity
and its definitive ascertainment is drawn from the author's reconstruction of
Dharmak\=irti's philosophy of cognition \cite{Yiannopoulos2020}.

These sources motivate the doctrinal identifications made in the text.  They
do not prove the operator-theoretic propositions, select unique operator
representatives for Buddhist technical terms, or entail the realization
axioms.

\subsection{Status of the present construction}

The distinctive contribution of this document is the organization of the
imported mathematical chassis into a typed cognitive realization with two
cyclic grounds, a bridge-selected neutral relation, a paired Born geometry,
an algebra of differentiated appearance, explicit realization and
deconditioning conditions, and an asymmetric return with phenomenal
final-defect space.  Its claims must be read at three levels:
\begin{enumerate}[label=(\roman*)]
  \item imported or standard mathematics, valid under the stated analytic
  hypotheses;
  \item realization axioms and model choices, which select one cognitive
  implementation of that mathematics;
  \item doctrinal identifications, which assign philosophical or contemplative
  meanings to the selected structures.
\end{enumerate}
No implication is asserted from the third level back to the first, and no
realization axiom is presented as a consequence of the interpretation-neutral
channel algebra.
% =============================================================================
% END INLINED FILE: appendices/B_provenance.tex
% =============================================================================

\bigskip
\begin{thebibliography}{99}

\bibitem{Yiannopoulos2012}
A.~Yiannopoulos,
\emph{Luminosity: Reflexive Awareness in Ratn\=akara\'s\=anti's Pith
Instructions for the Ornament of the Middle Way},
M.A. thesis, Rangjung Yeshe Institute at Kathmandu University, 2012.

\bibitem{Yiannopoulos2020}
A.~Yiannopoulos,
\emph{The Structure of Dharmak\=irti's Philosophy: A Study of
Object-Cognition in the Perception Chapter (pratyak\d{s}apariccheda) of the
Pram\=a\d{n}asamuccaya, the Pram\=a\d{n}av\=arttika, and Their Earliest
Commentaries},
Ph.D. dissertation, Emory University, 2020.

\bibitem{TAUM2026}
A.~Yiannopoulos,
\emph{Toward a Universal Model},
internal manuscript, July 2026.

\bibitem{Singh2026PairTransfer}
T.~P. Singh,
\emph{Interacting Generalised Trace Dynamics: Cohesive Emergence, Jordan
Spectra and the Problem of Fundamental Constants},
Paper 10 in the series P1--P10, manuscript dated 31 October 2026,
Zenodo, 2026.
\href{https://doi.org/10.5281/zenodo.23853851}{doi:10.5281/zenodo.23853851}.

\bibitem{GKS1976}
V.~Gorini, A.~Kossakowski, and E.~C.~G. Sudarshan,
Completely positive dynamical semigroups of $N$-level systems,
\emph{Journal of Mathematical Physics} \textbf{17} (1976), 821--825.

\bibitem{Lindblad1976}
G.~Lindblad,
On the generators of quantum dynamical semigroups,
\emph{Communications in Mathematical Physics} \textbf{48} (1976), 119--130.

\bibitem{FuchsVandeGraaf1999}
C.~A. Fuchs and J.~van de Graaf,
Cryptographic distinguishability measures for quantum-mechanical states,
\emph{IEEE Transactions on Information Theory} \textbf{45} (1999), 1216--1227.

\bibitem{BratteliRobinson}
O.~Bratteli and D.~W.~Robinson,
\emph{Operator Algebras and Quantum Statistical Mechanics, Vol.~1},
2nd ed., Springer, 1987.

\bibitem{Bognar}
J.~Bogn\'ar,
\emph{Indefinite Inner Product Spaces},
Springer, 1974.

\bibitem{Kato}
T.~Kato,
\emph{Perturbation Theory for Linear Operators},
Springer, 2nd ed., 1995.

\bibitem{ReedSimon}
M.~Reed and B.~Simon,
\emph{Methods of Modern Mathematical Physics, Vol. I: Functional Analysis},
Academic Press, 1980.

\end{thebibliography}

\end{document}
